% In this file you should put the actual content of the blueprint.
% It will be used both by the web and the print version.
% It should *not* include the \begin{document}
%
% If you want to split the blueprint content into several files then
% the current file can be a simple sequence of \input. Otherwise It
% can start with a \section or \chapter for instance.

This is the blueprint of a Lean~4 library on extreme value theory and heavy-tailed
distributions. It grew out of the Aalto ExtremeValueProject, developed for the course
\emph{MS-EV0029 Introduction to Formalized Mathematics in Lean}
at Aalto University (February--June 2025) by Kalle Kytölä and the course participants,
whose goal was to formalize the Fisher--Tippett--Gnedenko theorem
classifying univariate extreme value distributions.

The library covers:
\begin{itemize}
  \item cumulative distribution functions, their affine transformations, types and
    convergence, and the classification of extreme value distributions
    (Fisher--Tippett--Gnedenko);
  \item regular variation: slowly and regularly varying functions, the uniform convergence
    theorem (monotone and measurable cases), Potter bounds and the Karamata representation
    theorem;
  \item Karamata's theorem for tail integrals and Breiman's lemma;
  \item domains of attraction: Gnedenko's characterizations of the Fréchet and Weibull
    maximum domains of attraction, and de Haan's characterization of the Gumbel domain of
    attraction via auxiliary functions (with the von Mises condition);
  \item peaks over threshold: the Pickands--Balkema--de Haan theorem for all three domains of
    attraction and the asymptotics of the mean excess function;
  \item statistics of extremes: order statistics, binomial exceedance counts, and weak
    consistency of the Hill estimator;
  \item heavy tails: subexponential distributions, long-tailedness, convolution powers,
    the single big jump principle and Kesten's bound;
  \item stable laws: Gaussian, Cauchy and symmetric $\alpha$-stable laws, Lévy's continuity
    theorem, the generalized central limit theorem for symmetric regularly varying summands,
    and the existence of symmetric stable laws.
\end{itemize}

\textbf{WARNING: The blueprint is still work in progress!}

\chapter{Cumulative distribution functions}

\begin{definition}
  \label{def:cdf}
  \lean{CumulativeDistributionFunction}
  \leanok
  A function $F \colon \R \to \R$ is a
  \emph{cumulative distribution function} (c.d.f.) if
  \begin{enumerate}
  \item[(i)] $x \mapsto F(x)$ is increasing;
  \item[(ii)] $x \mapsto F(x)$ is right-continuous;
  \item[(iii)] $\lim_{x \to -\infty} F(x) = 0$ and $\lim_{x \to +\infty} F(x) = 1$.
  \end{enumerate}
\end{definition}

\begin{lemma}
  \label{lem:cdf-of-random-var}
  \uses{def:cdf}
  \lean{MeasureTheory.ProbabilityMeasure.cdf}
  \leanok
  If $X$ is a real-valued random variable, then the function
  $F \colon \R \to \R$ given by $F(x) = \PR \big[ X \le x \big]$ is a c.d.f.
\end{lemma}
\begin{proof}
  % \leanok
  Property (1.) in Definition~\ref{def:cdf} is obvious (by monotonicity of measures)
  and properties (2.) and (3.) are simple consequences of monotone convergence
  theorems for probability measures.
\end{proof}

\section{Degenerate distributions}

\begin{definition}
  \label{def:degenerate-cdf}
  \uses{def:cdf}
  \lean{CumulativeDistributionFunction.IsDegenerate}
  \leanok
  A c.d.f. $F$ is said to be
  \emph{degenerate} if for every $x \in \R$ we have either $F(x) = 0$ or $F(x) = 1$.
  Otherwise $F$ is said to be \emph{nondegenerate}.
\end{definition}

\begin{lemma}
  \label{lem:degenerate-cdf-iff-exists-jump}
  \uses{def:degenerate-cdf}
  \lean{CumulativeDistributionFunction.isDegenerate_iff}
  \leanok
  $F$ is a degenerate c.d.f. if and only if there exists a $x_0 \in \R$ such that
  \begin{align*}
  F(x) = \begin{cases}
          0 & \text{ for } x < x_0 \\
          1 & \text{ for } x \ge x_0 .
         \end{cases}
  \end{align*}
\end{lemma}
\begin{proof}
  \uses{}
  \leanok
  The ``if'' direction is clear. To prove the ``only if'' direction, assume that
  $F$ is a degenerated c.d.f., and let $x_0 = \inf \set{ x \in \R \; \big| \; F(x) = 1}$.
  Then it is straightforward to show by properties of a c.d.f. that $F$ has the asserted form.
\end{proof}

\begin{lemma}
  \label{lem:delta-has-degenerate-cdf}
  \uses{def:degenerate-cdf}
  \lean{CumulativeDistributionFunction.diracProba_is_degenerate}
  \leanok
  The c.d.f. of Dirac delta mass $\delta_{x_0}$ at $x_0 \in \R$ is degenerate.
\end{lemma}
\begin{proof}
  \uses{lem:degenerate-cdf-iff-exists-jump}
  \leanok
  \ldots
\end{proof}

\begin{lemma}
  \label{lem:degenerate-cdf-is-delta}
  \uses{def:degenerate-cdf}
  \lean{CumulativeDistributionFunction.eq_diracProba_of_isDegenerate}
  \leanok
  If a c.d.f. $F$ is degenerate, then it is the c.d.f. of a Dirac delta mass $\delta_{x_0}$
  at some point $x_0 \in \R$.
\end{lemma}
\begin{proof}
  \uses{lem:degenerate-cdf-iff-exists-jump}
  \leanok
  \ldots
\end{proof}

\section{Distributions of maxima of independent random variables}

% \begin{lemma}
% \label{lem:random-var-of-cdf}
% \uses{def:cdf}
% If $F \colon \R \to \R$ is a c.d.f., then there exists a probability
% space $(\Omega, \mathcal{F}, \PR)$ and a real-valued random variable
% $X \colon \Omega \to \R$ such that for every $x \in \R$
% we have $F(x) = \PR \big[ X \le x \big]$.
% \end{lemma}
% \begin{proof}
% Abstract nonsense (take $\Omega = [0,1]$ and $X = \idOf{\R}$ ).
% \end{proof}

\begin{lemma}
  \label{lem:cdf-of-max-two}
  \uses{def:cdf}
  % \leanok
  Let $X$ and $Y$ be two independent real-valued random variables with respective
  cumulative distribution functions $F$ and $G$,
  i.e. $F(x) = \PR \big[ X \le x \big]$ and $G(x) = \PR \big[ Y \le x \big]$.
  Then the c.d.f. of $M = \max (X, Y)$ is $x \mapsto F(x) \, G(x)$.
\end{lemma}
\begin{proof}
  \uses{}
  % \leanok
  Fix $x \in \R$. Note that $\max (X, Y) \le x$ if and only if both $X \le x$ and $Y \le x$.
  Calculate, using independence,
  \begin{align*}
  \PR\big[ \max (X, Y) \le x \big]
  \; = \; \PR\big[ X \le x \, \; Y \le x \big]
  \; = \; \PR\big[ X \le x \big] \; \PR\big[ Y \le x \big]
  \; = \; F(x) \, G(x) .
  \end{align*}
\end{proof}

\begin{lemma}
  \label{lem:cdf-of-max-many}
  \uses{def:cdf}
  % \leanok
  Let $X_0, X_1, \ldots, X_{n-1}$ be independent identically distributed real-valued random variables
  with cumulative distribution functions $F$,
  i.e. $F(x) = \PR \big[ X_j \le x \big]$ for every $j$.
  Then the c.d.f. of \[ M_n = \max_{0 \le j < n} X_j\] is the function $x \mapsto \big(F(x)\big)^n$.
\end{lemma}
\begin{proof}
  \uses{lem:cdf-of-max-two}
  % \leanok
  Induction on~$n$ using \ref{lem:cdf-of-max-two}.
\end{proof}

\begin{lemma}
  \label{lem:cdf-of-affine-max-many}
  \uses{def:cdf}
  % \leanok
  Let $X_0, X_1, \ldots, X_{n-1}$ be independent identically distributed real-valued random variables
  with cumulative distribution functions $F$,
  i.e. $F(x) = \PR \big[ X_j \le x \big]$ for every $j$,
  and let $a > 0$ and $b \in \R$.
  Then the c.d.f. of
  \begin{align*}
  \hat{M}_n = \frac{\max_{0 \le j < n} X_j - b}{a}
  \end{align*}
  is the function $x \mapsto \big(F(a x + b)\big)^n$.
\end{lemma}
\begin{proof}
  \uses{lem:cdf-of-max-many}
  % \leanok
  Use \ref{lem:cdf-of-max-many} and do a change of variables.
\end{proof}

\section{Distributions of minima of independent random variables}

\section{Equivalence classes modulo affine transformations}

\begin{definition}
  \label{def:oriented-affine-isomorphism}
  \lean{AffineIncrEquiv}
  \leanok
  The collection of all transformations $\R \to \R$ of the form
  $x \mapsto a x + b$, where $a>0$, $b \in \R$, forms a group.
  We call this the \emph{orientation preserving affine isomorphism group}
  and denote it by $\OriAffR$.
\end{definition}

\begin{definition}
  \label{def:oriented-affine-transform-of-cdf}
  \uses{def:oriented-affine-isomorphism, def:cdf}
  \lean{CumulativeDistributionFunction.affineTransform}
  \leanok
  The action of an orientation preserving affine isomorphism $A \in \OriAffR$
  on a cumulative distribution function $F$ is defined so
  that $A.F \colon \R \to \R$ is given by $(A.F)(x) = F(A^{-1}(x))$.
  Then $A.F$ is also a c.d.f.
\end{definition}

\begin{lemma}
  \label{lem:oriented-affine-action-on-cdf}
  \uses{def:oriented-affine-transform-of-cdf}
  \lean{CumulativeDistributionFunction.instMulActionAffineIncrEquiv}
  \leanok
  The actions of orientation preserving affine isomorphisms
  on a cumulative distribution functions is a group action, i.e.,
  $1.F = F$ and $(AB).F = A.(B.F)$ for any c.d.f. $F$ and any $A,B \in \OriAffR$.
\end{lemma}
\begin{proof}
  \uses{}
  % \leanok
  Direct calculations.
\end{proof}

\begin{lemma}
  \label{lem:degenerate-cdf-transform}
  \uses{lem:oriented-affine-action-on-cdf, def:degenerate-cdf}
  \lean{CumulativeDistributionFunction.affine_isDegenerate_iff}
  \leanok
  Let $F$ be a cumulative distribution function
  and $A \in \OriAffR$ an orientation preserving affine isomorphism.
  Then $A.F$ is degenerate if and only if $F$ is degenerate.
\end{lemma}
\begin{proof}
  \uses{}
  \leanok
  Straightforward from the definitions.
\end{proof}

\begin{lemma}
  \label{lem:cdf-continuity-pt-transform}
  \uses{lem:oriented-affine-action-on-cdf}
  \lean{CumulativeDistributionFunction.affine_continuousAt_of_continuousAt}
  \leanok
  Let $F$ be a cumulative distribution function,
  and $A \in \OriAffR$ an orientation preserving affine isomorphism.
  If a point $x \in \R$ is a continuity point of $F$, then
  the point $A(x) \in \R$ is a continuity point of $A.F$.
\end{lemma}
\begin{proof}
  \uses{}
  % \leanok
  Straightforward.
\end{proof}


% \begin{lemma}
% \label{lem:cdf-affine-group-action}
% \uses{def:cdf}
% Suppose $a > 0$ and $b \in \R$. Then for any c.d.f. $F$, also
% the function $x \mapsto F(a x + b)$ is a c.d.f.
% \end{lemma}
% \begin{proof}
% \ldots
% \end{proof}
%
% \begin{definition}
% \label{def:cdf}
% %\lean{CumulativeDistributionFunction}
% %\uses{def:degenerate-cdf}
% \uses{def:cdf}
% A function $F \colon \R \to \R$ is a
% \emph{cumulative distribution function} (c.d.f.) if
% \begin{enumerate}
% \item[(i)] $x \mapsto F(x)$ is increasing;
% \item[(ii)] $x \mapsto F(x)$ is right-continuous;
% \item[(iii)] $\lim_{x \to -\infty} F(x) = 0$ and $\lim_{x \to +\infty} F(x) = 1$.
% \end{enumerate}
% \end{definition}
%

\section{Miscellaneous results on cumulative distribution functions}

\begin{lemma}[Continuity points of c.d.f.s are those which carry no point mass]
  \label{lem:cdf-continuity-pt-iff-measure-singleton}
  \uses{def:cdf}
  \lean{CumulativeDistributionFunction.continuousAt_iff}
  \leanok
  Let $F$ be cumulative distribution function of a
  probability measure $\PRmeas$ on $\bR$. A point $x \in \bR$ is
  a continuity point of $F$ if and only if $\PRmeas[\set{x}] = 0$.
\end{lemma}
\begin{proof}
  % \uses{}
  % \leanok
  A c.d.f. is always continuous from the right.

  Continuity of $F$ from the left at $x$ means that for any
  sequence $(x_n)_{n \in \bN}$ increasing to $x$
  (i.e., $x_n \le x_{n+1} < x$ for all $n \in \bN$)
  \begin{align*}
    F(x_n) \to F(x) ,
  \end{align*}
  or equivalently in terms of measures
  \begin{align*}
    \PRmeas \big[ (-\infty,x_n] \big] \to \PRmeas \big[ (-\infty,x] \big] .
  \end{align*}
  But by monotone convergence of measures, we always have
  \begin{align*}
    \PRmeas \big[ (-\infty,x_n] \big]
    \to \; & \PRmeas \big[ (-\infty,x) \big] \\
    = \; & \PRmeas \big[ (-\infty,x] \big] - \PRmeas[\set{x}] .
  \end{align*}
  A comparison of these conditions shows that $F$ is also continuous from the left
  at $x$ if and only if $\PRmeas[\set{x}] = 0$.
\end{proof}

\begin{lemma}[A pair of nontrivial continuity points of nondegenerate c.d.f.]
  \label{lem:exists-two-nontrivial-continuity-pts-cdf}
  \uses{def:degenerate-cdf}
  \lean{CumulativeDistributionFunction.exists₂_continuousAt_of_not_isDegenerate}
  \leanok
  Let $G$ be a nondegenerate c.d.f. Then there exists
  continuity points $x_1 < x_2$ of $G$ such that
  $0 < G(x_1) \le G(x_2) < 1$.
\end{lemma}
\begin{proof}
  % \uses{}
  % \leanok
  Since $G$ is nondegenerate, there exists some $x_0 \in \bR$ such that
  $0 < G(x_0) < 1$. Since $G$ is continuous from the right, for some
  small $\delta > 0$ we have that $0 < G(x) < 1$ for
  all $x \in [x_0,x_0+\delta)$. Since the continuity points of $G$
  are dense, from any nonempty open interval we may pick a continuity point.
  First pick a continuity point $x_1 \in (x_0,x_0+\delta)$,
  and then pick another continuity point $x_2 \in (x_1,x_0+\delta)$.
\end{proof}


\begin{lemma}[Equality of c.d.f.s on a dense set suffices]
  \label{lem:cdf-equal-on-dense}
  \uses{def:cdf}
  \lean{CumulativeDistributionFunction.eq_of_forall_dense_eq}
  \leanok
  Suppose that $F,G$ are two c.d.f.s and $S \subseteq \bR$ is a dense
  subset of the real line. If $F(\xi) = G(\xi)$ for all $\xi \in S$,
  then we have $F = G$.
\end{lemma}
\begin{proof}
  \uses{}
  \leanok
  We must prove that for any $x \in \bR$ we have $F(x) = G(x)$.
  But by right-continuity of c.d.f.s, density of $S$, and coincidence of
  $F$ and $G$ on $S$, we have
  \begin{align*}
      F(x)
    = \; & \lim_{\xi \to x^+ \text{ along $S$}} F(\xi) \\
    = \; & \lim_{\xi \to x^+ \text{ along $S$}} G(\xi)
    = \, G(x) .
  \end{align*}
\end{proof}



\section{Topology on orientation-preserving affine isomorphisms}

\begin{definition}[]
  \label{def:affine-transform-topology}
  \uses{def:oriented-affine-isomorphism}
  \lean{instTopologicalSpaceAffineIncrEquiv}
  \leanok
  We equip the space $\OriAffR$ of orientation-preserving affine isomorphisms
  with the topology of pointwise convergence, i.e., with the
  coarsest topology which makes the evaluations $A \mapsto A(x)$ continuous
  for all $x \in \bR$.
\end{definition}

\begin{lemma}[The coefficients of affine map depend continuously on the map]
  \label{lem:affine-coefficients-continuous}
  \uses{def:oriented-affine-isomorphism}
  \lean{AffineIncrEquiv.continuous_coefs_fst, AffineIncrEquiv.continuous_coefs_snd}
  \leanok
  The coefficients $a$ and $b$ of an orientation-preserving affine isomorphism
  $A(x) = a x + b$ depend continuously on $A$.
\end{lemma}
\begin{proof}
  % \uses{}
  % \leanok
  We may first write $a = A(1) - A(0)$ and $b = A(0)$.
  These depend continuously on $A$,
  since the evaluations $A(1)$ and $A(0)$ do.
\end{proof}

\begin{lemma}[Metrizability of the topology on oriented affine isomorphisms]
  \label{lem:affine-metrizable}
  \uses{def:affine-transform-topology}
  \lean{AffineIncrEquiv.homeomorph_prod, AffineIncrEquiv.metrizableSpace}
  \leanok
  The topology of pointwise convergence makes $\OriAffR$
  homeomorphic to $\bR^2$, and in particular
  metrizable.
\end{lemma}
\begin{proof}
  \uses{lem:affine-coefficients-continuous}
  % \leanok
  The essential claim is that the function
  \begin{align*}
    \mathrm{cfs} \colon \OriAffR \to (0,+\infty) \times \bR
  \end{align*}
  obtained by mapping $A(x) = a x + b$ to its coefficients $(a,b)$ is a homeomorphism.
  (The homeomorphism to $\bR^2 = \bR \times \bR$ follows by combining with the
  homeomorphism $(a,b) \mapsto (\log a, b)$.)

  The continuity of $\mathrm{cfs}$ follows from Lemma~\ref{lem:affine-coefficients-continuous}.
  For the continuity of the inverse, we must only check that for any $x \in \bR$,
  its composition with the point evaluation at $x$ is continuous. But the composition
  is $(a,b) \mapsto a x + b$, and the continuity is clear.
\end{proof}

\begin{lemma}[Inversion of orientation preserving affine isomorphisms is continuous]
  \label{lem:affine-inversion-continuous}
  \uses{def:oriented-affine-isomorphism}
  \lean{AffineIncrEquiv.continuous_inv, AffineIncrEquiv.invHomeomorph}
  \leanok
  The map $A \mapsto A^{-1}$ is continuous on $\OriAffR$.
  In particular, inversion defines a homeomorphism of $\OriAffR$ to itself.
\end{lemma}
\begin{proof}
  \uses{lem:affine-coefficients-continuous}
  % \leanok
  Calculate and use Lemma~\ref{lem:affine-coefficients-continuous}.
\end{proof}

\begin{lemma}[The action of oriented affine transforms on c.d.f.s is continuous]
  \label{lem:action-on-cdf-continuous}
  \uses{lem:oriented-affine-action-on-cdf}
  \lean{CumulativeDistributionFunction.continuous_mulAction}
  \leanok
  The action $A . F$ of $A \in \OriAffR$ on a c.d.f $F$ depends
  jointly continuously on $A$ and $F$.

  \emph{(The topology on c.d.f.s is the topology of convergence in distribution, i.e.,
  convergence at all continuity points of the limit cdf.)}
\end{lemma}
\begin{proof}
  \uses{lem:affine-metrizable, lem:affine-inversion-continuous}
  % \leanok
  The spaces are metrizable, so it suffices to check sequential continuity.

  Suppose that $A_n \to B$ (oriented affine isomophisms)
  and $F_n \tendsinlaw G$ (c.d.f.s) as $n \to \infty$.
  Let $x \in \bR$ be a continuity point of $B.G$.
  Let $\eps > 0$. Note that $B^{-1}(x)$ is a continuity point of $G$.
  Then there exists a $\delta > 0$ such that $|G(y) - G\big(B^{-1}(x)\big)| < \frac{\eps}{2}$
  when $|y-B^{-1}(x)|<\delta$.

  By density of continuity points of $G$, pick continuity points
  $y_-,y_+$ such that
  \begin{align*}
    B^{-1}(x)-\delta \; < \; y_- \; < \; B^{-1}(x) \; < \; y_+ \; < \; B^{-1}(x)+\delta .
  \end{align*}
  Since $F_n \tendsinlaw G$, we have that
  $F_n(y_\pm) \to G(y_\pm)$ as $n \to \infty$,
  and in particular there exists some $N$ such that for $n \ge N$
  we have $\big| F_n(y_\pm) - G(y_\pm) \big| < \frac{\eps}{4}$ for both $y_\pm$.

  By Lemma~\ref{lem:affine-inversion-continuous}
  and $A_n \to B$, we get that $A_n^{-1} \to B^{-1}$, and in particular
  there exists some $N'$ such that for $n \ge N'$ we have
  \begin{align*}
    y_- < A_n^{-1}(x) < y_+ .
  \end{align*}

  For $n \ge \max \set{N,N'}$, we then have
  \begin{align*}
    (A_n.F_n)(x) \;
    = \; & F_n \big( A_n^{-1}(x) \big) \\
    \le \; & F_n (y_+) \\
    < \; & G(y_+) + \frac{\eps}{4} \\
    \le \; & G \big( B^{-1}(x) \big) + \frac{\eps}{2} + \frac{\eps}{4} \\
    = \; & (B.G) (x) + \frac{3\eps}{4} \\
    < \; & (B.G) (x) + \eps
  \end{align*}
  and similarly
  \begin{align*}
    (A_n.F_n)(x) \;
    = \; & F_n \big( A_n^{-1}(x) \big) \\
    \ge \; & F_n (y_-) \\
    > \; & G(y_-) - \frac{\eps}{4} \\
    \ge \; & G \big( B^{-1}(x) \big) - \frac{\eps}{2} - \frac{\eps}{4} \\
    = \; & (B.G) (x) - \frac{3\eps}{4} \\
    > \; & (B.G) (x) + \eps ,
  \end{align*}
  which together yield $\big| (A_n.F_n)(x) - (B.G)(x) \big| < \eps$.

  Since $x$ was an arbitrary continuity point of $B.G$ and
  $\eps > 0$ was arbitrary, this proves that $A_n.F \tendsinlaw B.G$.
\end{proof}

% lem:oriented-affine-action-on-cdf


\chapter{Extreme value distributions}

\section{Definition of extreme value distributions}

Lemma~\ref{lem:cdf-of-affine-max-many} motivates the following definition.

\begin{definition}
  \label{def:extr-val-distr}
  \uses{lem:oriented-affine-action-on-cdf}
  \lean{CumulativeDistributionFunction.IsExtremeValueDistr}
  \leanok
  A c.d.f. $G$ is said to be an \emph{extreme value distribution} if
  $G$ is nondegenerate and
  there exists a c.d.f. $F$ and a sequence $(A_n)_{n \in \N}$ of orientation
  preserving affine isomorphisms $A_n \in \OriAffR$, such that
  for every continuity point $x \in \R$ of~$G$ we have
  \begin{align*}
  \lim_{n \to \infty} \big( (A_n.F)(x) \big)^n \; = \; G(x) .
  \end{align*}
\end{definition}

\begin{lemma}
  \label{lem:extr-val-distr-transform}
  \uses{lem:oriented-affine-action-on-cdf, def:extr-val-distr}
  \lean{CumulativeDistributionFunction.IsExtremeValueDistr.affineTransform}
  \leanok
  Let $G$ be an extreme value distribution and
  and $A \in \OriAffR$ an orientation preserving affine isomorphism.
  Then also $A.G$ is an extreme value distribution.
\end{lemma}
\begin{proof}
  \uses{lem:cdf-continuity-pt-transform, lem:degenerate-cdf-transform}
  \leanok
  Straightforward using Lemmas~\ref{lem:cdf-continuity-pt-transform}
  and~\ref{lem:degenerate-cdf-transform}.
\end{proof}



\section{Three types of extreme value distributions}

\begin{definition}
  \label{def:std-Gumbel-cdf}
  \uses{def:cdf}
  \lean{standardGumbelCDF}
  % \leanok
  \emph{The standard Gumbel distribution} is the c.d.f. $\GumbelCDF$ given by
  \begin{align*}
  \GumbelCDF (x) = \exp\big(-\exp(-x)\big) .
  \end{align*}

  (In the parametrization of extreme value distribution types by one
  index $\gamma \in \bR$, this case corresponds to $\gamma = 0$.)
\end{definition}

\begin{definition}
  \label{def:std-Weibull-cdf}
  \uses{def:cdf}
  \lean{standardWeibullCDF}
  % \leanok
  \emph{The standard (reverse) Weibull distribution} of parameter
  $\alpha > 0$ is the c.d.f. $\WeibullCDF{\alpha}$ given by
  \begin{align*}
  \WeibullCDF{\alpha} (x) = \begin{cases}
      \exp \big( - (-x)^\alpha \big) \quad & \text{ for } x < 0 \\
      1 \quad & \text{ for } x \ge 0 .
    \end{cases}
  \end{align*}

  (In the parametrization of extreme value distribution types by one
  index $\gamma \in \bR$, this case corresponds to $\gamma < 0$
  via $\gamma = -1/\alpha$.)
\end{definition}

\begin{definition}
  \label{def:std-Frechet-cdf}
  \uses{def:cdf}
  \lean{standardFrechetCDF}
  % \leanok
  \emph{The standard Fr\'echet distribution} of parameter $\alpha > 0$
  is the c.d.f. $\FrechetCDF{\alpha}$ given by
  \begin{align*}
  \FrechetCDF{\alpha} (x) = \begin{cases}
    0 \quad & \text{ for } x \le 0 \\
    \exp \big( - x^{-\alpha} \big) \quad & \text{ for } x > 0 .
    \end{cases}
  \end{align*}

  (In the parametrization of extreme value distribution types by one
  index $\gamma \in \bR$, this case corresponds to $\gamma > 0$
  via $\gamma = 1/\alpha$.)
\end{definition}

\begin{theorem}
  \label{thm:Gumbel-is-extr-val-distr}
  \uses{def:std-Gumbel-cdf, def:extr-val-distr}
  \lean{isExtremeValueDistr_standardGumbelCDF}
  \leanok
  The standard Gumbel distribution $\GumbelCDF$ is an extreme value
  distribution.
\end{theorem}
\begin{proof}
  %\uses{lem:cdf-of-max-many}
  % \leanok
  Set $A_n(x) = x - \log(n)$ for $n \in \bN$. Then
  $A_n^{-1}(x) = x + \log(n)$ and for any $n \ge 1$ and $x \in \bR$ we get
  \begin{align*}
    \big( (A_n . \GumbelCDF)(x) \big)^n
    = \; & \big( \GumbelCDF(x + \log(n)) \big)^n \\
    = \; & \Big( \exp \big( - \exp(-(x + \log n)) \big) \Big)^n \\
    = \; & \exp \big( - n \exp(-x - \log n) \big) \\
    = \; & \exp \big( - n \, e^{-x} \, e^{-\log n}  \big) \\
    %= \; & \exp \big( - n \, e^{-x} \, n^{-1} \big) \\
    = \; & \exp \big( - e^{-x} \big) \\
    = \; & \GumbelCDF(x) .
  \end{align*}
  Since the above is true for each $n$, we in particular have
  \begin{align*}
    \lim_{n \to \infty} \big( (A_n . \GumbelCDF)(x) \big)^n
    = \GumbelCDF(x)
  \end{align*}
  for all $x \in \bR$. Since $\GumbelCDF$ is also nondegenerate,
  this shows that it is an extreme value distribution.
\end{proof}

\begin{theorem}
  \label{thm:Weibull-is-extr-val-distr}
  \uses{def:std-Weibull-cdf, def:extr-val-distr}
  \lean{isExtremeValueDistr_standardWeibullCDF}
  % \leanok
  For any $\alpha > 0$, the standard Weibull distribution $\WeibullCDF{\alpha}$
  is an extreme value distribution.
\end{theorem}
\begin{proof}
  %\uses{lem:cdf-of-max-many}
  % \leanok
  \ldots
\end{proof}

\begin{theorem}
  \label{thm:Frechet-is-extr-val-distr}
  \uses{def:std-Frechet-cdf, def:extr-val-distr}
  \lean{isExtremeValueDistr_standardFrechetCDF}
  % \leanok
  For any $\alpha > 0$, the standard Fr\'echet distribution $\FrechetCDF{\alpha}$
  is an extreme value distribution.
\end{theorem}
\begin{proof}
  \uses{}
  % \leanok
  \ldots
\end{proof}



\section{Equivalent formulations of the limit relation}

\begin{lemma}[Logarithmic version of the limit relation]
  \label{lem:log-ev-limit}
  \uses{def:cdf, lem:oriented-affine-action-on-cdf}
  \lean{ev_limit_iff_log_ev_limit}
  \leanok
  Let $F$ and $G$ be c.d.f.s, and $(A_n)_{n \in \N}$ a sequence of orientation
  preserving affine isomorphisms $A_n \in \OriAffR$. Then for any $x \in \bR$
  such that $0 < G(x) < 1$, the two conditions
  \begin{align*}
  \text{(i)} & & \lim_{n \to \infty} \big( (A_n.F)(x) \big)^n \, = \; & G(x) \\
  \text{(ii)} & & \lim_{n \to \infty} n \, \log F (A_n^{-1}(x)) \, = \; & \log G(x)
  \end{align*}
  are equivalent.
\end{lemma}
\begin{proof}
  \uses{}
  \leanok
  Recall that $(A_n.F)(x) = F (A_n^{-1}(x))$.
  Then just take logarithms (and use continuity) to get from (i) to (ii),
  and take exponentials (and use continuity) to get from (ii) to (i).
\end{proof}

\begin{lemma}[Relation implies $F$ tending to one]
  \label{lem:ev-limit-cdf-affine-tendsto-one}
  \uses{def:cdf, lem:oriented-affine-action-on-cdf}
  \lean{tendsto_one_of_ev_limit}
  \leanok
  Let $F$ and $G$ be c.d.f.s, and $(A_n)_{n \in \N}$ a sequence of orientation
  preserving affine isomorphisms $A_n \in \OriAffR$. Then for any $x \in \bR$
  such that $0 < G(x) < 1$, if
  \begin{align*}
  \text{(i)} & & \lim_{n \to \infty} \big( (A_n.F)(x) \big)^n \, = \; & G(x)
  \end{align*}
  holds, then necessarily
  \begin{align*}
    \lim_{n \to \infty} F (A_n^{-1}(x)) \; = \; 1 .
  \end{align*}
\end{lemma}
\begin{proof}
  % \uses{}
  % \leanok
  Otherwise $\big( F (A_n^{-1}(x)) \big)^n$ would have $0 \ne G(x)$
  as an accumulation point, contradicting the assumed limit (i).

  To wit, if for some $\delta > 0$ we would have
  $F (A_n^{-1}(x)) \le 1 - \delta$ for infinitely many $n$, then
  $0 \le \big( F (A_n^{-1}(x)) \big)^n \le (1 - \delta)^n $
  for those $n$, and since $(1 - \delta)^n \to 0$, we would
  get $\big( F (A_n^{-1}(x)) \big)^n \to 0 \ne G(x)$ along the
  subsequence of those $n$; a contradiction.
\end{proof}

\begin{lemma}[Taylor expansion limit modification]
  \label{lem:modify-limit-taylor}
  \lean{tendsto_smul_apply_smul_deriv_of_tendsto_atTop_of_tendsto_smul_apply_smul_deriv,
  tendsto_zero_of_tendsto_atTop_of_tendsto_smul}
  \leanok
  Let $S \subset \bR$ be a subset with $0 \in S$, and let $f_1, f_2 \colon S \to \bR$ be
  functions. Let also $(t_n)_{n \in \bN}$ be a sequence in $S$,
  and let $(m_n)_{n \in \bR}$ be a sequence of real numbers tending to infinity,
  $\lim_{n \to \infty} m_n = +\infty$.
  Assume that for $j = 1,2$
  \begin{itemize}
    \item $f_j(0)=0$;
    \item the derivative $f_j'(0)$ exists (derivative taken within the set $S$)
  \end{itemize}
  Assume further about $j = 1$ that
  \begin{itemize}
    \item $f_1'(0) \ne 0$;
    \item the limit $\lim_{n \to \infty} \big( m_n \, f_1 (t_n) \big)$ exists;
    \item for any $\delta > 0$ there exists an $\eps > 0$ such that
      if $|f_1(t)| < \eps$ (with $t \in S$) then $|t| < \delta$.
  \end{itemize}
  Denote $c = \frac{1}{f_1'(0)} \, \lim_{n \to \infty} \big( m_n \, f_1 (t_n) \big)$.
  Then we have $\lim_{n \to 0} t_n = 0$ and
  \begin{align*}
    \lim_{n \to \infty} \Big( m_n \, f_2 (t_n) \Big) \; = \; & c \, f_2'(0) .
  \end{align*}
\end{lemma}
\begin{proof}
  % \uses{}
  \leanok
  This is in principle straightforward:
  the assumptions are first checked to imply that $\lim_{n \to 0} t_n = 0$,
  and then one can just consider the first order Taylor expansions of the
  functions $f_j$ at $0$ given by the assumed existence of the derivatives
  (the key is $t_n \, = \, \frac{c}{m_n} + \oo \big(\frac{1}{m_n}\big)$).
\end{proof}

\begin{lemma}[Taylored version of the limit relation]
  \label{lem:taylored-ev-limit}
  \uses{def:cdf, lem:oriented-affine-action-on-cdf}
  \lean{log_ev_limit_iff_taylored_ev_limit}
  \leanok
  Let $F$ and $G$ be c.d.f.s, and $(A_n)_{n \in \N}$ a sequence of orientation
  preserving affine isomorphisms $A_n \in \OriAffR$. Then for any $x \in \bR$
  such that $0 < G(x) < 1$, the two conditions
  \begin{align*}
    \text{(ii)} & & \lim_{n \to \infty} n \, \log F (A_n^{-1}(x)) \, = \; & \log G(x) \\
    \text{(iii)} & & \lim_{n \to \infty} \Big( n \, \big(1 - F (A_n^{-1}(x)) \big) \Big)
      \, = \; & - \log G(x)
  \end{align*}
  are equivalent.
\end{lemma}
\begin{proof}
  \uses{lem:modify-limit-taylor}
  % \leanok

  Both implications {(ii) $\, \Rightarrow \,$ (iii)} and {(iii) $\, \Rightarrow \,$ (ii)}
  are proven similarly using Lemma~\ref{lem:modify-limit-taylor}.

  Assume (ii). Let $f_1(t) = - \log (1-t)$ and $f_2(t) = t$ and
  $S = [0,1) \subset \bR$, and $m_n = n$ and $t_n = 1 - F (A_n^{-1}(x))$.
  It is straightforward to check the assumptions of Lemma~\ref{lem:modify-limit-taylor},
  with
  \begin{align*}
    f_1'(0) \; = \; & \der{t} \Big|_{t=0} \Big( - \log (1 - t) \Big) \; = \; 1
  \end{align*}
  and $f_2'(0) = \mathrm{id}'(0) = 1$.
  The key assumption about the existence of
  the limit $\lim_{n \to \infty} \big( m_n \, f_1(t_n) \big)$ is given by (ii),
  and the conclusion is (iii).

  Similarly assuming (iii) we derive (ii) with Lemma~\ref{lem:modify-limit-taylor}
  just interchanging the roles of the two functions, i.e., now using
  $f_1(t) = t$ and $f_2(t) = - \log (1-t)$ instead.
\end{proof}

\begin{lemma}[Inverted Taylored version of the limit relation]
  \label{lem:inv-taylored-ev-limit}
  \uses{def:cdf, lem:oriented-affine-action-on-cdf}
  \lean{taylored_ev_limit_iff_oneDivOneSub_limit}
  \leanok
  Let $F$ and $G$ be c.d.f.s, and $(A_n)_{n \in \N}$ a sequence of orientation
  preserving affine isomorphisms $A_n \in \OriAffR$. Then for any $x \in \bR$
  such that $0 < G(x) < 1$, the two conditions
  \begin{align*}
    \text{(iii)} & & \lim_{n \to \infty} \Big( n \, \big(1 - F (A_n^{-1}(x)) \big) \Big)
      \, = \; & - \log G(x) \\
    \text{(iv)} & & \lim_{n \to \infty} \frac{1}{n \, \big(1 - F (A_n^{-1}(x)) \big)}
      \, = \; & \frac{1}{-\log G(x)}
  \end{align*}
  are equivalent.
\end{lemma}
\begin{proof}
  % \uses{}
  \leanok
  By assumption $G(x) \in (0,1)$ we have $-\log G(x) > 0$.

  The implication {(iii)~$\,\Rightarrow \,$~(iv)}
  can therefore be seen by applying $t \mapsto \frac{1}{t}$
  and using its continuity at $t = -\log G(x)$, and the
  converse implication {(iv)~$\,\Rightarrow \,$~(iii)} similarly by using continuity of
  $t \mapsto \frac{1}{t}$ at $t = \frac{1}{-\log G(x)}$.
\end{proof}

\begin{lemma}[Transformed version of the limit relation]
  \label{lem:transformed-ev-limit}
  \uses{def:one-div-one-sub-cdf, def:one-div-neg-log-cdf, lem:oriented-affine-action-on-cdf}
  \lean{oneDivSub_limit_iff}
  \leanok
  Let $F$ and $G$ be c.d.f.s, and $(A_n)_{n \in \N}$ a sequence of orientation
  preserving affine isomorphisms $A_n \in \OriAffR$. Then for any $x \in \bR$
  such that $0 < G(x) < 1$, the two conditions
  \begin{align*}
    \text{(iv)} & & \lim_{n \to \infty} \frac{1}{n \, \big(1 - F (A_n^{-1}(x)) \big)}
      \, = \; & \frac{1}{-\log G(x)} \\
    \text{(v)} & & \lim_{n \to \infty} \frac{1}{n} \oneDivOneSub{A_n . F} (x)
      \, = \; & \oneDivNegLog{G} (x)
  \end{align*}
  are equivalent.

  (See Definitions~\ref{def:one-div-one-sub-cdf}
  and \ref{def:one-div-neg-log-cdf} for the transforms involved in condition (v)).
\end{lemma}
\begin{proof}
  % \uses{}
  \leanok
  This is in principle straightforward, although certain cases need to be checked
  separately and the continuity of various natural extensions need to addressed.
\end{proof}

% def:one-div-neg-log-cdf


\begin{theorem}[Equivalent versions of the limit relation]
  \label{thm:tfae-ev-limit}
  \uses{def:one-div-one-sub-cdf, def:one-div-neg-log-cdf, lem:oriented-affine-action-on-cdf}
  \lean{tfae_ev_limit_relation}
  \leanok
  Let $F$ and $G$ be c.d.f.s, and $(A_n)_{n \in \N}$ a sequence of orientation
  preserving affine isomorphisms $A_n \in \OriAffR$. Then for any $x \in \bR$
  such that $0 < G(x) < 1$, the conditions
  \begin{align*}
    \text{(i)} & & \lim_{n \to \infty} \big( (A_n.F)(x) \big)^n \, = \; & G(x) \\
    \text{(ii)} & & \lim_{n \to \infty} n \, \log F (A_n^{-1}(x)) \, = \; & \log G(x) \\
    \text{(iii)} & & \lim_{n \to \infty} \Big( n \, \big(1 - F (A_n^{-1}(x)) \big) \Big)
      \, = \; & - \log G(x) \\
    \text{(iv)} & & \lim_{n \to \infty} \frac{1}{n \, \big(1 - F (A_n^{-1}(x)) \big)}
      \, = \; & \frac{1}{-\log G(x)} \\
    \text{(v)} & & \lim_{n \to \infty} \frac{1}{n} \oneDivOneSub{A_n . F} (x)
      \, = \; & \oneDivNegLog{G} (x)
  \end{align*}
  are equivalent.

  (See Definitions~\ref{def:one-div-one-sub-cdf}
  and \ref{def:one-div-neg-log-cdf} for the transforms involved in condition (v)).

  \emph{(More equivalent conditions are to be added; this is just a theorem to collect
  various equivalent phrasings.)}
\end{theorem}
\begin{proof}
  \uses{lem:taylored-ev-limit, lem:log-ev-limit, lem:inv-taylored-ev-limit,
  lem:transformed-ev-limit}
  \leanok
  This is just a combination of earlier results.
\end{proof}


\chapter{Classification of extreme value distributions}

\section{Auxiliary classification I}

\begin{lemma}[Second order differential equation for Q]
  \label{lem:ode-of-order-two-for-Q}
  %\uses{}
  \lean{ExtremeValueTheory.ode_of_order_two_for_Q}
  \leanok
  Suppose that
  \begin{align*}
  Q \colon \bR \to \bR
  \end{align*}
  is differentiable and satisfies
  \begin{align*}
  Q(0) = 0 \qquad \text{ and } \qquad Q'(0) = 1
  \end{align*}
  and
  \begin{align*}
  Q(h+s) = Q(h) \alpha(s) + Q(s)
  \end{align*}
  for some $\alpha \colon \bR \to \bR$ and every $s,h \in \bR$.
  Then $Q$ is twice continuously differentiable and satisfies
  \begin{align}
    Q''(s) = Q'(s) \, Q''(0) \qquad \text{ for every } s \in \bR .
  \end{align}
\end{lemma}
\begin{proof}
  %\uses{}
  % \leanok
  Note first that the equation implies (rearranging and dividing by $h$),
  for any $s$ and $h \ne 0$,
  \begin{align*}
    \frac{Q(h+s) - Q(s)}{h} = \frac{Q(h)}{h} \alpha(s) .
  \end{align*}
  Taking the limit as $h \to 0$ and using $Q(0)=0$ yields
  $Q'(s) = Q'(0) \, \alpha(s) = \alpha(s)$, where we also
  took into account $Q'(0) = 1$. Therefore necessarily
  $\alpha = Q'$, and the equation can be rewritten in the form
  \begin{align*}
  Q(h+s) = Q(h) \, Q'(s) + Q(s) \qquad \text{ for } s,h \in \bR .
  \end{align*}

  Rearranging the equation, we find
  \begin{align*}
    Q(h+s) - Q(s) = Q(h) \, Q'(s)
  \end{align*}
  and interchanging the role of $s$ and $h$ also
  \begin{align*}
    Q(h+s) - Q(h) = Q(s) \, Q'(h) .
  \end{align*}
  Subtracting the last two equations yields
  \begin{align*}
    Q(s) - Q(h) = Q(s) \, Q'(h) - Q(h) \, Q'(s) ,
  \end{align*}
  which by rearranging and dividing by $h \ne 0$ yields
  \begin{align*}
     \frac{Q(h)}{h} \big( Q'(s) - 1 \big) = Q(s) \frac{Q'(h) - 1}{h} .
  \end{align*}
  Taking the limit $h \to 0$, recalling $Q(0)=1$ and $Q'(0)=1$,
  gives the existence of the second derivative $Q''(0)$ and the equation
  \begin{align*}
    Q'(s) - 1 = Q(s) Q''(0) .
  \end{align*}
  Solving $Q'(s) = 1 + Q''(0) \, Q(s)$ and recalling that $Q$
  is differentiable shows that $Q'$ is also differentiable, so
  $Q$ is indeed twice differentiable. Differentiating, we get
  the asserted equation
  \begin{align*}
    Q''(s) = Q'(s) \, Q''(0) .
  \end{align*}
  Since $Q'$ is differentiable and in particular continuous, this also shows
  that $Q''$ is continuous, i.e., that $Q$ is twice continuously differentiable.
\end{proof}

\begin{lemma}[Solution for Q]
  \label{lem:solve-Q}
  %\uses{}
  \lean{ExtremeValueTheory.solve_Q}
  \leanok
  Suppose that $Q \colon \bR \to \bR$ is twice continuously differentiable
  and $Q'$ is positive and $Q$
  and satisfies $Q(0)=0$, $Q'(0) = 1$, and the equation concluded in
  Lemma~\ref{lem:ode-of-order-two-for-Q} with $\gamma = Q''(0)$, i.e.,
  \begin{align}\label{eq: Q eqn dd}
  Q''(s) = \gamma \, Q'(s) \qquad \text{ for every } s \in \bR .
  \end{align}
  Then $Q$ is given by
  \begin{align*}
  Q(s) = \begin{cases}
    \frac{e^{\gamma s} - 1}{\gamma} & \; \text{ if $\gamma \ne 0$} \\
    \;\; s & \; \text{ if $\gamma = 0$}
    \end{cases}
    \qquad \text{for $s \in \bR$.}
  \end{align*}
%  if $\gamma \ne 0$ and
%  \begin{align*}
%  Q(s) = s
%  \end{align*}
%  if $\gamma = 0$.
  \end{lemma}
\begin{proof}
  %\uses{}
  % \leanok
  Since $Q$ is differentiable and $Q'(s)>0$ for any $s \in \bR$,
  we can write \eqref{eq: Q eqn dd} as
  \begin{align*}
  \der{s} \log Q'(s)
    = \frac{Q''(s)}{Q'(s)}
    = \gamma .
  \end{align*}
  Integrating and noting $\log Q'(0) = \log 1 = 0$, this
  yields $\log Q'(s) = \gamma s$ for all $s \in \bR$, or equivalently
  $Q'(s) = e^{\gamma s}$. Integrating a second time
  and noting $Q(0) = 0$, this yields the desired formulas
  in the two cases $\gamma \ne 0$ and $\gamma = 0$, respectively, as follows.

  If $\gamma \ne 0$, then we get
  \begin{align*}
  Q(s) = \int_0^s e^{\gamma u} \, \ud u = \frac{e^{\gamma s} - 1}{\gamma} .
  \end{align*}

  If $\gamma = 0$, then we get
  \begin{align*}
  Q(s) = \int_0^s e^0 \, \ud u = s .
  \end{align*}
\end{proof}

% TODO: Move to a different file.
\begin{theorem}[Monotone functions are a.e. differentiable]
  \label{thm:monotone-ae-differentiable}
  %\uses{}
  \lean{Monotone.ae_differentiableAt}
  \leanok
  \mathlibok
  If $f : \bR \to \bR$ is nondecreasing, then the derivative
  $f'(x)$ exists at almost every $x \in \bR$.
  In particular there exists points $x$ where $f'(x)$ exists.
\end{theorem}
\begin{proof}
  %\uses{}
  \leanok
  \mathlibok
  (The proof is already in Mathlib.)
\end{proof}

\begin{lemma}[Solution for E]
  \label{lem:solve-E}
  %\uses{lem:solve-Q}
  \lean{ExtremeValueTheory.solve_E}
   \leanok
  Suppose that $E \colon (0,\infty) \to \bR$
  is nondecreasing and nonconstant function which satisfies $E(1) = 0$ and
  \begin{align*}
  E(\lambda \sigma) = E(\lambda) A(\sigma) + E(\sigma)
  \end{align*}
  for some $A \colon (0,\infty) \to (0,\infty)$ and all $\lambda, \sigma > 0$.
  Then, denoting $c = E'(1)$ and $\gamma = \frac{1}{E'(1)} \dder{s} E(e^s) \big|_{s=0}$,
  for all $\lambda \in \bR$ we have
  \begin{align*}
  E(\lambda) = \begin{cases}
    c \big( \lambda^\gamma - 1 \big) & \;\;\text{ if $\gamma \ne 0$.} \\
    c \, \log (\lambda) & \;\;\text{ if $\gamma = 0$.}
    \end{cases}
  \end{align*}
\end{lemma}
\begin{proof}
  \uses{lem:ode-of-order-two-for-Q, lem:solve-Q, thm:monotone-ae-differentiable}
  % \leanok
  Denote $H(s) = E(e^s)$ for $s \in \bR$. Then $H(0) = E(1) = 0$ and
  $H$ is also nondecreasing and nonconstant.
  By Theorem~\ref{thm:monotone-ae-differentiable}, there
  exists some $s_0 \in \bR$ such that the derivative $H'(s_0)$ exists.

  The equation for $E$ yields an equation for $H$,
  \begin{align*}
  H(h+s) = H(h) A(e^{s}) + H(s) \qquad \text{ for any } s, h \in \bR.
  \end{align*}
  We can rearrange this equation and divide by $h$,
  and we get for any $s, h \in \bR$, $h \ne 0$,
  \begin{align*}
  \frac{H(h+s) - H(s)}{h} = \frac{H(h)}{h} A(e^{s}) .
  \end{align*}
  If $s = s_0$, then the LHS tends to $H'(s_0)$ as $h \to 0$.
  The RHS must therefore also have a limit as $s \to 0$, and
  observing that $\frac{H(h)}{h} = \frac{H(h) - H(0)}{h}$,
  that limit is $H'(0) A(e^{s_0})$, which shows that the derivative $H'(0)$ exists.
  Then applying the same equation at general $s \in \bR$ shows that
  $H'(s) = H'(0) A(e^{s})$ must exist, so $H \colon \bR \to \bR$ is
  in fact everywhere differentiable.

  Note that we must have $H'(0) > 0$, because if $H'(0) = 0$
  then by the above equation $H'(s) = 0$ for every $s$ and then
  $H$ is a constant function, which is a contradiction.
  Denote $c = H'(0) = \der{s} E(e^s) \big|_{s=0} = E'(1)$.

  Also the equation and positivity of the function $A$ give $H'(s)>0$
  for all $s$. %, so $H$ is strictly increasing.

  Now consider $Q(s) = \frac{H(s)}{c}$. This $H$ is
  obviously also differentiable, since $H$ is. The equation
  for $H$ yields the following equation for $Q$,
  \begin{align*}
    Q(h+s) = Q(h) \alpha(s) + Q(s) \qquad \text{ for any } s,h \in \bR ,
  \end{align*}
  where $\alpha(s) = A(e^{s})$.
  By Lemma~\ref{lem:ode-of-order-two-for-Q} $Q$ is then twice
  continuously differentiable and satisfies
  \begin{align}
  Q'(s) Q'' (0) = Q'' (s) \qquad \text{ for all } s \in \bR .
  \end{align}
  By Lemma~\ref{lem:solve-Q} we then get
  a formula for $Q$, which involves
  \begin{align*}
  \gamma = Q''(0) = \frac{H''(0)}{c}
    = \frac{1}{c} \dder{s} E(e^s) \big|_{s = 0} .
  \end{align*}

  If $\gamma \ne 0$ then the formula reads
  $Q(s) = \frac{e^{\gamma s} - 1}{\gamma}$.
  Now just tracing the definitions, we get the asserted formula
  \begin{align*}
  E(\lambda) = H(\log \lambda)
    = c \, Q(\log \lambda) = c \, \big( e^{\gamma \, \log(\lambda)} - 1 \big)
    = c \, \big( \lambda^{\gamma} - 1 \big) .
  \end{align*}

  If $\gamma = 0$ then the formula reads $Q(s) = s$.
  Now just tracing the definitions, we get instead
  \begin{align*}
  E(\lambda) = H(\log \lambda)
    = c \, Q(\log \lambda) = c \, \log(\lambda) ,
  \end{align*}
  and the proof is complete.
\end{proof}


\section{Inverting to recover a cumulative distribution function}

%\begin{lemma}[Classification of cdfs of evds]
%  \label{lem:cdf-of-evd-classification}
%  %\uses{}
%  %\lean{}
%  % \leanok
%  Suppose that $G \colon \bR \to [0,1]$ is a c.d.f.
%  whose right-continuous inverse is given by $G^\leftarrow = \cdots$
%\end{lemma}


\chapter{Convergence in distribution with cdfs}

This chapter provides a standard characterization of convergence in distribution
(weak convergence of probability measures) on the real line in terms of
cumulative distribution functions.

\section{Convergence in distribution}

Convergence in distribution for random variables can be defined when the random variables
take values in a topological space, and it amounts to the weak convergence of the
probability measures that are the laws of those random variables. In the special case
of real-valued random variables, or probability measures on the real line, the definition
reads:
\begin{definition}[Weak convergence of probability measures]
  \label{def:convergence-in-distribution}
  %\uses{}
  \lean{MeasureTheory.ProbabilityMeasure.instTopologicalSpace,
  MeasureTheory.ProbabilityMeasure.tendsto_iff_forall_lintegral_tendsto}
  \leanok
  \mathlibok
  A sequence $(\PRmeas_n)_{n \in \bN}$ of Borel probability measures on $\bR$ converges
  weakly to a Borel probability measure $\PRmeas$ on $\bR$ if for all bounded continuous
  functions $f \colon \bR \to [0,+\infty)$ we have
  \begin{align*}
    \lim_{n \to \infty} \int_{\bR} f(x) \, \ud \PRmeas_n(x)
    = \int_{\bR} f(x) \, \ud \PRmeas(x) .
  \end{align*}
\end{definition}

\section{Auxiliary results}

\begin{lemma}[Monotone real functions have only countably many points of discontinuity]
  \label{lem:monotone-discontinuities-countable}
  %\uses{}
  \lean{Set.Countable.dense_compl, Monotone.countable_not_continuousAt}
  \leanok
  \mathlibok
  A monotone function $f \colon \bR \to \bR$ can have at most countably
  many points of discontinuity. In particular the set $D \subset \bR$
  of continuity points of $f$ is dense in $\bR$.
\end{lemma}
\begin{proof}
  %\uses{}
  \leanok
  \mathlibok
  (The proof should already be in Mathlib.)
\end{proof}

\begin{lemma}[Tightness of a cumulative distribution function]
  \label{lem:cdf-tight}
  \uses{def:cdf}
  \lean{CumulativeDistributionFunction.forall_pos_exists_lt_gt_continuousAt}
  \leanok
  Let $F$ be a cumulative distribution function. Then for any
  $\eps > 0$ there exists points $a,b \in \bR$ with $a < b$ such that
  $F(b) - F(a) > 1 - \eps$ and $F$ is continuous at the points $a$ and $b$.
\end{lemma}
\begin{proof}
  \uses{def:cdf, lem:monotone-discontinuities-countable}
  % \leanok
  Cumulative distribution functions satisfy
  $F(x) \downarrow 0$ as $x \downarrow - \infty$ and
  $F(x) \uparrow 1$ as $x \uparrow + \infty$. The
  required large difference $F(b) - F(a)$ is
  obtained by choosing $a$ small enough so that $F(a) < \frac{\eps}{2}$
  and $b$ large enough so that $F(b) > 1 - \frac{\eps}{2}$.
  In order to guarantee that $a < b$ and
  that $a$ and $b$ are continuity points of $F$, we
  recall that continuity points of the monotone function $F$ are dense
  by Lemma~\ref{lem:monotone-discontinuities-countable}, so we
  may decrease $a$ and increase $b$ as appropriate.
\end{proof}

\begin{lemma}[Subdivision with small mesh and within dense set]
  \label{lem:subdivision-dense}
  %\uses{}
  \lean{forall_exists_subdivision_diff_lt_of_dense}
  \leanok
  Let $D \subset \bR$ be a dense set and $a,b \in D$ with $a < b$.
  Then for any $\delta > 0$ there exists a $k \in \bN$
  and $a = c_0, c_1, \ldots, c_{k-1}, c_k = b \in D$ such that
  $|c_j - c_{j-1}| < \delta$ for all $j=1,\ldots,k$.
\end{lemma}
\begin{proof}
  %\uses{}
  % \leanok
  \ldots
\end{proof}

\begin{lemma}[Subdivision for continuous function approximation]
  \label{lem:continuous-function-approximation-subdivision}
  %\uses{}
  \lean{forall_exists_subdivision_dist_apply_lt_of_dense_of_continuous}
  \leanok
  Let $D \subset \bR$ be a dense set, let $f \colon \bR \to \bR$ be continuous,
  let $a, b \in D$ with $a < b$, and let $\eps > 0$.
  Then there exists a $k \in \bN$ and points
  $a=c_0 < c_1 < \cdots < c_{k-1} < c_k = b$ such that for
  each $j = 1, \ldots, k$ we have $c_j \in D$ and
  \begin{align*}
  \big| f(x) - f(c_j) \big| < \eps
  \qquad \text{ for } \qquad x \in [c_{j-1} , c_j] .
  \end{align*}
\end{lemma}
\begin{proof}
  \uses{lem:subdivision-dense}
  % \leanok
  On the compact interval
  $[a,b] \subset \bR$, the continuous function~$f$ is uniformly continuous, so for some
  $\delta > 0 $ we have $|f(x)-f(y)|<\eps$ whenever $|x-y|<\delta$ and $x,y \in
  [a,b]$.
  Now apply Lemma~\ref{lem:subdivision-dense} to choose $k$ and points
  $a=c_0 < c_1 < \cdots < c_{k-1} < c_k$ such that $c_j - c_{j-1} < \delta$
  and $c_j \in D$ for all $j = 1 , \ldots , k$.
  Now for any $j = 1 , \ldots , k$, since for $x \in [c_{j-1} , c_j]$ we have
  $|x - c_j| < \delta$, we get
  \begin{align*}
  \big| f(x) - f(c_j) \big| < \eps
  \end{align*}
  as desired.
\end{proof}

\begin{lemma}[Simple function integral as linear combination of cdf differences]
  \label{lem:simple-integral-cdf-difference}
  \uses{def:cdf}
  \lean{CumulativeDistributionFunction.integral_sum_indicator_eq}
  \leanok
  Let $a = c_0 < c_1 < \cdots < c_k = b$ and consider the
  linear combination of indicator functions
  \begin{align*}
    h(x) = \sum_{j=1}^k \alpha_j \; \indof{(c_{j-1},c_j]}(x) .
  \end{align*}
  Then the integral of $h$ with respect to a Borel probability measure $\PRmeas$ on $\bR$ whose
  can be written as
  \begin{align*}
    \int_{\bR} h(x) \, \ud \PRmeas (x) = \sum_{j=1}^k \alpha_j \, \big( F(c_j) - F(c_{j-1}) \big) ,
  \end{align*}
  where $F$ is the c.d.f. of $\PRmeas$.
\end{lemma}
\begin{proof}
  % \uses{}
  % \leanok
  \begin{align*}
  \int_{\bR} h \, \ud \PRmeas
   = \; & \int_{\bR}
          \Big( \sum_{j=1}^k \alpha_j \, \indof{(c_{j-1},c_j]}(x) \Big) \, \ud \PRmeas(x) \\
   = \; & \sum_{j=1}^k \alpha_j \; \int_{\bR} \indof{(c_{j-1},c_j]}(x) \, \ud \PRmeas(x) \\
   = \; & \sum_{j=1}^k \alpha_j \; \PRmeas \big[ (c_{j-1},c_j] \big] \\
   = \; & \sum_{j=1}^k \alpha_j \, \big( F_n(c_j) - F_n(c_{j-1}) \big)
  \end{align*}
\end{proof}

\begin{lemma}[One of the portmanteau implications]
  \label{lem:portmanteau-convergence-for-borel}
  \uses{def:convergence-in-distribution}
  \lean{MeasureTheory.ProbabilityMeasure.tendsto_measure_of_null_frontier_of_tendsto',
  MeasureTheory.ProbabilityMeasure.tendsto_measure_of_null_frontier_of_tendsto}
  \leanok
  \mathlibok
  Weak convergence of probability measures implies that if the boundary of a Borel set carries
  no probability mass under the limit measure, then the limit of the measures of the set equals
  the measure of the set under the limit probability measure.

  In other words, if $\lim_{n \to \infty} \PRmeas_n = \PRmeas$ in the sense of weak convergence
  of measures, Definition~\ref{def:convergence-in-distribution}, and if $A \subset \bR$
  is a Borel set such that $\PRmeas[\partial A] = 0$, then
  \begin{align*}
    \lim_{n \to \infty} \PRmeas_n [ A ] = \PRmeas [ A ] .
  \end{align*}
\end{lemma}
\begin{proof}
  % \uses{}
  \leanok
  \mathlibok
  (The proof is in Mathlib.)
\end{proof}

\section{Convergence in distribution from pointwise convergence of cdfs}

\begin{theorem}[Sufficient condition for convergence in distribution with cdfs]
  \label{thm:convergence-in-distribution-with-cdf}
  \uses{def:convergence-in-distribution, def:cdf}
  \lean{tendsto_of_forall_continuousAt_tendsto_cdf}
  \leanok
  Let $F$ and $F_n$, $n \in \bN$, be cumulative distribution functions of
  probability measures $\PRmeas$ and $\PRmeas_n$, $n \in \bN$, respectively, i.e.,
  \begin{align*}
      F(x) = \; & \PRmeas \big[(-\infty,x]\big] & & \text{for $x \in \bR$} \\
      F_n(x) = \; & \PRmeas_n \big[(-\infty,x]\big] & & \text{for $x \in \bR$ and $n \in \bN$.}
  \end{align*}
  If $\lim_{n \to \infty} F_n(x) = F(x)$ for all continuity points $x$ of $F$,
  then $\lim_{n \to \infty} \PRmeas_n = \PRmeas$ in the sense of weak convergence
  of measures, Definition~\ref{def:convergence-in-distribution}.
\end{theorem}
\begin{proof}
  \uses{lem:monotone-discontinuities-countable, lem:cdf-tight,
  lem:continuous-function-approximation-subdivision,
  lem:simple-integral-cdf-difference}
  % \leanok
  Let $D \subset \bR$ denote the set of continuity points of $F$. By
  Lemma~\ref{lem:monotone-discontinuities-countable}, $D$ is dense in $\bR$.
  Assume that $\lim_{n \to \infty} F_n(x) = F(x)$ for all $x \in D$.

  Let~$\eps > 0$. Choose, by Lemma~\ref{lem:cdf-tight}, points
  $a,b \in D$, $a<b$, such that $F(b) - F(a) > 1 - \eps$.

  Observe also that since $\lim_{n \to \infty} F_n(a) = F(a)$
  and $\lim_{n \to \infty} F_n(b) = F(b)$, there exists some $N_1$ such that
  we have
  \begin{align*}
  F_n(b) - F_n(a) > 1 - 2\eps \qquad \text{ for all } n \geq N_1 .
  \end{align*}

  Let $f \colon \bR \to \bR$ be bounded and continuous.
  By Lemma~\ref{lem:continuous-function-approximation-subdivision}
  we can choose points
  $a=c_0 < c_1 < \cdots < c_{k-1} < c_k = b$ such that
  for all $j = 1 , \ldots , k$ we have $c_j \in D$ and
  \begin{align*}
  \big| f(x) - f(c_j) \big| < \eps
  \qquad \text{ for } \qquad x \in [c_{j-1} , c_j] .
  \end{align*}
  Define the simple function $h \colon \bR \to \bR$ by
  \begin{align*}
  h(x) = \sum_{j=1}^k f(c_j) \; \indof{(c_{j-1},c_j]}(x)
  \end{align*}
  The above estimate shows that $|f(x) - h(x)| < \eps$ for all $x \in [a,b]$.
  By boundedness of~$f$, there exists a constant $K>0$ such that $|f(x)| \leq K$
  for all $x \in \bR$. Since $h$ vanishes outside $(a,b]$,
  the triangle inequality for integral with respect to $\PRmeas_n$ gives
  \begin{align*}
  \Big| \int_{\bR} f \, \ud \PRmeas_n - \int_{\bR} h \, \ud \PRmeas_n \Big|
  \, \leq \; & \, \underbrace{\int_{(a,b]} |f-h| \, \ud \PRmeas_n}_{\leq \eps}
      + \underbrace{\int_{\bR \setminus (a,b]} |f| \, \ud \PRmeas_n}_{%
      \leq K \, \PRmeas_n\big[ \bR \setminus (a,b] \big]} .
  \end{align*}
  When $n \geq N_1$, we have
  $\PRmeas_n\big[ \bR \setminus (a,b] \big]
  = 1 - \PRmeas_n\big[ (a,b] \big] = 1 - (F_n(b) - F_n(a)) < 2 \eps$,
  and thus the triangle inequality implies
  \begin{align*}
  \Big| \int_{\bR} f \, \ud \PRmeas_n - \int_{\bR} h \, \ud \PRmeas_n \Big|
  \, \leq \; & \eps + K \, 2 \eps = (1 + 2K) \, \eps .
  \end{align*}
  Similarly, integrating now with respect to $\PRmeas$ instead, one shows that
  \begin{align*}
  \Big| \int_{\bR} f \, \ud \PRmeas - \int_{\bR} h \, \ud \PRmeas \Big|
  \, \leq \; & (1 + K) \, \eps .
  \end{align*}
  It remains to consider the integrals of the function~$h$ with respect to both
  $\PRmeas_n$ and $\PRmeas$. By Lemma~\ref{lem:simple-integral-cdf-difference},
  these integrals are expressible in terms of the cumulative distribution functions,
  \begin{align*}
  \int_{\bR} h \, \ud \PRmeas_n
  = \; & \sum_{j=1}^k f(c_j) \, \big( F_n(c_j) - F_n(c_{j-1}) \big)
  \end{align*}
  and
  \begin{align*}
  \int_{\bR} h \, \ud \PRmeas
  % = \; & \sum_{j=1}^k f(c_j) \; \PRmeas_n \big[ (c_{j-1},c_j] \big]
  %      \\
  = \; & \sum_{j=1}^k f(c_j) \, \big( F(c_j) - F(c_{j-1}) \big) .
  \end{align*}
  The difference of the integrals of~$h$ with respect to these two can
  therefore be estimated as
  \begin{align*}
  \Big| \int_{\bR} h \, \ud \PRmeas - \int_{\bR} h \, \ud \PRmeas_n \Big|
  \, = \; & \, \Big| \sum_{j=1}^k f(c_j) \,
      \big( F(c_j) - F_n(c_j) - F(c_{j-1}) +  F_n (c_{j-1}) \big) \Big| \\
  \leq \; & \, \sum_{j=1}^k |f(c_j)| \; \Big(
      \big| F(c_j) - F_n(c_j) \big| + \big| F(c_{j-1}) +  F_n (c_{j-1}) \big|
      \Big) \\
  \leq \; & \, 2 k K \max_{j = 0 , \ldots , k} \big| F(c_j) - F_n(c_j) \big| .
  \end{align*}
  By our assumption~\textup{(ii)}, we have
  $\lim_{n \to \infty} F_n(c_j) = F(c_j)$ for each $j = 1 , \ldots, k$, so
  there exists $N_2$ such that for $n \geq N_2$ we have
  $\max_{j = 1 , \ldots , k} | F(c_j) - F_n(c_j) | < \frac{\eps}{k}$, and
  thus
  \begin{align*}
  \Big| \int_{\bR} h \, \ud \PRmeas - \int_{\bR} h \, \ud \PRmeas_n \Big|
  \leq \; & \, 2 K \eps .
  \end{align*}
  Combining the estimates we have obtained, for $n \geq \max(N_1 , N_2)$, we have
  \begin{align*}
  & \Big| \int_{\bR} f \, \ud \PRmeas - \int_{\bR} f \, \ud \PRmeas_n \Big| \\
  \leq \; \, &
      \underbrace{
      \Big| \int_{\bR} f \, \ud \PRmeas - \int_{\bR} h \, \ud \PRmeas \Big|}_{
          \leq (1+K) \eps}
      + \underbrace{
      \Big| \int_{\bR} h \, \ud \PRmeas - \int_{\bR} h \, \ud \PRmeas_n \Big|}_{
          \leq 2 K \eps}
      + \underbrace{
      \Big| \int_{\bR} h \, \ud \PRmeas_n - \int_{\bR} f \, \ud \PRmeas_n	\Big|}_{
          \leq (1+2K) \eps } \\
  \leq \; \, & (2 + 5 K) \eps .
  \end{align*}
  Since $\eps > 0$ was arbitrary, this shows that
  $\int f \, \ud \PRmeas_n \to \int f \, \ud \PRmeas$ as $n \to \infty$, so we
  have established the weak convergence $\PRmeas_n \to \PRmeas$
  according to Definition~\ref{def:convergence-in-distribution}.
\end{proof}

\begin{lemma}[Necessary condition for convergence in distribution with cdfs]
  \label{lem:cdf-convergence-from-convergence-in-distribution}
  \uses{def:convergence-in-distribution, def:cdf}
  \lean{CumulativeDistributionFunction.tendsto_apply_of_tendsto_of_continuousAt}
  \leanok
  Let $\PRmeas$ and $\PRmeas_n$, $n \in \bN$, be Borel
  probability measures on $\bR$, and let $F$ and $F_n$, $n \in \bN$,
  be their cumulative distribution functions, respectively, i.e.,
  \begin{align*}
      F(x) = \; & \PRmeas \big[(-\infty,x]\big] & & \text{for $x \in \bR$} \\
      F_n(x) = \; & \PRmeas_n \big[(-\infty,x]\big] & & \text{for $x \in \bR$ and $n \in \bN$.}
  \end{align*}
  If $\lim_{n \to \infty} \PRmeas_n = \PRmeas$ in the sense of weak convergence
  of measures, Definition~\ref{def:convergence-in-distribution}, then
  for all continuity points $x$ of $F$ we have
  $\lim_{n \to \infty} F_n(x) = F(x)$.
\end{lemma}
\begin{proof}
  \uses{lem:cdf-continuity-pt-iff-measure-singleton, lem:portmanteau-convergence-for-borel}
  \leanok
  Let $x \in \bR$ be a continuity point of $F$. Then we have $\mu [\set{x}] = 0$,
  by Lemma~\ref{lem:cdf-continuity-pt-iff-measure-singleton}.
  Note that the boundary of the Borel set $(-\infty,x] \subset \bR$ is
  the singleton $\partial (-\infty,x] = \set{x}$.
  Therefore the assumption $\lim_{n \to \infty} \PRmeas_n = \PRmeas$ implies that
  \begin{align*}
    \PRmeas_n \big[ (-\infty,x] \big] \to \PRmeas \big[ (-\infty,x] \big] ,
  \end{align*}
  by a general fact (Lemma~\ref{lem:portmanteau-convergence-for-borel})
  about weakly converging sequences of measures
  that a for Borel sets whose boundary carries no mass in the limit measure.
  In terms of the c.d.f.s, the above reads
  \begin{align*}
    F_n(x) \to F(x)
  \end{align*}
  as asserted.
\end{proof}


\chapter{Transforms of cumulative distribution functions}

In this part, we introduce certain transforms and extensions of
cumulative distribution functions, which are used in the classification
calculation of the extreme value distributions.


\section{Extended cumulative distribution function}

\begin{definition}[Extended cumulative distribution function]
  \label{def:cdf-extend}
  \uses{def:cdf}
  \lean{CumulativeDistributionFunction.extend}
  \leanok
  The extension $\extendcdf{F}$ of a c.d.f. $F$ is the function
  \begin{align*}
    \extendcdf{F} \colon [-\infty,+\infty] \to [0,1]
  \end{align*}
  given by
  \begin{align*}
    \extendcdf{F}(x) = \begin{cases}
      0 & \; \text{ if } x = -\infty \\
      F(x) & \; \text{ if } -\infty < x < +\infty \\
      1 & \; \text{ if } x = +\infty .
      \end{cases}
  \end{align*}
\end{definition}

\begin{lemma}[Continuity points of extended c.d.f.]
  \label{lem:cdf-extend-continuity-pts}
  \uses{def:cdf-extend}
  \lean{CumulativeDistributionFunction.extend_continuousAt,
  CumulativeDistributionFunction.extend_continuousAt_top,
  CumulativeDistributionFunction.extend_continuousAt_bot}
  \leanok
  The extension $\extendcdf{F}$ of a c.d.f. $F$ is continuous
  at $x = -\infty$, $x = + \infty$, and at any
  $x \in (-\infty,+\infty)$ where $F$ is continuous.
\end{lemma}
\begin{proof}
  % \uses{}
  % \leanok
  Since $\lim_{x \to +\infty} F(x) = 1$ by properties of c.d.f.s
  and $\extendcdf{F}(+\infty) = 1$ by definition of the extension,
  continuity at $x = + \infty$ follows. Continuity at $x=-\infty$
  is similar.

  Suppose $F$ is continuous at $x \in \bR$.
  %Let $\iota \colon \bR \to [-\infty,+\infty]$ be the inclusion
  %of the real line into the extended real line.
  Then since $\extendcdf{F}$ coincides with $F$
  in a neighborhood of $x$
  (indeed on the open set $\bR \subsetneq [-\infty,+\infty]$),
  the continuity of $F$ at $x \in \bR$ implies
  continuity of $\extendcdf{F}$ at $x$.
\end{proof}

\section{One over one minus cumulative distribution function}

\begin{definition}[One over one minus cumulative distribution function]
  \label{def:one-div-one-sub-cdf}
  \uses{def:cdf-extend}
  \lean{CumulativeDistributionFunction.oneDivOneSub}
  \leanok
  The transform $\oneDivOneSub{F}$ of a c.d.f. $F$ is the function
  \begin{align*}
    \oneDivOneSub{F} \colon [-\infty,+\infty] \to [1,+\infty]
  \end{align*}
  given by
  \begin{align*}
    \oneDivOneSub{F}(x) = \begin{cases}
      1 & \; \text{ if } \extendcdf{F}(x) = 0 \\
      \frac{1}{1-\extendcdf{F}(x)} & \; \text{ if } 0 < \extendcdf{F}(x) < 1 \\
      +\infty & \; \text{ if } \extendcdf{F}(x) = 1 ,
      \end{cases}
  \end{align*}
  where $\extendcdf{F} \colon [-\infty,+\infty] \to [0,1]$ is the extension of the c.d.f. $F$.
\end{definition}

\begin{lemma}[Continuity points of one over one minus c.d.f.]
  \label{lem:one-div-one-sub-cdf-continuity-pts}
  \uses{def:one-div-one-sub-cdf}
  \lean{CumulativeDistributionFunction.oneDivOneSub_continuousAt,
  CumulativeDistributionFunction.oneDivOneSub_continuousAt_top,
  CumulativeDistributionFunction.oneDivOneSub_continuousAt_bot}
  \leanok
  The transform $\oneDivOneSub{F}$ of a c.d.f. $F$ is continuous
  at $x = -\infty$, $x = + \infty$, and at any
  $x \in (-\infty,+\infty)$ where $F$ is continuous.
\end{lemma}
\begin{proof}
  \uses{lem:cdf-extend-continuity-pts}
  % \leanok
  Since $\lim_{x \to +\infty} \extendcdf{F}(x) = \extendcdf{F}(+\infty) = 1$
  by Lemma~\ref{lem:cdf-extend-continuity-pts} and the continuous extension
  of $p \mapsto \frac{1}{1-p}$ to a function $[0,1] \to [0,+\infty]$
  tends to $+\infty$ as $p \to 1$, we have
  \begin{align*}
    \lim_{x \to +\infty} \oneDivOneSub{F}(x)
    \; = \; +\infty \; = \; \oneDivOneSub{F}(+\infty) .
  \end{align*}
  Therefore $\oneDivOneSub{F}$ is continuous at $+\infty$.
  Continuity at $-\infty$ similarly follows from
  $\lim_{x \to -\infty} \extendcdf{F}(x) = \extendcdf{F}(-\infty) = 0$
  and $\frac{1}{1-p}$ tending to $1$
  as $p \to 0$, which give
  \begin{align*}
    \lim_{x \to -\infty} \oneDivOneSub{F}(x)
    \; = \; 1 \; = \; \oneDivOneSub{F}(-\infty) .
  \end{align*}

  Suppose $F$ is continuous at $x \in \bR$,
  and recall from Lemma~\ref{lem:cdf-extend-continuity-pts} that
  $\extendcdf{F}$ is then also continuous at $x$.
  Now $\oneDivOneSub{F}$ is a composition of the continuous function
  $p \mapsto \frac{1}{1-p} \colon [0,1] \to [0,+\infty]$ with
  $\extendcdf{F}$, and as such also becomes continuous at $x$.
\end{proof}

\section{One over negative logarithm cumulative distribution function}

\begin{definition}[One over negative logarithm cumulative distribution function]
  \label{def:one-div-neg-log-cdf}
  \uses{def:cdf-extend}
  \lean{CumulativeDistributionFunction.oneDivNegLog}
  \leanok
  The transform $\oneDivNegLog{F}$ of a c.d.f. $F$ is the function
  \begin{align*}
    \oneDivNegLog{F} \colon [-\infty,+\infty] \to [0,+\infty]
  \end{align*}
  given by
  \begin{align*}
    \oneDivNegLog{F}(x) = \begin{cases}
      0 & \; \text{ if } \extendcdf{F}(x) = 0 \\
      \frac{1}{\log \big( 1/\extendcdf{F}(x) \big)} & \; \text{ if } 0 < \extendcdf{F}(x) < 1 \\
      +\infty & \; \text{ if } \extendcdf{F}(x) = 1 ,
      \end{cases}
  \end{align*}
  where $\extendcdf{F} \colon [-\infty,+\infty] \to [0,1]$ is the extension of the c.d.f. $F$.
\end{definition}

\begin{lemma}[Continuity points of one over negative logarithm c.d.f.]
  \label{lem:one-div-neg-log-cdf-continuity-pts}
  \uses{def:one-div-neg-log-cdf}
  \lean{CumulativeDistributionFunction.oneDivNegLog_continuousAt,
  CumulativeDistributionFunction.oneDivNegLog_continuousAt_top,
  CumulativeDistributionFunction.oneDivNegLog_continuousAt_bot}
  \leanok
  The transform $\oneDivNegLog{F}$ of a c.d.f. $F$ is continuous
  at $x = -\infty$, $x = + \infty$, and at any
  $x \in (-\infty,+\infty)$ where $F$ is continuous.
\end{lemma}
\begin{proof}
  \uses{lem:cdf-extend-continuity-pts}
  % \leanok
  \uses{lem:cdf-extend-continuity-pts}
  % \leanok
  Since $\lim_{x \to +\infty} \extendcdf{F}(x) = \extendcdf{F}(+\infty) = 1$
  by Lemma~\ref{lem:cdf-extend-continuity-pts} and the continuous extension
  of $p \mapsto \frac{1}{\log(1/p)}$ to a function $[0,1] \to [0,+\infty]$
  tends to $+\infty$ as $p \to 1$, we have
  \begin{align*}
    \lim_{x \to +\infty} \oneDivNegLog{F}(x)
    \; = \; +\infty \; = \; \oneDivNegLog{F}(+\infty) .
  \end{align*}
  Therefore $\oneDivNegLog{F}$ is continuous at $+\infty$.
  Continuity at $-\infty$ similarly follows from
  $\lim_{x \to -\infty} \extendcdf{F}(x) = \extendcdf{F}(-\infty) = 0$
  and $\frac{1}{\log(1/p)}$ tending to $0$
  as $p \to 0$, which give
  \begin{align*}
    \lim_{x \to -\infty} \oneDivNegLog{F}(x)
    \; = \; 0 \; = \; \oneDivNegLog{F}(-\infty) .
  \end{align*}

  Suppose $F$ is continuous at $x \in \bR$,
  and recall from Lemma~\ref{lem:cdf-extend-continuity-pts} that
  $\extendcdf{F}$ is then also continuous at $x$.
  Now $\oneDivNegLog{F}$ is a composition of the continuous function
  $p \mapsto \frac{1}{\log(1/p)} \colon [0,1] \to [0,+\infty]$ with
  $\extendcdf{F}$, and as such also becomes continuous at $x$.
\end{proof}


\chapter{One-parameter subgroups of affine isomorphisms}

%The group $\OriAffR$ is isomorphic to the matrix group
%\begin{align*}
%  \left\{ \left[ \begin{array}{cc} a & b \\ 0 & 1 \end{array} \right] \; \Bigg|
%    \; a>0, b \in \bR \right\} .
%\end{align*}
%This is a 2-dimensional Lie group with Lie algebra\ldots


\begin{lemma}[Functional equation in one parameter subgroups of affine isomorphisms]
  \label{lem:affine-one-parameter-subgroup-functional-eqn}
  \uses{def:oriented-affine-isomorphism}
  % \lean{}
  % \leanok
  Suppose that $t \mapsto A_t$ is a homomorphism of multiplicative groups
  $(0,+\infty) \to \OriAffR$, i.e., for any $s, t > 0$ we have
  \begin{align*}
    A_{s t} = A_s \circ A_t .
  \end{align*}
  Write $A_t(x) = a_t x + b_t$, with $a_t > 0$ and $b_t \in \bR$.
  Then we have, for any $s, t > 0$,
  \begin{align*}
    a_{t s} = \; & a_t \, a_s
      \qquad \text{ and } \\
    b_{t s} = \; & a_t \, b_s + b_t .
  \end{align*}
  (Also by symmetry $b_{t s} = a_s \, b_t + b_s$.)
\end{lemma}
\begin{proof}
  % \uses{}
  % \leanok
\end{proof}

\begin{lemma}[Functional equation scaling coefficient solution]
  \label{lem:solution-functional-eqn-scaling}
  %\uses{def:lem:oriented-affine-action-on-cdf}
  \lean{eq_exp_const_mul_of_multiplicative_of_measurable}
  \leanok
  Suppose that $a \colon (0,+\infty) \to (0,+\infty)$ is a
  measurable function satisfying, for any $s, t > 0$,
  \begin{align*}
    a(t s) = \; & a(t) \, a(s) .
  \end{align*}
  Then there exists a $\rho \in \bR$ such that for all $t > 0$,
  \begin{align*}
    a(t) = t^{\rho} .
  \end{align*}
\end{lemma}
\begin{proof}
  \uses{lem:different-difference-set-contains-interval}
  % \leanok
\end{proof}

\begin{lemma}[Functional equation translation coefficient solution with $\rho = 0$]
  \label{lem:solution-functional-eqn-translation-no-scaling}
  %\uses{def:lem:oriented-affine-action-on-cdf}
  \lean{eq_const_mul_of_additive_of_measurable,
        eq_exp_const_mul_of_multiplicative_of_measurable}
  \leanok
  Suppose that
  $b \colon (0,+\infty) \to \bR$ is a measurable function
  satisfying, for any $s, t > 0$,
  \begin{align*}
    b(t s) = \; & b(s) + b(t) .
  \end{align*}
  Then there exists a constant $c$ such that for $t \in (0,+\infty)$
  we have
  \begin{align*}
    b(t) = \; & -c \, \log (t) .
  \end{align*}
\end{lemma}
\begin{proof}
  \uses{lem:solution-functional-eqn-scaling}
  \leanok
\end{proof}

\begin{lemma}[Functional equation translation coefficient solution with $\rho \ne 0$]
  \label{lem:solution-functional-eqn-translation-with-scaling}
  %\uses{def:lem:oriented-affine-action-on-cdf}
  \lean{eq_const_mul_one_sub_exp_of_multiplicative_with_scaling_of_measurable}
  \leanok
  Suppose that $\rho \in \bR \setminus \set{0}$ and
  $b \colon (0,+\infty) \to \bR$ is a measurable function
  satisfying, for any $s, t > 0$,
  \begin{align*}
    b(t s) = \; & t^{\rho} \, b(s) + b(t) .
  \end{align*}
  Then there exists a constant $c$ such that for $t \in (0,+\infty) \setminus \set{1}$
  we have
  \begin{align*}
    b(t) = \; & c (1 - t^{-\rho}) .
  \end{align*}
\end{lemma}
\begin{proof}
  \uses{lem:solution-functional-eqn-scaling}
  \leanok
\end{proof}

The above solutions to functional equations can be used to classify all
one-parameter subgroups of the group of oriented affine isomorphisms of $\bR$.
Such a subgroup can be given in a parametrized form as a group homomorphism
$\bR \to \OriAffR$ (from the additive group $\bR$) or alternatively
as group homomorphisms $(0,+\infty) \to \OriAffR$ (from the multiplicative
group $(0,+\infty)$). The additive and multiplicative versions are
related by the change of variable
$\bR \ni t \leftrightarrow \lambda := e^t \in (0,+\infty)$
(conversely, $t = \log(\lambda)$).
\emph{(In Lean the type $\bR$ is more convenient than the type $(0,+\infty)$,
so in formal statements we prefer the former choice.)}

\begin{theorem}[One-parameter subgroups of affine isomorphisms of $\bR$]
  \label{thm:one-parameter-subgroups-of-affine-isomorphisms}
  \uses{def:oriented-affine-isomorphism, def:translation-subgroup, def:fixing-subgroup}
  \lean{AffineIncrEquiv.homomorphism_from_Real_characterization,
    AffineIncrEquiv.homomorphism_from_Real_characterization_of_nontrivial}
  \leanok
  \emph{[TODO: Switch to additive notation and $\bR$ rather than multiplicative
  notation and $(0,+\infty)$, to match the most convenient formalized statements.]}

  Suppose that $t \mapsto A_t$ is a measurable homomorphism of multiplicative groups
  $(0,+\infty) \to \OriAffR$, i.e., for any $s, t > 0$ we have
  \begin{align*}
    A_{s t} = A_s \circ A_t
  \end{align*}
  and $A_t(x) = a_t x + b_t$, with $t \mapsto a_t$ and $t \mapsto b_t$
  measurable functions.
  Then either
  \begin{description}
    \item[(0)] there exists a $\beta \in \bR$ such that for all $t > 0$ and $x \in \bR$
      \begin{align*}
        A_t(x) = x + \beta \, \log(t) ;
      \end{align*}
      or
    \item[(1)] there exists a $\rho \ne 0$ and $c \in \bR$ such that
      for all $t > 0$ and $x \in \bR$
      \begin{align*}
        A_t(x) = t^{\rho} (x - c) + c .
      \end{align*}
  \end{description}
\end{theorem}
\begin{proof}
  \uses{lem:affine-one-parameter-subgroup-functional-eqn,
    lem:solution-functional-eqn-scaling,
    lem:solution-functional-eqn-translation-no-scaling,
    lem:solution-functional-eqn-translation-with-scaling}
  \leanok
\end{proof}


\chapter{Types}

\begin{definition}[Type (of distribution on $\bR$)]
  \label{def:cdf-type}
  \uses{def:cdf, lem:oriented-affine-action-on-cdf}
  %\lean{lcinv}
  %\leanok
  Two c.d.f.s $F, G$ are said to be of the same type, if
  there exists an order-preserving affine isomorphism $A \in \OriAffR$
  such that $G = A . F$.

  Being of the same type is an equivalence relation, and
  the equivalence classes are called types (of distributions on $\bR$).
\end{definition}


\section{Convergence to types}

The notion of convergence of cumulative distribution functions
considered here is always taken to be pointwise convergence
on the set of continuity points of the limit c.d.f.
By Theorem~\ref{thm:convergence-in-distribution-with-cdf} and
Lemma~\ref{lem:cdf-convergence-from-convergence-in-distribution},
this corresponds to convergence in distribution (weak convergence
of probability measures). Below, when we write $F_n \tendsinlaw F$
for c.d.f.s $F_n$, $n \in \bN$, and $F$, this is always what we
mean.

For affine maps $A \colon \bR \to \bR$, we use the topology
of pointwise convergence of functions. Equivalently, convergence of affine maps
means the convergence of the coefficients $a, b \in \bR$ in the expression
$x \mapsto a x + b$ of the affine functions
(so $A_n \to A$ if and only if the functions are of the form
$A_n(x) = a_n x + b_n$ and $A(x) = a x + b$, and $a_n \to a$ and $b_n \to b$).

\begin{lemma}[Unique affine relation among two nondegenerate c.d.f.s]
  \label{lem:unique-affine-relation-cdf}
  \uses{def:cdf, def:degenerate-cdf, lem:oriented-affine-action-on-cdf}
  \lean{CumulativeDistributionFunction.unique_orientationPreservingAffineEquiv_smul_eq_not_isDegenerate}
  \leanok
  Let $F, G$ be two c.d.f.s of the same type, and $A \in \OriAffR$
  an affine isomorphism such that $G = A.F$. If $F$ is nondegenerate,
  then $A$ is the only element of $\OriAffR$ for which the
  relation $G = A.F$ holds.
\end{lemma}
\begin{proof}
  % \uses{}
  % \leanok
  Since $F$ is nondegenerate, we can find two different points $x_1 < x_2$
  such that $0 < F(x_1) < F(x_2) \le 1$. By right continuity of $F$, we can
  assume these points to be taken minimal with the given values, i.e.,
  $x_j = \inf \big\{ x \in \bR \; \big| \; F(x) = F(x_j) \big\}$ for
  $j=1,2$.

  The assumption $G = A.F$ means $G(x) = F \big( A^{-1}(x) \big)$ for all $x \in \bR$.
  Therefore $G \big( A(x_1) \big) = F(x_1) < F(x_2) = G \big( A(x_2) \big)$.
  We also get $A(x_j) = \inf \big\{ y \in \bR \; \big| \; G(y) = F(x_j) \big\}$ for
  $j=1,2$ by strict monotonicity and bijectivity of $A$
  (if, for example, there would exist a $y'_2 < A(x_2)$ such that $G(y'_2) = F(x_2)$,
  then the point $x'_2 = A^{-1}(y'_2) < x_2$ would be such that
  $F\big(A^{-1}(y'_2)\big) = G(y'_2) = F(x_2)$, contradicting the minimality of $x_2$).

  If $\widetilde{A} \in \OriAffR$ is also such that $G = \widetilde{A}.F$, then
  the same holds for it:
  $\widetilde{A}(x_j) = \inf \big\{ y \in \bR \; \big| \; G(y) = F(x_j) \big\}$
  for $j=1,2$. We conclude that
  \begin{align*}
    \widetilde{A}(x_1) \, = \; &
        \inf \big\{ y \in \bR \; \big| \; G(y) = F(x_1) \big\} \; = \, A(x_1) \\
    \widetilde{A}(x_2) \, = \; &
        \inf \big\{ y \in \bR \; \big| \; G(y) = F(x_2) \big\} \; = \, A(x_2) .
  \end{align*}
  But an affine map of $\bR$ is determined by its values at two distinct points:
  from $a x_1 + b = y_1$ and $a x_2 + b = y_2$ with $x_1 \ne x_2$ one can solve $a, b$.
  Therefore we must have $\widetilde{A} = A$.
\end{proof}

\begin{lemma}[Degeneration by shrinking affine transformations]
  \label{lem:degenerate-shrinking-limit}
  \uses{def:cdf-degenerate, lem:oriented-affine-action-on-cdf}
  \lean{CumulativeDistributionFunction.isDegenerate_of_tendsto_shrinking}
  \leanok
  Let $(F_n)_{n \in \bN}$ be a sequence of c.d.f.s which converges to a c.d.f. $G$,
  $F_n \tendsinlaw G$. Consider affine transformations of the form $A_n(x) = a_n x + b_n$,
  with $a_n > 0$ and $b_n \in \bR$, such that $a_n \to 0$ and $b_n \to \beta \in \bR$
  as $n \to \infty$. Then $A_n . F_n \tendsinlaw \widetilde{G}$, where
  $\widetilde{G}$ is the degenerate c.d.f. of the delta mass at $\beta$.
\end{lemma}
\begin{proof}
  % \uses{}
  % \leanok
  It suffices to prove that for any $x < \beta$ we have $\widetilde{G}(x) = 0$
  and for any $x > \beta$ we have $\widetilde{G}(x) = 1$.

  Let us focus on the latter, so let $x > \beta$. Let also $\eps > 0$; we will prove
  that $\widetilde{G}(x) > 1 - \eps$, and the claim will follow.

  By density of continuity points of $\widetilde{G}$, we can choose $x'$ such that
  $\beta < x' < x$ and $\widetilde{G}$ is continuous at $x'$. Then the assumed
  convergence $A_n . F_n \tendsinlaw \widetilde{G}$ implies that
  $(A_n . F_n)(x') \to \widetilde{G}(x')$. Since $\widetilde{G}(x) \ge \widetilde{G}(x')$,
  it suffices to prove that $\widetilde{G}(x') > 1 - \eps$.

  Since $G$ is a c.d.f., we can choose a continuity point $z$ of $G$
  large enough so that $G(z) > 1 - \eps$.
  Then by the assumed convergence $F_n \tendsinlaw G$, we have
  $F_n(z) \to G(z)$. By definition we have
  \begin{align*}
    (A_n . F_n)(A_n(z)) \; = \; F_n \big( A_n^{-1} (A_n(z)) \big) \; = \; F_n(z)
      \; \longrightarrow \; G(z) . %\; > \;  1-\eps .
  \end{align*}
  Note that $A_n(z) = a_n z + b_n \to \beta$ as $n \to \infty$ by the assumptions
  $a_n \to 0$ and $b_n \to \beta$. In particular, for $n$ large enough,
  we have $A_n(z) < x'$.
  Therefore, for $n$ large enough
  \begin{align*}
    F_n(z) \; = \; (A_n . F_n) \big( A_n(z) \big) \; \le \; (A_n . F_n)(x') .
  \end{align*}
  The LHS tends to $G(z)$ and the RHS tends to $\widetilde{G}(x')$,
  showing
  \begin{align*}
    1 - \eps \; < \; G(z) \; \le \; \widetilde{G}(x') \; \le \; \widetilde{G}(x) .
  \end{align*}
  This concludes the proof that $\widetilde{G}(x) = 1$ for all $x > \beta$.

  The proof that $\widetilde{G}(x) = 0$ for all $x < \beta$ is similar.
\end{proof}

\begin{lemma}[Impossibility of expanding affine transformations]
  \label{lem:impossible-expanding-limit}
  \uses{def:cdf, lem:oriented-affine-action-on-cdf}
  \lean{CumulativeDistributionFunction.not_tendsto_cdf_of_expanding_of_tendsto_not_isDegenerate}
  \leanok
  Let $(F_n)_{n \in \bN}$ be a sequence of c.d.f.s which converges to a nondegenerate
  c.d.f. $G$, $F_n \tendsinlaw G$. Consider affine transformations of the form $A_n(x) = a_n x + b_n$,
  with $a_n > 0$ and $b_n \in \bR$, such that $a_n \to +\infty$
  as $n \to \infty$. Then $A_n . F_n$ cannot converge to any c.d.f.
\end{lemma}
\begin{proof}
  \uses{lem:exists-two-nontrivial-continuity-pts-cdf}
  \leanok
  Since $G$ is assumed nondegenerate, by Lemma~\ref{lem:exists-two-nontrivial-continuity-pts-cdf}
  we can pick two continuity points $x_1 < x_2$
  of $G$ such that $0 < G(x_1) \le G(x_2) < 1$. Then from the assumption
  $F_n \tendsinlaw G$ we get $F_n(x_1) \to G(x_1)$ and $F_n(x_2) \to G(x_2)$.

  Assume, by way of contradiction, that we have convergence
  $A_n . F \tendsinlaw \widetilde{G}$ to some c.d.f. $\widetilde{G}$.

  We claim that then $\big(A_n(x_1)\big)_{n \in \bN}$ is bounded from below and
  $\big(A_n(x_2)\big)_{n \in \bN}$ is bounded from above.
  Since
  \begin{align*}
    a_n = \frac{A_n(x_2) - A_n(x_1)}{x_2 - x_1} ,
  \end{align*}
  this will show that $(a_n)_{n \in \bN}$ is bounded from above, contradicting
  the assumption $a_n \to +\infty$, and finishing the proof.

  To show that $\big(A_n(x_1)\big)_{n \in \bN}$ is bounded from below,
  choose a continuity point $z$ of $\widetilde{G}$ such that
  $\widetilde{G}(z) < G(x_1)$. Then the assumed convergence
  $A_n . F \tendsinlaw \widetilde{G}$ implies
  $(A_n . F)(z) \to \widetilde{G}(z)$.
  On the other hand, if $\big(A_n(x_1)\big)_{n \in \bN}$ is not
  bounded from below, then along some subsequence $(n_k)_{k \in \bN}$ of indices
  we have $A_{n_k}(x_1) < z$, and for those indices we then have
  \begin{align*}
    F_{n_k}(x_1) \; = \; (A_{n_k} . F_{n_k}) \big( A_{n_k}(x_1) \big)
      \; \le \; (A_{n_k} . F_{n_k})(z) .
  \end{align*}
  The LHS tends to $G(x_1)$ as $k \to \infty$, whereas the RHS tends
  to $\widetilde{G}(z)$. We get $G(x_1) \le \widetilde{G}(z)$, contradicting
  the choice of $z$. This shows that $\big(A_n(x_1)\big)_{n \in \bN}$ must
  in fact be bounded from below.

  The proof that $\big(A_n(x_2)\big)_{n \in \bN}$ must be bounded from above
  is similar.
\end{proof}

\begin{theorem}[Convergence to types]
  \label{thm:convergence-to-types}
  \uses{def:degenerate-cdf, lem:oriented-affine-action-on-cdf, def:cdf-type}
  % \lean{}
  % \leanok
  Suppose that $(F_n)_{n \in \bN}$ is a sequence of
  c.d.f.s which converges to a nondegenerate c.d.f. $G$, i.e.,
  $F_n \tendsinlaw G$ as $n \to \infty$.
  Let $(A_n)_{n \in \bN}$ be a sequence of oriented
  affine isomorphisms of $\bR$, $A_n \in \OriAffR$
  such that $A_n.F_n \tendsinlaw \widetilde{G}$ for some c.d.f.
  $\widetilde{G}$.

  If we write $A_n(x) = a_n x + b_n$, then $(a_n)_{n \in \bN}$
  and $(b_n)_{n \in \bN}$ are bounded sequences.

  If $\widetilde{G}$ is nondegenerate, then
  $A_n \to A \in \OriAffR$ and $A.G = \widetilde{G}$.
  In particular $G$ and $\widetilde{G}$ are of the same type.
  Moreover, $A$ is the unique affine transformation for which
  the equality $A.G = \widetilde{G}$ holds.
\end{theorem}
\begin{proof}
  \uses{lem:impossible-expanding-limit, lem:degenerate-shrinking-limit,
    lem:unique-affine-relation-cdf, lem:action-on-cdf-continuous}
  % \leanok
  Let us first argue that $(a_n)_{n \in \bN}$ are bounded.
  If not, then by passing to a subsequence, we have $a_{n_k} \to +\infty$.
  But since $F_{n_k} \tendsinlaw G$ and $G$ is nondegenerate,
  it contradicts Lemma~\ref{lem:impossible-expanding-limit}
  to have $A_{n_k} . F_{n_k} \tendsinlaw \widetilde{G}$.
  Therefore $(a_n)_{n \in \bN}$ must be bounded: there exists
  some $M>0$ such that $a_n \le M$ for all $n \in \bN$.

  Let us then argue that $(b_n)_{n \in \bN}$ is bounded.
  If not, then we can extract a subsequence such
  that either $b_{n_k} \to -\infty$ or $b_{n_k} \to +\infty$.
  Let us prove the impossibility of the second one, the first is similar.
  So assume that $b_{n_k} \to +\infty$.
  Since $G$ is nondegenerate, we may pick a continuity point $x_0$
  of $G$ such that $0 < G(x_0) < 1$. Then we have $F_n(x_0) \to G(x_0)$
  by the assumption $F_n \tendsinlaw G$.
  We may also pick a continuity point $z$ of $\widetilde{G}$
  such that $\widetilde{G}(z) > G(x_0)$. Then we have
  $(A_n.F_n)(z) \to \widetilde{G}(z)$
  by the assumption $A_n.F_n \tendsinlaw \widetilde{G}$.
  But $A_{n_k}(x_0) = a_{n_k} x_0 + b_{n_k} \to +\infty$, since $0 < a_{n_k} \le M$
  and $b_{n_k} \to +\infty$. Therefore we have for all large enough $k$
  that $A_{n_k}(x_0) > z$. And then
  \begin{align*}
    (A_{n_k}.F_{n_k})(z) \; \le \; (A_{n_k}.F_{n_k})\big(A_{n_k}(x_0)\big) \; = \; F_{n_k}(x_0) .
  \end{align*}
  The LHS tends to $\widetilde{G}(z)$ as $k \to \infty$, and the RHS tends to $G(x_0)$.
  Therefore we get $\widetilde{G}(z) \le G(x_0)$, contradicting the choice of $z$.
  This shows that we cannot have $b_{n_k} \to +\infty$. Similarly one proves that
  we cannot have $b_{n_k} \to -\infty$. We conclude that $(b_n)_{n \in \bN}$
  is indeed bounded.

  From now on, suppose furthermore that also $\widetilde{G}$ is nondegenerate.
  We claim that then $(a_n)_{n \in \bN}$ is bounded away from $0$:
  for some $\eps > 0$ we have $a_n \ge \eps$ for all $n \in \bN$.
  If not, then we could extract a subsequence such that $a_{n_k} \to 0$
  and also $b_n \to \beta$ (since $b_n$ is bounded).
  But since $F_{n_k} \tendsinlaw G$ and $G$ is nondegenerate,
  from Lemma~\ref{lem:degenerate-shrinking-limit}
  we would get $A_{n_k} . F_{n_k} \tendsinlaw \widetilde{G}_0$,
  where $\widetilde{G}_0$ is degenerate. But by assumption
  $A_{n_k} . F_{n_k} \tendsinlaw \widetilde{G}$, where
  $\widetilde{G}$ is nondegenerate; this is impossible by uniqueness
  of limits for convergence in distribution. Therefore
  $(a_n)_{n \in \bN}$ must indeed be bounded away from $0$.

  Note that since $(a_n)_{n \in \bN}$ is bounded away from $0$ and $+\infty$,
  and $(b_n)_{n \in \bN}$ is bounded, we can extract a subsequence such that
  $a_{n_k} \to \alpha \in (0,+\infty)$ and $b_{n_k} \to \beta \in \bR$.
  By assumption, we have $A_{n_k} . F_{n_k} \tendsinlaw \widetilde{G}$.
  But since $A_{n_k} \to A$ where $A(x) = \alpha x + \beta$
  and we have also assumed $F_{n_k} \tendsinlaw G$, this implies by
  continuity (Lemma~\ref{lem:action-on-cdf-continuous}) that $A_{n_k} . F_{n_k} \tendsinlaw A . G$.
  By uniqueness of limits, we get $A . G = \widetilde{G}$.

  To prove that in fact $A_n \to A$, not just along a subsequence, note the
  following. From any subsequence $A_{n_k}$, we can extract a further
  convergent subsequence of values of $a_{n_{k_\ell}}$ and $b_{n_{k_\ell}}$
  values as above, with limits $\alpha' \in (0,+\infty)$
  and $\beta' \in \bR$. The same argument as above shows that
  $A' . G = \widetilde{G}$ where $A'(x) = \alpha' x + \beta'$.
  Lemma~\ref{lem:unique-affine-relation-cdf} then says that we must have
  $A' = A$, i.e., $\alpha' = \alpha$ and $\beta' = \beta$. Since any
  subsequence has a convergent further subsequence with the same limit,
  the entire sequence must converge, $A_n \to A$.
  The proof is complete.
\end{proof}

\begin{theorem}[Convergence to types again]
  \label{thm:convergence-to-types-different-affine}
  \uses{def:cdf, lem:oriented-affine-action-on-cdf, def:cdf-type}
  % \lean{} % also `AffineIncrEquiv.inv_mul_eq_mkOfCoefs`
  % \leanok
  Let $(A_n)_{n \in \bN}$ and
  $(\widetilde{A}_n)_{n \in \bN}$ be two sequences of oriented
  affine isomorphisms of $\bR$, $A_n, \widetilde{A}_n \in \OriAffR$.
  Write $A_n(x) = a_n x + b_n$
  and $\widetilde{A}_n(x) = \tilde{a}_n x + \tilde{b}_n$,
  and for the inverses $A_n^{-1}(x) = c_n x + d_n$
  and $\widetilde{A}^{-1}_n(x) = \tilde{c}_n x + \tilde{d}_n$.

  Let $(F_n)_{n \in \bN}$ be a sequence of c.d.f.s
  such that $A_n.F_n \tendsinlaw G$, with
  $G$ a nondegenerate c.d.f.

  Then the convergence of also $\widetilde{A}_n.F_n \tendsinlaw G$
  holds if and only if the coefficients of the affine maps
  satisfy the relations
  \begin{align*}
    \frac{\tilde{a}_n}{a_n} \to 1
    \quad \text{ and } \quad
    \frac{a_n \tilde{b}_n - \tilde{a}_n b_n}{a_n} \to 0,
  \end{align*}
  or equivalently,
  \begin{align*}
    \frac{\tilde{c}_n}{c_n} \to 1
    \quad \text{ and } \quad
    \frac{\tilde{d}_n - d_n}{c_n} \to 0 .
  \end{align*}
\end{theorem}
\begin{proof}
  \uses{thm:convergence-to-types}
  % \leanok
  \emph{(This is actually just a special case of what is stated as
  Corollary~\ref{cor:convergence-to-types-different-limit} below.
  The better organization is to prove that corollary
  directly, and obtain this lemma as a special case.)}

  We will apply the convergence to types with the reference sequence
  $(A_n.F_n)_{n \in \bN}$, which by assumption tends to a
  nondegenerate $G$.

  To express the other sequence in terms of the reference sequence, we write
  \begin{align*}
    \widetilde{A}_n.F_n \; = \; (\widetilde{A}_n A_n^{-1}).(A_n.F_n) .
  \end{align*}
  By assumption this also tends to $G$.

  Theorem~\ref{thm:convergence-to-types} applies, and guarantees convergence
  $\widetilde{A}_n A_n^{-1} \to A$ to some $A \in \OriAffR$,
  and it also implies $A.G = G$ (note that the other limit is also $G$
  by our assumptions).
  However, the unique $A$ for which we have $A.G = G$ is $A = \mathrm{id}$.
  We therefore get $\widetilde{A}_n A_n^{-1} \to \id$.
  To explicitly see the coefficients, note that
  \begin{align*}
    A_n^{-1}(x) \; = \; & a_n^{-1} (x - b_n)
  \end{align*}
  and
  \begin{align*}
    \widetilde{A}_n \big( A_n^{-1}(x) \big)
    \; = \; & \tilde{a}_n \big(a_n^{-1} (x - b_n)\big) + \tilde{b}_n \\
    \; = \; & \tilde{a}_n a_n^{-1} x - \tilde{a}_n a_n^{-1} b_n + \tilde{b}_n .
  \end{align*}
  The convergence $\widetilde{A}_n A_n^{-1} \to \id$ is equivalent to
  the convergence of the coefficients,
  \begin{align*}
    \tilde{a}_n a_n^{-1} \; \longrightarrow \; & 1 \\
    - \tilde{a}_n a_n^{-1} b_n + \tilde{b}_n \; \longrightarrow \; & 0 .
  \end{align*}
  The second one can be rewritten as
  \begin{align*}
    \frac{a_n \tilde{b}_n - \tilde{a}_n b_n}{a_n} \; \longrightarrow \; & 0 .
  \end{align*}
\end{proof}

\begin{corollary}[Convergence to types with different limits]
  \label{cor:convergence-to-types-different-limit}
  \uses{def:cdf, lem:oriented-affine-action-on-cdf, def:cdf-type}
  % \lean{} % also `AffineIncrEquiv.inv_mul_eq_mkOfCoefs`
  % \leanok
  Let $(A_n)_{n \in \bN}$ and
  $(\widetilde{A}_n)_{n \in \bN}$ be two sequences of oriented
  affine isomorphisms of $\bR$, $A_n, \widetilde{A}_n \in \OriAffR$.
  Write $A_n(x) = a_n x + b_n$
  and $\widetilde{A}_n(x) = \tilde{a}_n x + \tilde{b}_n$,
  and for the inverses $A_n^{-1}(x) = c_n x + d_n$
  and $\widetilde{A}^{-1}_n(x) = \tilde{c}_n x + \tilde{d}_n$.

  Let $(F_n)_{n \in \bN}$ be a sequence of c.d.f.s
  such that $A_n.F_n \tendsinlaw G$ and
  $\widetilde{A}_n.F_n \tendsinlaw \widetilde{G}$, with
  $G$ and $\widetilde{G}$ nondegenerate c.d.f.s.
  Then for some $\alpha>0$ and $\beta \in \bR$ we have
  \begin{align*}
    \frac{\tilde{a}_n}{a_n} \to \alpha
    \quad \text{ and } \quad
    \frac{a_n \tilde{b}_n - \tilde{a}_n b_n}{a_n} \to \beta,
  \end{align*}
  and we have
  \begin{align*}
    A.G = \widetilde{G}
    \qquad \text{ where } \; A(x) = \alpha x + \beta .
  \end{align*}
  Equivalently, with $\gamma = \alpha^{-1}$ and $\delta = -\alpha^{-1} \beta$
  so that $A^{-1}(x) = \gamma x + \delta$, we have
  \begin{align*}
    \frac{\tilde{c}_n}{c_n} \to \gamma
    \quad \text{ and } \quad
    \frac{\tilde{d}_n - d_n}{c_n} \to \delta .
  \end{align*}

  In particular, $G$ and $\widetilde{G}$ have the same type.
\end{corollary}
\begin{proof}
  \uses{thm:convergence-to-types-different-affine}
  % \leanok
  We will apply the convergence to types with the reference sequence
  $(A_n.F_n)_{n \in \bN}$, which by assumption tends to a
  nondegenerate $G$.

  To express the other sequence in terms of the reference sequence, we write
  \begin{align*}
    \widetilde{A}_n.F_n \; = \; (\widetilde{A}_n A_n^{-1}).(A_n.F_n) .
  \end{align*}
  By assumption this tends to a nondegenerate $\widetilde{G}$.

  Theorem~\ref{thm:convergence-to-types} applies, and guarantees convergence
  $\widetilde{A}_n A_n^{-1} \to A$ to some $A \in \OriAffR$,
  and it also implies $A.G = \widetilde{G}$.

  Write $A(x) = \alpha x + \beta$. To explicitly see the coefficients of
  $\widetilde{A}_n A_n^{-1}$, note that
  \begin{align*}
    A_n^{-1}(x) \; = \; & a_n^{-1} (x - b_n)
  \end{align*}
  and
  \begin{align*}
    \widetilde{A}_n \big( A_n^{-1}(x) \big)
    \; = \; & \tilde{a}_n \big(a_n^{-1} (x - b_n)\big) + \tilde{b}_n \\
    \; = \; & \tilde{a}_n a_n^{-1} x - \tilde{a}_n a_n^{-1} b_n + \tilde{b}_n .
  \end{align*}
  The convergence $\widetilde{A}_n A_n^{-1} \to A$ is equivalent to
  the convergence of the coefficients,
  \begin{align*}
    \tilde{a}_n a_n^{-1} \; \longrightarrow \; & \alpha \\
    - \tilde{a}_n a_n^{-1} b_n + \tilde{b}_n \; \longrightarrow \; & \beta .
  \end{align*}
  The second one can be rewritten as
  \begin{align*}
    \frac{a_n \tilde{b}_n - \tilde{a}_n b_n}{a_n} \; \longrightarrow \; & \beta .
  \end{align*}
\end{proof}

\begin{lemma}[A choice of normalizing constants for convergence to types]
  \label{lemma:convergence-to-types-normalization}
  \uses{def:cdf, lem:oriented-affine-action-on-cdf, def:cdf-type, def:lc-inverse}
  % \lean{}
  % \leanok
  (It is possible to choose normalization constants for the affine
  transformations using the left-continuous inverses of the c.d.f.s.
  TODO: Precise statement.)
%  Suppose that $(F_n)_{n \in \bN}$ is a sequence of c.d.f.s
%  and $(A_n)_{n \in \bN}$ %and $(\widetilde{A}_n)_{n \in \bN}$
%  be a sequence of oriented
%  affine isomorphisms of $\bR$, $A_n \in \OriAffR$,
%  such that $A_n.F_n \tendsinlaw G$ where
%  $G$ is a nondegenerate c.d.f.
%
%  (TODO: The following is not quite correct. Need to assume one chooses
%  continuity points of the left-continuous inverse.)
%
%  Then for any $0 < p_1 < p_2 < 1$ and $0 < p_0 < 1$, defining
%  \begin{align*}
%    \tilde{a}_n = \frac{1}{\lcinv{F}_n(p_2) - \lcinv{F}_n(p_1)}
%    \quad \text{ and } \quad
%    \tilde{b}_n = \frac{\lcinv{F}_n(p_0)}{\lcinv{F}_n(p_2) - \lcinv{F}_n(p_1)}
%  \end{align*}
%  we have that
%  \begin{align*}
%    \widetilde{A}_n.F_n \tendsinlaw \widetilde{G}
%    \qquad \text{ where } \; \widetilde{A}_n(x) = \tilde{a}_n x + \tilde{b}_n .
%  \end{align*}
\end{lemma}
\begin{proof}
  \uses{thm:convergence-to-types-different-affine}
  % \leanok
  \ldots
\end{proof}


\section{One-parameter subgroups of affine transformations}

\begin{definition}[Subgroup of translations]
  \label{def:translation-subgroup}
  \uses{def:oriented-affine-isomorphism}
  \lean{AffineIncrEquiv.homOfIndex₀, AffineIncrEquiv.subGroupOfIndex₀}
  % \leanok
  The mapping $s \mapsto A_s$ with
  \begin{align*}
    A_s(x) = x + s
  \end{align*}
  is a homomorphism $\bR \to \OriAffR$. The image of this
  homomorphism is the \emph{subgroup of translations} in $\OriAffR$.
\end{definition}

\begin{lemma}[Only translations have no fixed points]
  \label{lem:no-fixed-point-implies-translation}
  \uses{def:translation-subgroup}
  \lean{AffineIncrEquiv.mem_subGroupOfIndex₀_of_no_fixed_point}
  \leanok
  If $A \in \OriAffR$ has no fixed points (no $x \in \bR$ such that $A(x) = x$)
  then $A$ belongs to the subgroup of translations, i.e., $A(x) = x + s$
  for some $s \in \bR$ (in fact $s \ne 0$).
\end{lemma}
\begin{proof}
  % \uses{}
  % \leanok
  Let us prove this by contrapositive: that any element $A$ which is not a
  translation must have a fixed point.
  So assume that $A$ is not a translation, i.e.,
  $A(x) = a x + b$ with some $a \ne 1$. Then the fixed point equation
  $A(x) = x$ reads
  \begin{align*}
    a x + b = x  ,
  \end{align*}
  and it has a solution $x = \frac{-b}{a-1} \in \bR$, which then is a
  fixed point of $A$.
\end{proof}

\begin{lemma}[Conjugate of translation is translation]
  \label{lem:conjugate-translation}
  \uses{def:translation-subgroup}
  \lean{AffineIncrEquiv.conjugate_homOfIndex₀}
  % \leanok
  Let $A^{(\beta)}_s = x + \beta s$ for $s, \beta \in \bR$ as in
  Definition~\ref{def:translation-subgroup}. Let
  also $B \in \OriAffR$ be given by $B(x) = a x + b$. Then
  \begin{align*}
    B \, A^{(\beta)}_s \, B^{-1} = A^{(a \beta)}_{s} .
  \end{align*}
\end{lemma}
\begin{proof}
  % \uses{}
  % \leanok
  Calculate, for $x \in \bR$
  \begin{align*}
    (B A^{(\beta)}_s B^{-1})(x) \; = \; & (B A^{(\beta)}_s)\big( \frac{x-b}{a} \big) \\
    \; = \; & B\big( \frac{x-b}{a} + \beta s \big) \\
    \; = \; & a \big( \frac{x-b}{a} + \beta s \big) + b \\
    \; = \; & x - b + a \beta s + b \\
    \; = \; & x + a \beta s \; = \; A^{(a \beta)}_s(x).
  \end{align*}
\end{proof}

\begin{definition}[Subgroup fixing a point]
  \label{def:fixing-subgroup}
  \uses{def:oriented-affine-isomorphism}
  \lean{AffineIncrEquiv.homOfIndex, AffineIncrEquiv.subGroupOfIndex}
  % \leanok
  The mapping $s \mapsto A_s$ with
  \begin{align*}
    A_s(x) = e^{s} (x - c) + c
  \end{align*}
  is a homomorphism $\bR \to \OriAffR$. The image of this
  homomorphism is the \emph{subgroup fixing $c$} in $\OriAffR$.
\end{definition}

\begin{lemma}[Characterization of the subgroup fixing a point]
  \label{lem:fixing-subgroup-characterization}
  \uses{def:fixing-subgroup}
  \lean{AffineIncrEquiv.mem_subGroupOfIndex_iff_fixed_point}
  \leanok
  An orientation-preserving affine transformation $A \in \OriAffR$
  belongs to the subgroup fixing $c \in \bR$ if and only if $A(c) = c$.

  \emph{(Note that the subgroup is a priori defined as the image of a homomorphism,
  so the statement indeed requires a proof.)}
\end{lemma}
\begin{proof}
  % \uses{}
  % \leanok
  Suppose first that $A$ is an element of the said subgroup, i.e.,
  $A(x) = e^{s} (x - c) + c$ for some $s \in \bR$. Then clearly $A(c) = c$.

  Suppose then that $A(c) = c$. Write $A(x) = a x + b$ for $a>0$ and $b \in \bR$.
  Plug in $x=c$ in the assumed fixed point property to obtain
  \begin{align*}
    a c + b = c .
  \end{align*}
  The above can be solved to give $b = (1-a)c$.
  Also since $a>0$, we can write $a = e^s$ with $s \in \bR$.
  With these, the formula for $A$ simplifies to
  \begin{align*}
    A(x) = e^s x + c (1 - e^s) = e^s \big( x - c \big) + c .
  \end{align*}
  This shows $A = A_s$ as desired (with $A_s$ as in Definition~\ref{def:fixing-subgroup}).
\end{proof}

\begin{lemma}[Conjugate of fixing is fixing image]
  \label{lem:conjugate-fixing}
  \uses{def:fixing-subgroup}
  \lean{AffineIncrEquiv.conjugate_homOfIndex}
  % \leanok
  Let $A^{(\alpha;c)}_s = e^{\alpha s} (x - c) + c$ for $\alpha, c \in \bR$ as in
  Definition~\ref{def:fixing-subgroup}. Let
  also $B \in \OriAffR$ be given by $B(x) = a x + b$. Then
  \begin{align*}
    B \, A^{(\alpha;c)}_s \, B^{-1} = A^{(\alpha;B(c))}_{s} .
  \end{align*}
\end{lemma}
\begin{proof}
  % \uses{}
  % \leanok
  Calculate, for $x \in \bR$
  \begin{align*}
    (B A^{(\alpha;c)}_s B^{-1})(x) \; = \; & (B A^{(\alpha;c)}_s)\big( \frac{x-b}{a} \big) \\
    \; = \; & B\big( e^{\alpha s} (\frac{x-b}{a} - c) + c \big) \\
    \; = \; & a e^{\alpha s} (\frac{x-b}{a} - c) + a c + b \\
    \; = \; & e^{\alpha s} \big( x - (a c + b) \big) + (a c + b) \\
    \; = \; & e^{\alpha s} \big( x - B(c) \big) + B(c) \\
    \; = \; & A^{(\alpha;B(c))}_s (x) .
  \end{align*}
\end{proof}


\section{Self-similarity characterizations of the extreme value distributions}

\begin{lemma}[Continuous parameter extreme value limit relation]
  \label{lem:continuous-parameter-ev-limit}
  \uses{def:cdf, lem:oriented-affine-action-on-cdf}
  \lean{continuous_parameter_ev_limit_relation}
  \leanok
  Let $F$ be a c.d.f.

  \emph{(Note that below we use
  the sequence $(F^n)_{n \in \bN}$ of $n$th powers of a fixed c.d.f.,
  not a sequence of arbitrary c.d.f.s.
  Recall that the $n$th power $F^n$ is the c.d.f. of the maximum
  of $n$ independent random variables with the distribution $F$.)}

  Suppose that for a sequence $(A_n)_{n \in \bN}$ of oriented
  affine isomorphisms of $\bR$, $A_n \in \OriAffR$, we have
  \begin{align*}
    A_n.F^n \tendsinlaw G ,
  \end{align*}
  where $G$ is a c.d.f.

  (In Lean the statement is pointwise: if $\big((A_n.F)(x)\big)^n \to G(x)$ at a point $x$,
  then $\big((A_n.F)(x)\big)^{\max(1, \lfloor n t \rfloor)} \to G(x)^t$.)

  Then, for any $t > 0$, denoting by $G^t$ the c.d.f. given by
  $G^t(x) = \big( G(x) \big)^t$, we have
  \begin{align*}
    A_{n}.F^{\lfloor n t \rfloor} \tendsinlaw G^t ,
  \end{align*}
  where, for $x \in \bR$, the floor notation $\lfloor x \rfloor$ stands
  for the greatest integer $k \in \bZ$ such that $k \le x$.
\end{lemma}
\begin{proof}
  % \uses{}
  \leanok
  Let $t > 0$ and let $x \in \bR$ be a continuity point of $G$.
  For $n \in \bN$, calculate
  \begin{align*}
      \big((A_n . F) (x) \big)^{\lfloor n t \rfloor}
    %\; = \; \big(F(A_n^{-1}(x)) \big)^{\lfloor n t \rfloor} \\
    \; = \; & \Big( \big((A_n . F) (x) \big)^n \Big)^{\lfloor n t \rfloor / n} .
  \end{align*}
  By assumption, we have $\big((A_n . F) (x) \big)^n \to G(x)$ as $n \to \infty$.
  Also $\lfloor n t \rfloor / n \to t$ as $n \to \infty$.
  By (joint) continuity of the power function
  $(x,y) \mapsto x^y = \exp\big( y \, \log(x) \big)$,
  we get that the expression above tends to $G(x)^t$.

  Finally noting that the continuity points of $G^t$ are the same as the
  continuity points of $G$, the above in fact proves the asserted
  $A_{n}.F^{\lfloor n t \rfloor} \tendsinlaw G^t$.
\end{proof}

\begin{lemma}[Self-similarity of extreme value distributions]
  \label{lem:self-similarity-of-extreme-value-distributions}
  \uses{def:cdf, lem:oriented-affine-action-on-cdf, def:extr-val-distr}
  \lean{CumulativeDistributionFunction.IsExtremeValueDistr.self_similar}
  \leanok
  %Let $F$ be a c.d.f.
  %
  %\emph{(Note that below we use
  %the sequence $(F^n)_{n \in \bN}$ of $n$th powers of a fixed c.d.f.,
  %not a sequence of arbitrary c.d.f.s.
  %Recall that the $n$th power $F^n$ is the c.d.f. of the maximum
  %of $n$ independent random variables with the distribution $F$.)}
  %
  %Suppose that for a sequence $(A_n)_{n \in \bN}$ of oriented
  %affine isomorphisms of $\bR$, $A_n \in \OriAffR$, we have
  %\begin{align*}
  %  A_n.F^n \tendsinlaw G ,
  %\end{align*}
  %where $G$ is a c.d.f.
  Suppose that $G$ is an extreme-value distribution.
  Then there exists a family $(A_t)_{t > 0}$ of
  oriented affine isomorphisms of $\bR$, $A_t \in \OriAffR$,
  such that for any $t > 0$
  \begin{align*}
    G^t = A_t . G .
  \end{align*}

  Moreover, $t \mapsto A_t$ is a measurable homomorphism
  of multiplicative groups $(0,+\infty) \to \OriAffR$.
\end{lemma}
\begin{proof}
  \uses{lem:continuous-parameter-ev-limit, thm:convergence-to-types}
  % \leanok
  \ldots
\end{proof}

%\begin{lemma}[Continuous parameter limit relation]
%  \label{lem:continuous-parameter-ev-limit}
%  \uses{def:cdf, lem:oriented-affine-action-on-cdf}
%  % \lean{}
%  % \leanok
%  Let $F$ be a c.d.f.
%
%  Suppose that for a sequence $(A_n)_{n \in \bN}$ of oriented
%  affine isomorphisms of $\bR$, and for a family $(G_t)_{t > 0}$
%  of c.d.f.s, \ldots.
%
%
%  and any $t > 0$, we have
%  we have
%  \begin{align*}
%    A_n.F^n \tendsinlaw G_t ,
%  \end{align*}
%  where $G_t$ is a nondegenerate c.d.f.
%
%  (In Lean the statement is pointwise: if $\big((A_n.F)(x)\big)^n \to G(x)$ at a point $x$,
  then $\big((A_n.F)(x)\big)^{\max(1, \lfloor n t \rfloor)} \to G(x)^t$.)

  Then, for any $t > 0$, denoting by $G^t$ the c.d.f. given by
%  $G^t(x) = \big( G(x) \big)^t$, we have
%  \begin{align*}
%    A_{n}.F^{\lfloor n t \rfloor} \tendsinlaw G^t .
%  \end{align*}
%\end{lemma}

%%\begin{lemma}[Functional equation in one parameter subgroups of affine isomorphisms]
%%  \label{lem:affine-one-parameter-subgroup}
%%  \uses{def:lem:oriented-affine-action-on-cdf}
%%  % \lean{}
%%  % \leanok
%%  Suppose that $(A_t)_{t > 0}$ is a family of affine isomorphisms of $\bR$
%%  such that
%%  \begin{align*}
%%    A_{s t} = A_s \circ A_t  \qquad \text{ for any } s, t > 0 .
%%  \end{align*}
%%  Write $A_t(x) = a_t x + b_t$, with $a_t > 0$ and $b_t \in \bR$.
%%  Then we have, for any $s, t > 0$,
%%  \begin{align*}
%%    a_{t s} = \; & a_t \, a_s
%%      \qquad \text{ and } \\
%%    b_{t s} = \; & a_t \, b_s + b_t .
%%  \end{align*}
%%  (Also by symmetry $b_{t s} = a_s \, b_t + b_s$.)
%%\end{lemma}
%%\begin{proof}
%%  % \uses{}
%%  % \leanok
%%\end{proof}
%%
%%\begin{lemma}[Functional equation scaling coefficient solution]
%%  \label{lem:solution-functional-eqn-scaling}
%%  %\uses{def:lem:oriented-affine-action-on-cdf}
%%  % \lean{}
%%  % \leanok
%%  Suppose that $a \colon (0,+\infty) \to (0,+\infty)$ is a
%%  measurable function satisfying, for any $s, t > 0$,
%%  \begin{align*}
%%    a(t s) = \; & a(t) \, a(s) .
%%  \end{align*}
%%  Then there exists a $\theta \in \bR$ such that for all $t > 0$,
%%  \begin{align*}
%%    a(t) = t^{\theta} .
%%  \end{align*}
%%\end{lemma}
%%\begin{proof}
%%  % \uses{}
%%  % \leanok
%%\end{proof}
%%
%%\begin{lemma}[Functional equation translation coefficient solution with $\theta = 0$]
%%  \label{lem:solution-functional-eqn-translation-no-scaling}
%%  %\uses{def:lem:oriented-affine-action-on-cdf}
%%  % \lean{}
%%  % \leanok
%%  Suppose that
%%  $b \colon (0,+\infty) \to \bR$ is a measurable function
%%  satisfying, for any $s, t > 0$,
%%  \begin{align*}
%%    b(t s) = \; & b(s) + b(t) .
%%  \end{align*}
%%  Then there exists a constant $c$ such that for $t \in (0,+\infty)$
%%  we have
%%  \begin{align*}
%%    b(t) = \; & -c \, \log (t) .
%%  \end{align*}
%%\end{lemma}
%%\begin{proof}
%%  % \uses{}
%%  % \leanok
%%\end{proof}
%%
%%\begin{lemma}[Functional equation translation coefficient solution with $\theta \ne 0$]
%%  \label{lem:solution-functional-eqn-translation-with-scaling}
%%  %\uses{def:lem:oriented-affine-action-on-cdf}
%%  % \lean{}
%%  % \leanok
%%  Suppose that $\theta \in \bR \setminus 0$ and
%%  $b \colon (0,+\infty) \to \bR$ is a measurable function
%%  satisfying, for any $s, t > 0$,
%%  \begin{align*}
%%    b(t s) = \; & t^{\theta} \, b(s) + b(t) .
%%  \end{align*}
%%  Then there exists a constant $c$ such that for $t \in (0,+\infty) \setminus \set{1}$
%%  we have
%%  \begin{align*}
%%    b(t) = \; & c (1 - t^{-\theta}) .
%%  \end{align*}
%%\end{lemma}
%%\begin{proof}
%%  % \uses{}
%%  % \leanok
%%\end{proof}

\begin{lemma}[Self-similar continuous c.d.f. family characterization $\gamma = 0$]
  \label{lem:characterization-self-similar-family-zero-index}
  \uses{def:lem:oriented-affine-action-on-cdf, def:translation-subgroup}
  \lean{gumbel_type_of_selfSimilar_index_zero', gumbel_type_of_selfSimilar_index_zero}
  \leanok
  Suppose that $G$ is a nondegenerate c.d.f. such that
  \begin{align*}
    G^{t} = A_t . G \qquad \text{ for any } t > 0 ,
  \end{align*}
  where
  \begin{align*}
    A_t(x) = x + \beta \, \log t
  \end{align*}
  with $\beta > 0$.

  Then with $d = \log \big(-\log G(0) \big)$,
  for all $x \in \bR$ we have
  \begin{align*}
    G(x) = \exp \Big( - \exp \big( -\beta^{-1} x + d \big) \Big) .
  \end{align*}
  (In particular, $G$ is of Gumbel type: there exists an $A \in \OriAffR$
  such that $G = A.\GumbelCDF$.)
\end{lemma}
\begin{proof}
  %\uses{lem:solution-functional-eqn-translation-no-scaling}
  \uses{lem:no-fixed-point-implies-translation}
  % \leanok

  Since $G$ is nondegenerate, there exists an $x_0 \in \bR$
  with $0 < G(x_0) < 1$. Write $q = -\log G(x_0) > 0$,
  so that $G(x_0) = e^{-q}$.

  For $t > 0$, from the equation $G^{t} = A_t . G$
  we get for any $x \in \bR$
  \begin{align*}
  G(x)^t = G \big( A_t^{-1}(x) \big) = G \big( x - \beta \, \log(t) \big) .
  \end{align*}
  In particular, with $x = x_0 + \beta \, \log(t)$, we get
  \begin{align*}
    G\big( x_0 + \beta \, \log(t) \big)^{t} = G(x_0) = e^{- q},
  \end{align*}
  from which we can solve
  \begin{align*}
    G\big( x_0 + \beta \, \log(t) \big) = e^{-q/t} .
  \end{align*}
  The above holds for any $t>0$, and any $x \in \bR$
  can be written as $x = x_0 + \beta \, \log\big( e^{(x - x_0)/\beta} \big)$.
  We therefore get, for any $x \in \bR$,
  \begin{align*}
    G(x) = \exp \Big( - q \, \exp \big( -(x-x_0)/\beta \big) \Big) .
  \end{align*}
  This is of the desired form, with $d = \frac{x_0}{\beta} + \log(q)$.
  Plugging in $x = 0$ then shows $d = \log \big(-\log G(0) \big)$.
\end{proof}

\begin{lemma}[Self-similar continuous c.d.f. family characterization $\gamma > 0$]
  \label{lem:characterization-self-similar-family-pos-index}
  \uses{def:lem:oriented-affine-action-on-cdf, def:fixing-subgroup}
  \lean{frechet_type_of_selfSimilar_index_pos', frechet_type_of_selfSimilar_index_pos}
  \leanok
  Suppose that $G$ is a nondegenerate c.d.f. such that
  \begin{align*}
    G^{t} = A_t . G \qquad \text{ for any } t > 0 ,
  \end{align*}
  where
  \begin{align*}
    A_t(x) = t^{\alpha} x + c \, (1 - t^{\alpha}) = t^{\alpha} (x-c) + c
  \end{align*}
  with $c \in \bR$ and $\alpha > 0$.

  Then with $\sigma = \big(- \log G(c+1)\big)^{\alpha}$,
  for all $x \ge c$ we have
  \begin{align*}
    G(x) = \exp \Big( - \big( \frac{x - c}{\sigma} \big)^{-1/\alpha} \big) \Big) .
  \end{align*}

  (It easily follows that $G$ is of Fr\'echet type: there exists
  an $A \in \OriAffR$ such that $G = A.\FrechetCDF{\alpha}$.)
\end{lemma}
\begin{proof}
  %\uses{lem:solution-functional-eqn-translation-with-scaling}
  \uses{lem:fixing-subgroup-characterization, lem:cdf-equal-on-dense}
  % \leanok

  Since $G$ is nondegenerate, there exists an $x_0 \in \bR$
  with $0 < G(x_0) < 1$. Write $q = -\log G(x_0) > 0$,
  so that $G(x_0) = e^{-q}$.

  %Note that
  %\begin{align*}
  %  A_t^{-1}(x) = t^{-\alpha} x + c (1 - t^{-\alpha}) = t^{-\alpha} (x-c) + c .
  %\end{align*}
  For $t > 0$, from the equation $G^{t} = A_t . G$
  we get for any $x \in \bR$
  \begin{align*}
  G(x)^t = G \big( A_t^{-1}(x) \big) . %= G \big( t^{-\alpha} (x-c) + c \big) .
  \end{align*}

  Note that for $x \le c$ and $t = 2$ we have $A_2^{-1}(x) = 2^{-\alpha}(x-c) + c \ge x$,
  so the above implies $G(x)^2 = G(A_2^{-1}(x)) \ge G(x)$.
  This not possible if $0 < G(x) < 1$, so we must in particular have $x_0 > c$.

  With $x = A_t(x_0)$ in the above equation, %= t^{\alpha} (x_0-c) + c$,
  we get
  \begin{align*}
    G\big( A_t(x_0) \big)^{t} = G(x_0) = e^{- q},
  \end{align*}
  from which we can solve
  \begin{align*}
    G(x) = G\big( A_t(x_0) \big) = e^{-q/t} .
  \end{align*}

  %The above holds for any $t>0$.
  Any $x \ge c$ can be written as $x = A_t(x_0) = t^{\alpha} (x_0-c) + c$
  with $t = \big( \frac{x - c}{x_0 - c} \big)^{1/\alpha}$
  (recall that $x_0 - c > 0$).
  We therefore get, for any $x \ge c$,
  \begin{align*}
    G(x) = \exp \Big( - q \, \big( \frac{x - c}{x_0 - c} \big)^{-1/\alpha} \big) \Big) .
  \end{align*}
  This is of the form
  \begin{align*}
    G(x) = \exp \Big( - \big( \frac{x - c}{\sigma} \big)^{-1/\alpha} \big) \Big) ,
  \end{align*}
  and plugging in $x = c + 1$ yields $\sigma = \big(- \log G(c+1)\big)^{\alpha}$.
  %This is of the desired form, with $d = \frac{x_0}{c} + \log(q)$.
  %Plugging in $x = 0$ then shows $d = \log \big(-\log G(0) \big)$.
\end{proof}

\begin{lemma}[Self-similar continuous c.d.f. family characterization $\gamma < 0$]
  \label{lem:characterization-self-similar-family-neg-index}
  \uses{def:lem:oriented-affine-action-on-cdf, def:fixing-subgroup}
  \lean{weibull_type_of_selfSimilar_index_neg', weibull_type_of_selfSimilar_index_neg}
  \leanok
  Suppose that $G$ is a nondegenerate c.d.f. such that
  \begin{align*}
    G^{t} = A_t . G \qquad \text{ for any } t > 0 ,
  \end{align*}
  where
  \begin{align*}
    A_t(x) = t^{-\alpha} x + c \, (1 - t^{-\alpha}) = t^{-\alpha} (x-c) + c
  \end{align*}
  with $c \in \bR$ and $\alpha > 0$.

  Then with $\sigma = \big(- \log G(c-1)\big)^{-\alpha}$,
  for all $x \le c$ we have
  \begin{align*}
    G(x) = \exp \Big( - \big( \frac{c - x}{\sigma} \big)^{1/\alpha} \big) \Big) .
  \end{align*}

  (It easily follows that $G$ is of Weibull type: there exists
  an $A \in \OriAffR$ such that $G = A.\WeibullCDF{\alpha}$.)
\end{lemma}
\begin{proof}
  %\uses{lem:solution-functional-eqn-scaling,
  %  lem:solution-functional-eqn-translation-exp-zero}
  \uses{lem:fixing-subgroup-characterization, lem:cdf-equal-on-dense}
  % \leanok

  \emph{(Note: The Lean formalized statement uses the opposite sign of $\alpha$:
  it is assumed that $\alpha < 0$ and $A_t(x) = t^{+\alpha} (x-c) + c$.)}

  Since $G$ is nondegenerate, there exists an $x_0 \in \bR$
  with $0 < G(x_0) < 1$. Write $q = -\log G(x_0) > 0$,
  so that $G(x_0) = e^{-q}$.

  %Note that
  %\begin{align*}
  %  A_t^{-1}(x) = t^{-\alpha} x + c (1 - t^{-\alpha}) = t^{-\alpha} (x-c) + c .
  %\end{align*}
  For $t > 0$, from the equation $G^{t} = A_t . G$
  we get for any $x \in \bR$
  \begin{align*}
  G(x)^t = G \big( A_t^{-1}(x) \big) . %= G \big( t^{-\alpha} (x-c) + c \big) .
  \end{align*}

  Note that for $x \ge c$ and $t = 2$ we have $A_2^{-1}(x) = 2^{\alpha}(x-c) + c \ge x$,
  so the above implies $G(x)^2 = G(A_2^{-1}(x)) \ge G(x)$.
  This not possible if $0 < G(x) < 1$, so we must in particular have $x_0 < c$.

  With $x = A_t(x_0)$ in the above equation, %= t^{\alpha} (x_0-c) + c$,
  we get
  \begin{align*}
    G\big( A_t(x_0) \big)^{t} = G(x_0) = e^{- q},
  \end{align*}
  from which we can solve
  \begin{align*}
    G(x) = G\big( A_t(x_0) \big) = e^{-q/t} .
  \end{align*}

  %The above holds for any $t>0$.
  Any $x \le c$ can be written as $x = A_t(x_0) = t^{-\alpha} (x_0-c) + c$
  with $t = \big( \frac{c - x}{c - x_0} \big)^{-1/\alpha}$
  (recall that $c - x_0 > 0$).
  We therefore get, for any $x \le c$,
  \begin{align*}
    G(x) = \exp \Big( - q \, \big( \frac{c - x}{c - x_0} \big)^{1/\alpha} \big) \Big) .
  \end{align*}
  This is of the form
  \begin{align*}
    G(x) = \exp \Big( - \big( \frac{c - x}{\sigma} \big)^{1/\alpha} \big) \Big) ,
  \end{align*}
  and plugging in $x = c - 1$ yields $\sigma = \big(- \log G(c-1)\big)^{-\alpha}$.
  %This is of the desired form, with $d = \frac{x_0}{c} + \log(q)$.
  %Plugging in $x = 0$ then shows $d = \log \big(-\log G(0) \big)$.
\end{proof}


\begin{theorem}[{Three types of extreme value distributions [Fisher-Tippett-Gnedenko]}]
  \label{thm:three-types-of-extr-val-distr}
  \uses{def:extr-val-distr,
    def:std-Gumbel-cdf, def:std-Frechet-cdf, def:std-Weibull-cdf}
  \lean{CumulativeDistributionFunction.IsExtremeValueDistr.classification}
  \leanok
  For any extreme value distribution $G$, one of the following holds:
  \begin{description}
    \item{($\GumbelCDF$)}
      $G = A . \GumbelCDF$ for some $A \in \OriAffR$;
    \item{($\FrechetCDF{{}}$)}
      $G = A . \FrechetCDF{\alpha}$ for some $A \in \OriAffR$ and $\alpha > 0$;
    \item{($\WeibullCDF{{}}$)}
      $G = A . \WeibullCDF{\alpha}$ for some $A \in \OriAffR$ and $\alpha > 0$.
  \end{description}
  In particular, the only three possible types of extreme value distributions
  are the type of the Gumbel c.d.f., the type of
  the Fr\'echet c.d.f. $\FrechetCDF{\alpha}$ for $\alpha > 0$, and
  the type of the Weibull c.d.f. $\WeibullCDF{\alpha}$ for $\alpha > 0$.
\end{theorem}
\begin{proof}
  \uses{thm:one-parameter-subgroups-of-affine-isomorphisms,
    lem:self-similarity-of-extreme-value-distributions,
    lem:characterization-self-similar-family-zero-index,
    lem:characterization-self-similar-family-pos-index,
    lem:characterization-self-similar-family-neg-index}
  \leanok
  \ldots
\end{proof}


\chapter{Left-continuous inverses}

Let $R$ and $S$ be complete linear orders, for example $[0,1]$, $[0,+\infty]$, or
$[-\infty,+\infty] =: \overline{\bR}$.

\begin{definition}
  \label{def:lc-inverse}
  \lean{lcInv}
  \leanok
  Let $f \colon R \to S$ be a function (usually assumed nondecreasing).
  The left-continuous inverse of $f$ is the function $\lcinv{f} \colon S \to R$
  given by
  \begin{align*}
    \lcinv{f}(y) := \inf \big\{ x \in R \; \big| \; f(x) \ge y \big\}
    \qquad \text{ for } y \in S.
  \end{align*}
  The right-continuous inverse $\rcinv{f} \colon S \to R$ is analoguously defined by
  \begin{align*}
    \rcinv{f}(y) := \sup \big\{ x \in R \; \big| \; f(x) \le y \big\}
    \qquad \text{ for } y \in S.
  \end{align*}
\end{definition}


\chapter{Cauchy-Hamel functional equation}

\section{Positive measure additive subgroups of the reals}

\begin{lemma}[Countably many connected components for an open set]
  \label{lem:countably-many-connected-components-of-open}
  %\uses{}
  \lean{IsOpen.countable_setOf_connectedComponentIn}
  \leanok
  Let $X$ be a locally connected separable space. Then any open subset $U \subseteq X$ has at
  most countably many connected components.
\end{lemma}
\begin{proof}
  % \uses{}
  \leanok
  (The proof is already formalized, see:
  \texttt{IsOpen.countable-setOf-connectedComponentIn}.)
\end{proof}

For subsets $A, B \subseteq \bR$, we use the following notation for pointwise
sum sets and difference sets:
\begin{align*}
  A + B & = \set{ a + b \; \big| \; a \in A, \, b \in B } , \\
  A - B & = \set{ a - b \; \big| \; a \in A, \, b \in B } .
\end{align*}

We denote by $\Leb$ the Lebesgue measure on $\bR$.

\begin{lemma}[Finding an interval with high overlap]
  \label{lem:exists-high-overlap-interval}
  %\uses{}
  \lean{exists_interval_measure_inter_gt_mul_measure}
  \leanok
  Let $A \subset \bR$ be a measurable set such that $0 < \Leb[A] < +\infty$.
  Then for any $r \in [0,1)$, there exists a nontrivial interval $J \subset \bR$
  (a subset of the real line which is connected and has nonempty interior)
  such that
  \begin{align*}
    \Leb[A \cap J] \; > \; r \, \Leb[J] .
  \end{align*}
\end{lemma}
\begin{proof}
  % \uses{}
  \leanok
  (The Lean proof instead takes a Lebesgue density point $x$ of $A$ and a small
  closed interval $J = [x-\eps, x+\eps]$.)

  Assume, without loss of generality, $0 < r < 1$.
  Since the Lebesgue measure is outer regular, we can find an open set $U \subset \bR$
  such that $A \subseteq U$ and $\Leb[U] < r^{-1} \, \Leb[A]$.

  The open set $U$ has at most countably many connected components
  (which are in fact open intervals); denote
  by $(U_i)_{i \in I}$ the indexed collection of them.

  Note that for at least one index $j \in I$ we
  have $\Leb[A \cap U_j] > r \Leb[U_j]$,
  because otherwise we get
  \begin{align*}
    \Leb[A] = \Leb[A \cap U]
    = \, & \sum_{i \in I} \Leb[A \cap U_i] \\
    \le \, & r \sum_{i \in I} \Leb[U_i] \\
    = \, & r \Leb[U],
  \end{align*}
  contradicting the choice of $U$.

  Now $J = U_j$ has the desired properties.
\end{proof}

\begin{lemma}[Shifts of a smaller interval contained in a larger interval]
  \label{lem:shift-interval-contained-in-larger}
  %\uses{}
  %\lean{}
  % \leanok % needs also `isPreconnected_of_Ioo_subset_of_subset_Icc` etc. before really done
  Let $I, J$ be nontrivial intervals, whose length satisfy
  \begin{align*}
    0 \; < \; \Leb[J] \; < \; \Leb[I] \; < \, +\infty .
  \end{align*}
  Then there exists an open interval $\Delta$ of length
  $\Leb[\Delta] = \Leb[I] - \Leb[J] > 0$ such that
  \begin{align*}
    t + J \; \subset \; I
    \qquad \text{ for any $t \in \Delta$.}
  \end{align*}
\end{lemma}
\begin{proof}
  % \uses{}
  % \leanok
  \ldots
\end{proof}

\begin{lemma}[Overlapping union of copies of an interval]
  \label{lem:overlapping-union-interval-copies}
  %\uses{}
  \lean{volume_union_add_self_ge_of_Ioo_subset, volume_union_add_self_le_of_subset_Icc,
    isConnected_of_Ioo_subset_of_subset_Icc}
  % \leanok % needs also `isPreconnected_of_Ioo_subset_of_subset_Icc` etc. before really done
  Let $J$ be a nontrivial interval of finite length ($0 \, < \, \Leb[J] \, < +\infty$).
  Let $c < 1$ and denote $\delta = c \Leb[J] \, < \, \Leb[J]$.
  Then for any $t \in (-\delta,\delta)$, the set
  $J' = (t+J) \cup J$ is an interval (connected set with nonempty interior)
  whose length satisfies the bound $\Leb[J'] \, < \, (1+c) \, \Leb[J]$.
\end{lemma}
\begin{proof}
  \uses{lem:shift-interval-contained-in-larger}
  % \leanok
  Denote $a = \inf J$ and $b = \sup J$. We have $-\infty < a < b < +\infty$ and
  \begin{align*}
    (a,b) \; \subseteq \; J \; \subseteq \; [a,b]
  \end{align*}
  and we have $\Leb[J] = b - a$.

  Let $t \in (+\delta,\delta)$, with $\delta = c \, \Leb[J] = c(b-a)$.
  Then we have the containments
  \begin{align*}
    \big( \min \{a,a+t\} , \max \{b,b+t\} \big)
      \; \subseteq \; (t+J) \cup J
      \; \subseteq \; \big[ \min \{a,a+t\} , \max \{b,b+t\} \big] .
  \end{align*}
  It in particular follows that $J' := (t+J) \cup J$ is an interval.

  If $t \ge 0$, then $J'$ is contained in
  $[a,b+t]$ which has length $b+t-a = (b-a) + t < \Leb[J] + c \, \Leb[J] = (1+c) \Leb[J]$.
  If $t < 0$, then $J'$ is contained in $[a+t,b]$ and one similarly gets a length bound.
\end{proof}


\begin{lemma}[Difference set of positive measure set contains an interval]
  \label{lem:difference-set-contains-interval}
  %\uses{}
  \lean{exists_Ioo_subset_diff_self_of_measure_pos}
  \leanok
  Let $A \subset \bR$ be a measurable set of positive Lebesgue measure.
  Then there exists a $\delta > 0$ such that
  \begin{align*}
    (-\delta,\delta) \, \subseteq \, A - A .
  \end{align*}
\end{lemma}
\begin{proof}
  \uses{lem:countably-many-connected-components-of-open, lem:exists-high-overlap-interval,
    lem:overlapping-union-interval-copies}
  % \leanok
  Pick a measurable subset $A_0 \subseteq A$ such that
  $0 < \Leb[A_0] < +\infty$.

  By Lemma~\ref{lem:exists-high-overlap-interval}, we can find a nontrivial interval~$J$
  such that $\Leb[A_0 \cap J] > \frac{3}{4} \Leb[J]$.
  Let $\delta = \frac{1}{2} \Leb[J]$.
  We claim that $(-\delta,\delta) \subseteq A_0 - A_0$ (which then
  clearly implies the assertion of the lemma).
  Indeed, suppose $t \in (-\delta,\delta)$.
  Then by Lemma~\ref{lem:overlapping-union-interval-copies},
  $(t+J) \cup J$ is an interval of length less than $\frac{3}{2} \Leb[J]$.
  Moreover, we have $A_0 \cap J \subseteq (t+J) \cup J$ and
  $t+(A_0 \cap J) \subseteq (t+J) \cup J$. Now note that
  \begin{align*}
    \Leb[t + (A_0 \cap J)] = \Leb[A_0 \cap J] > \frac{3}{4} \Leb[J] .
  \end{align*}
  If the sets $t + (A_0 \cap J)$ and $(A_0 \cap J)$ were disjoint, then
  the measure of $(A_0 \cap J) \cup (t+(A_0 \cap J))$ would thus be
  greater than $\frac{3}{2} \Leb[J]$, which is impossible given the length
  of the interval $(t+J) \cup J$ is less than $\frac{3}{2} \Leb[J]$.
  Therefore there exists a point $a \in (t + (A_0 \cap J)) \cap (A_0 \cap J)$.
  Denoting $a' = a - t$, we have $a, a' \in A_0 \cap J \subseteq A_0$
  and $t = a - a' \in A_0 - A_0 \subseteq A - A$. Since $t \in (-\delta,\delta)$
  was arbitrary, this proves the assertion.
\end{proof}

\begin{lemma}[Dividing a high overlap interval]
  \label{lem:dividing-high-operlap-interval}
  %\uses{}
  % \lean{}
  % \leanok
  Let $A$ be a measurable. Suppose that for some $a < b$,
  the interval $J = (a,b)$ satisfies
  $\Leb[A \cap J] \; > \;  r \, \Leb[J]$. %Denote $\delta = \Leb[J] = b-a$.
  Let $m \in \Npos$, and consider the subintervals
  \begin{align*}
    J_i = \Big( a + \frac{i}{(b-a)m} , \; a + \frac{i+1}{(b-a)m} \Big)
    \qquad \text{ for $i = 0, \ldots, m-1$.}
  \end{align*}
  Then for some $i$ we have $\Leb[A \cap J_i] \; > \;  r \, \Leb[J_i]$.
\end{lemma}
\begin{proof}
  % \uses{}
  % \leanok
\end{proof}

\begin{lemma}[Difference of two positive measure sets contains an interval]
  \label{lem:different-difference-set-contains-interval}
  % \uses{}
  % \lean{}
  % \leanok
  Let $A, B \subset \bR$ be two measurable sets of positive Lebesgue measure.
  Then the difference set $A - B$ contains a nontrivial open interval.
  %\begin{align*}
  %  (-\delta,\delta) \, \subseteq \, A - B .
  %\end{align*}
\end{lemma}
\begin{proof}
  \uses{lem:dividing-high-operlap-interval}
  % \leanok
  Pick measurable subsets $A_0 \subseteq A$ and $B_0 \subseteq B$ such that
  $0 < \Leb[A_0] < +\infty$ and $0 < \Leb[B_0] < +\infty$.

  By Lemma~\ref{lem:exists-high-overlap-interval}, we can find nontrivial intervals~$J, I$
  such that $\Leb[A_0 \cap J] > \frac{3}{4} \Leb[J]$
  and $\Leb[B_0 \cap I] > \frac{3}{4} \Leb[I]$.

  By Lemma~\ref{lem:dividing-high-operlap-interval}
  % we may divide
  % $I$ and $J$ to equal length subintervals, at least one of which inherits
  % this property, i.e.,
  for any $n,m \in \bN$ we can find intervals
  $I' \subset I$ and $J' \subset J$
  with $\Leb[I'] = \frac{1}{n} \Leb[I]$ and $\Leb[J'] = \frac{1}{m} \Leb[J]$
  and such that $\Leb[A_0 \cap J'] > \frac{3}{4} \Leb[J']$
  and $\Leb[B_0 \cap I'] > \frac{3}{4} \Leb[I']$.
  By choosing $n$ and $m$ suitably, we can ensure that
  \begin{align*}
    \frac{1}{2} \Leb[I'] \; \le \; \Leb[J'] \; < \; \Leb[I'] .
  \end{align*}

  Then by Lemma~\ref{lem:shift-interval-contained-in-larger}
  there exists an open interval $\Delta$ of length
  $\Leb[\Delta] \, > \, \Leb[I'] - \Leb[J']$ such that
  for all $t \in \Delta$ we have $t + J' \subset I'$.

  Consider a fixed $t \in \Delta$.
  Observe that $t + (B_0 \cap J) \subset I'$ and
  \begin{align*}
    \Leb[t + (B_0 \cap J')] \; = \; \Leb[B_0 \cap J']
      \; > \; \frac{3}{4} \Leb[J'] \; \ge \; \frac{3}{8} \Leb[I'] .
  \end{align*}
  We now claim that $A_0 \cap I'$ and $t + (B_0 \cap J')$ intersect.
  Their measures add up to at least
  \begin{align*}
    \Leb[A_0 \cap I'] + \Leb[t + (B_0 \cap J')]
    \; = \; & \Leb[A_0 \cap I'] + \Leb[B_0 \cap J'] \\
    \; > \; & \frac{3}{4} \Leb[I'] + \frac{3}{4} \Leb[J'] \\
    \; \ge \; & \frac{3}{4} \Leb[I'] + \frac{3}{8} \Leb[I']
    \; = \; & \frac{9}{8} \Leb[I'] .
  \end{align*}
  and both have been shown to be subsets of $I'$; therefore they
  cannot be disjoint.
  In particular there is a point $z \in (A_0 \cap I') \cap (t + B_0 \cap J')$,
  which means that
  \begin{align*}
    a = z = t + b
  \end{align*}
  for some $a \in A_0 \cap I'$ and $b \in B_0 \cap J'$. We solve
  $t = a - b \in A_0 - B_0 \subseteq A - B$. Since $t \in \Delta$
  was arbitrary, we have shown that the nontrivial open interval $\Delta$
  is contained in $A - B$.
\end{proof}


\chapter{Regular variation}

This chapter develops the basic theory of regularly varying functions,
following Bingham, Goldie and Teugels, \emph{Regular Variation} (Cambridge
University Press, 1987), henceforth [BGT]. All limits are taken as $x \to +\infty$.

\section{Definitions and basic properties}

\begin{definition}[Slowly varying function]
  \label{def:slowly-varying}
  \lean{ExtremeValueTheory.IsSlowlyVarying}
  \leanok
  A function $L \colon \bR \to \bR$ is \emph{slowly varying} (at $+\infty$)
  if $L(x) > 0$ for all sufficiently large $x$ and
  \begin{align*}
    \lim_{x \to \infty} \frac{L(c x)}{L(x)} = 1 \qquad \text{for every } c > 0 .
  \end{align*}
\end{definition}

\begin{definition}[Regularly varying function]
  \label{def:regularly-varying}
  \lean{ExtremeValueTheory.IsRegularlyVarying}
  \leanok
  Let $\rho \in \bR$. A function $f \colon \bR \to \bR$ is \emph{regularly varying}
  (at $+\infty$) \emph{with index $\rho$} if $f(x) > 0$ for all sufficiently large $x$ and
  \begin{align*}
    \lim_{x \to \infty} \frac{f(c x)}{f(x)} = c^\rho \qquad \text{for every } c > 0 .
  \end{align*}
\end{definition}

\begin{lemma}[Slow variation is regular variation of index zero]
  \label{lem:slowly-varying-iff-rv-zero}
  \uses{def:slowly-varying, def:regularly-varying}
  \lean{ExtremeValueTheory.isSlowlyVarying_iff_isRegularlyVarying_zero}
  \leanok
  A function is slowly varying if and only if it is regularly varying with index $0$.
\end{lemma}
\begin{proof}
  \leanok
  Immediate from the definitions, since $c^0 = 1$.
\end{proof}

\begin{lemma}[Germ invariance]
  \label{lem:rv-congr}
  \uses{def:regularly-varying}
  \lean{ExtremeValueTheory.IsRegularlyVarying.congr}
  \leanok
  If $f$ is regularly varying with index $\rho$ and $g(x) = f(x)$ for all
  sufficiently large $x$, then $g$ is regularly varying with index $\rho$.
\end{lemma}
\begin{proof}
  \leanok
  For fixed $c > 0$, both $x$ and $cx$ are eventually in the region where $f = g$,
  so the ratios $g(cx)/g(x)$ and $f(cx)/f(x)$ eventually agree.
\end{proof}

\begin{lemma}[Basic examples]
  \label{lem:rv-examples}
  \uses{def:regularly-varying, def:slowly-varying}
  \lean{ExtremeValueTheory.IsRegularlyVarying.rpow_id, ExtremeValueTheory.IsRegularlyVarying.const,
    ExtremeValueTheory.IsSlowlyVarying.const, ExtremeValueTheory.isSlowlyVarying_log}
  \leanok
  The power function $x \mapsto x^\rho$ is regularly varying with index $\rho$;
  every positive constant is slowly varying; and the logarithm $\log$ is slowly varying.
\end{lemma}
\begin{proof}
  \leanok
  For $x > 0$ we have $(cx)^\rho / x^\rho = c^\rho$ exactly. For the logarithm,
  $\log(cx)/\log(x) = 1 + \log(c)/\log(x) \to 1$ since $\log x \to \infty$.
\end{proof}

\begin{lemma}[Closure properties]
  \label{lem:rv-closure}
  \uses{def:regularly-varying, def:slowly-varying, lem:slowly-varying-iff-rv-zero}
  \lean{ExtremeValueTheory.IsRegularlyVarying.mul, ExtremeValueTheory.IsRegularlyVarying.inv,
    ExtremeValueTheory.IsRegularlyVarying.div, ExtremeValueTheory.IsRegularlyVarying.rpow,
    ExtremeValueTheory.IsRegularlyVarying.const_mul,
    ExtremeValueTheory.IsSlowlyVarying.mul, ExtremeValueTheory.IsSlowlyVarying.inv,
    ExtremeValueTheory.IsSlowlyVarying.div, ExtremeValueTheory.IsSlowlyVarying.rpow}
  \leanok
  Let $f$ and $g$ be regularly varying with indices $\rho$ and $\sigma$, respectively,
  and let $r \in \bR$, $a > 0$. Then
  $f g$ is regularly varying with index $\rho + \sigma$,
  $1/f$ with index $-\rho$,
  $f/g$ with index $\rho - \sigma$,
  $f^r$ with index $\rho r$, and
  $a f$ with index $\rho$.
  In particular, slowly varying functions are closed under products, reciprocals,
  quotients and real powers.
\end{lemma}
\begin{proof}
  \leanok
  Eventual positivity is preserved by all these operations. The limits follow from
  the algebra of limits, e.g.\
  $\frac{f(cx) g(cx)}{f(x) g(x)} = \frac{f(cx)}{f(x)} \cdot \frac{g(cx)}{g(x)} \to c^\rho c^\sigma
  = c^{\rho+\sigma}$, and for real powers continuity of $t \mapsto t^r$ at $c^\rho > 0$.
  The slowly varying statements are the case of index $0$.
\end{proof}

\begin{theorem}[Representation through a slowly varying function]
  \label{thm:rv-iff-exists-sv}
  \uses{def:regularly-varying, def:slowly-varying, lem:rv-examples, lem:rv-closure, lem:rv-congr}
  \lean{ExtremeValueTheory.isRegularlyVarying_iff_exists_isSlowlyVarying}
  \leanok
  A function $f$ is regularly varying with index $\rho$ if and only if there
  is a slowly varying function $L$ such that $f(x) = x^\rho L(x)$ for all
  sufficiently large $x$ [BGT, Theorem~1.4.1].
\end{theorem}
\begin{proof}
  \leanok
  If $f$ is regularly varying with index $\rho$, take $L(x) = x^{-\rho} f(x)$,
  which is regularly varying with index $-\rho + \rho = 0$ by
  Lemma~\ref{lem:rv-closure}. Conversely $x \mapsto x^\rho L(x)$ is regularly
  varying with index $\rho + 0$, and regular variation depends only on the germ
  at $+\infty$ (Lemma~\ref{lem:rv-congr}).
\end{proof}


\section{The uniform convergence theorem: monotone case}

For eventually monotone functions the uniform convergence theorem has an elementary
proof by a sandwich argument, and no measurability is needed.

\begin{theorem}[Uniform convergence theorem, monotone case]
  \label{thm:rv-uct-monotone}
  \uses{def:regularly-varying}
  \lean{ExtremeValueTheory.IsRegularlyVarying.tendstoLocallyUniformlyOn,
    ExtremeValueTheory.IsRegularlyVarying.tendstoUniformlyOn_of_isCompact,
    ExtremeValueTheory.IsRegularlyVarying.tendstoUniformlyOn_Icc}
  \leanok
  Let $f$ be regularly varying with index $\rho$, and monotone or antitone on some
  half-line $[a, \infty)$. Then $f(cx)/f(x) \to c^\rho$ locally uniformly in
  $c \in (0, \infty)$; equivalently, uniformly in $c$ on every compact subset
  $K \subseteq (0,\infty)$ [BGT, Theorem~1.5.2].
\end{theorem}
\begin{proof}
  \leanok
  Fix $c_0 > 0$ and $\eps > 0$. By continuity of $c \mapsto c^\rho$ at $c_0$ choose
  $\eta \in (0, c_0)$ such that $|c^\rho - c_0^\rho| < \eps/3$ for $|c - c_0| \le \eta$.
  By monotonicity, for $x$ large and $c \in [c_0 - \eta, c_0 + \eta]$, the ratio
  $f(cx)/f(x)$ lies between the two ratios $f((c_0 \pm \eta)x)/f(x)$. Each of these
  is within $\eps/3$ of $(c_0 \pm \eta)^\rho$ for large $x$, hence within $2\eps/3$ of
  $c_0^\rho$, and so $f(cx)/f(x)$ is within $2\eps/3$ of $c_0^\rho$ and within $\eps$ of
  $c^\rho$, uniformly for $c$ in the neighbourhood. Locally uniform convergence on
  $(0,\infty)$ upgrades to uniform convergence on compact subsets.
\end{proof}

\begin{theorem}[Potter bounds, monotone case]
  \label{thm:rv-potter-monotone}
  \uses{def:regularly-varying, thm:rv-uct-monotone}
  \lean{ExtremeValueTheory.IsRegularlyVarying.potter_bound}
  \leanok
  Let $f$ be regularly varying with index $\rho$, and monotone or antitone on some
  half-line $[a,\infty)$. Then for every $\delta > 0$ and $A > 1$ there is $x_0 > 0$
  such that $f > 0$ on $[x_0, \infty)$ and
  \begin{align*}
    \frac{f(y)}{f(x)} \; \le \; A \max\Big\{ \Big(\frac{y}{x}\Big)^{\rho+\delta},
      \Big(\frac{y}{x}\Big)^{\rho-\delta} \Big\}
    \qquad \text{for all } x, y \ge x_0
  \end{align*}
  [BGT, Theorem~1.5.6].
\end{theorem}
\begin{proof}
  \leanok
  Choose $\eta > 0$ with $1 + \eta \le A$ and $A^{-1} \le 1 - \eta$. By the uniform
  convergence theorem on $[1,2]$ (Theorem~\ref{thm:rv-uct-monotone}), for $x$ large and
  $t \in [1,2]$ we have
  $(1-\eta) t^{\rho - \delta} \le f(tx)/f(x) \le (1+\eta) t^{\rho+\delta}$.
  Moreover $2^{\rho - \delta} < f(2x)/f(x) < 2^{\rho+\delta}$ for $x$ large.
  By induction on $n$, writing $f(tx)/f(x) = \frac{f(\frac{t}{2} \cdot 2x)}{f(2x)} \cdot \frac{f(2x)}{f(x)}$,
  the two-sided bound extends to all $t \in [1, 2^n]$, hence to all $t \ge 1$.
  For $y \ge x$ apply the upper bound with $t = y/x$; for $y < x$ invert the lower bound
  with $t = x/y$.
\end{proof}


\section{The uniform convergence theorem: measurable case}

Without monotonicity, the uniform convergence theorem requires measurability. The proof
works additively, with $h(t) = \log f(e^t) - \rho t$.

\begin{theorem}[Uniform convergence theorem, additive form]
  \label{thm:uct-additive}
  \lean{ExtremeValueTheory.tendstoUniformlyOn_sub_of_measurable}
  \leanok
  Let $h \colon \bR \to \bR$ be measurable, and assume that $h(x+u) - h(x) \to 0$ as
  $x \to \infty$ for every $u \in \bR$. Then this convergence is uniform in $u$ on
  every compact set $K \subseteq \bR$.
\end{theorem}
\begin{proof}
  \leanok
  It suffices to treat $K = [-N, N]$, and by chaining translates it suffices to treat
  $u \in [0,1]$. Fix $\eps > 0$ and let
  $A_x = \set{ u \in [0,2] \; \big| \; |h(x+u) - h(x)| \ge \eps/2 }$. These are measurable
  sets whose indicators tend to $0$ pointwise, so by dominated convergence their Lebesgue
  measures tend to $0$. For $x$ large, $\Leb(A_x) < 1/2$ and $\Leb(A_{x+\xi}) < 1/2$ for all
  $\xi \in [0,1]$. Since $[1,2]$ has measure $1$, some $u \in [1,2]$ lies neither in $A_x$ nor
  in $A_{x+\xi} + \xi$, and then
  $|h(x+\xi) - h(x)| \le |h(x+u) - h(x)| + |h(x+u) - h(x+\xi)| < \eps$.
\end{proof}

\begin{lemma}[Additive form of a regularly varying function]
  \label{lem:rv-log-comp-exp}
  \uses{def:regularly-varying}
  \lean{ExtremeValueTheory.IsRegularlyVarying.tendsto_log_comp_exp_sub}
  \leanok
  If $f$ is regularly varying with index $\rho$ and $h(t) = \log f(e^t) - \rho t$,
  then $h(t+u) - h(t) \to 0$ as $t \to \infty$ for every $u \in \bR$.
\end{lemma}
\begin{proof}
  \leanok
  We have $h(t+u) - h(t) = \log \frac{f(e^u e^t)}{f(e^t)} - \rho u \to \log(e^{u\rho}) - \rho u = 0$,
  using continuity of $\log$ at $e^{u\rho} > 0$.
\end{proof}

\begin{theorem}[Uniform convergence theorem, measurable case]
  \label{thm:rv-uct-measurable}
  \uses{thm:uct-additive, lem:rv-log-comp-exp, lem:slowly-varying-iff-rv-zero}
  \lean{ExtremeValueTheory.IsRegularlyVarying.tendstoUniformlyOn_of_measurable,
    ExtremeValueTheory.IsRegularlyVarying.tendstoLocallyUniformlyOn_of_measurable,
    ExtremeValueTheory.IsSlowlyVarying.tendstoUniformlyOn_of_measurable,
    ExtremeValueTheory.IsSlowlyVarying.tendstoLocallyUniformlyOn_of_measurable}
  \leanok
  Let $f$ be measurable and regularly varying with index $\rho$. Then
  $f(cx)/f(x) \to c^\rho$ uniformly in $c$ on every compact subset $K \subseteq (0,\infty)$
  (and hence locally uniformly on $(0,\infty)$). In particular, for a measurable slowly
  varying $L$, $L(cx)/L(x) \to 1$ uniformly on compact subsets of $(0,\infty)$
  [BGT, Theorems~1.2.1 and~1.5.2].
\end{theorem}
\begin{proof}
  \leanok
  Apply Theorem~\ref{thm:uct-additive} to $h(t) = \log f(e^t) - \rho t$, which is
  measurable and satisfies its hypothesis by Lemma~\ref{lem:rv-log-comp-exp}. Uniform
  convergence $h(t + u) - h(t) \to 0$ for $u \in \log K$ translates, via $x = e^t$, $c = e^u$
  and uniform continuity of $\exp$ on bounded sets, into uniform convergence of
  $f(cx)/f(x) = c^\rho \exp(h(t+u) - h(t))$ to $c^\rho$.
\end{proof}

\begin{theorem}[Potter bounds, measurable case]
  \label{thm:rv-potter-measurable}
  \uses{thm:rv-uct-measurable}
  \lean{ExtremeValueTheory.IsRegularlyVarying.potter_bound_of_measurable,
    ExtremeValueTheory.IsRegularlyVarying.potter_bounds_of_measurable,
    ExtremeValueTheory.IsSlowlyVarying.potter_bound_of_measurable}
  \leanok
  Let $f$ be measurable and regularly varying with index $\rho$. For every $\delta > 0$ and
  $A > 1$ there is $x_0 > 0$ such that $f > 0$ on $[x_0, \infty)$ and, for all $x, y \ge x_0$,
  \begin{align*}
    A^{-1} \min\Big\{ \Big(\frac{y}{x}\Big)^{\rho+\delta}, \Big(\frac{y}{x}\Big)^{\rho-\delta} \Big\}
    \; \le \; \frac{f(y)}{f(x)} \; \le \;
    A \max\Big\{ \Big(\frac{y}{x}\Big)^{\rho+\delta}, \Big(\frac{y}{x}\Big)^{\rho-\delta} \Big\} .
  \end{align*}
  For a measurable slowly varying $L$ this reads
  $L(y)/L(x) \le A \max\{ (y/x)^\delta, (y/x)^{-\delta} \}$
  [BGT, Theorem~1.5.6].
\end{theorem}
\begin{proof}
  \leanok
  Identical to the proof of Theorem~\ref{thm:rv-potter-monotone}, using
  Theorem~\ref{thm:rv-uct-measurable} on $[1,2]$ instead of the monotone uniform
  convergence theorem. The lower bound follows from the upper bound with the roles
  of $x$ and $y$ exchanged.
\end{proof}


\section{The Karamata representation theorem}

\begin{lemma}[Representation implies slow variation]
  \label{lem:sv-of-representation}
  \uses{def:slowly-varying}
  \lean{ExtremeValueTheory.IsSlowlyVarying.of_representation}
  \leanok
  Let $a \in \bR$, $c, \eps \colon \bR \to \bR$ with $c(x) \to c_0 > 0$ and $\eps(x) \to 0$,
  such that $t \mapsto \eps(t)/t$ is integrable on $[a, x]$ for every $x \ge a$. If
  \begin{align*}
    L(x) = c(x) \exp\Big( \int_a^x \frac{\eps(t)}{t} \, \ud t \Big) \qquad \text{for } x \ge a ,
  \end{align*}
  then $L$ is slowly varying. (No measurability of $L$ is needed.)
\end{lemma}
\begin{proof}
  \leanok
  Eventual positivity is clear. For $\lambda > 0$,
  $\frac{L(\lambda x)}{L(x)} = \frac{c(\lambda x)}{c(x)} \exp \big( \int_x^{\lambda x} \frac{\eps(t)}{t} \, \ud t \big)$.
  The first factor tends to $c_0 / c_0 = 1$. If $|\eps| \le \eta$ on $[\min(x,\lambda x), \infty)$,
  then $\big| \int_x^{\lambda x} \frac{\eps(t)}{t} \, \ud t \big| \le \eta \, |\log \lambda|$, so the
  integral tends to $0$ and the exponential to $1$.
\end{proof}

\begin{theorem}[Karamata representation theorem]
  \label{thm:karamata-representation}
  \uses{lem:sv-of-representation, thm:rv-uct-measurable}
  \lean{ExtremeValueTheory.IsSlowlyVarying.exists_representation,
    ExtremeValueTheory.isSlowlyVarying_iff_exists_representation}
  \leanok
  A measurable function $L$ is slowly varying if and only if there are $a > 0$ and measurable
  functions $c, \eps$ with $c(x) \to c_0 \in (0,\infty)$ and $\eps(x) \to 0$ such that
  \begin{align*}
    L(x) = c(x) \exp\Big( \int_a^x \frac{\eps(t)}{t} \, \ud t \Big) \qquad \text{for all } x \ge a
  \end{align*}
  [BGT, Theorem~1.3.1].
\end{theorem}
\begin{proof}
  \leanok
  The direction ($\Leftarrow$) is Lemma~\ref{lem:sv-of-representation}.
  For ($\Rightarrow$), write $\ell = \log \circ L$ and $g(t) = \ell(t)/t$. By the uniform
  convergence theorem (Theorem~\ref{thm:rv-uct-measurable}), $\ell(t) - \ell(x) \to 0$ uniformly
  for $t \in [x, e x]$; in particular $\ell$ is locally bounded on some $[a, \infty)$, so $g$ is
  locally integrable there. Put $\eps(t) = \ell(e t) - \ell(t)$ for $t \ge a$; then
  $\eps(t) \to 0$, and the substitution $t \mapsto e t$ gives
  $\int_a^x \frac{\eps(t)}{t} \, \ud t = \int_x^{e x} g(t) \, \ud t - \int_a^{e a} g(t) \, \ud t$.
  Since $\int_x^{ex} t^{-1} \, \ud t = 1$, we get
  $\ell(x) - \int_x^{ex} g(t) \, \ud t = \int_x^{ex} \frac{\ell(x) - \ell(t)}{t} \, \ud t \to 0$.
  Hence $c(x) = L(x) \exp\big(-\int_a^x \frac{\eps(t)}{t} \, \ud t\big) \to \exp\big(\int_a^{ea} g(t)\, \ud t\big) > 0$.
\end{proof}

\section{Extended and second-order regular variation}

\begin{definition}[Extended regular variation]
  \label{def:extended-rv}
  \uses{def:regularly-varying}
  \lean{ExtremeValueTheory.extRVLimit, ExtremeValueTheory.IsExtendedRegularlyVarying}
  \leanok
  $f$ is extended regularly varying with index $\gamma \in \bR$ and auxiliary function $a > 0$ if
  $(f(tx) - f(t))/a(t) \to (x^\gamma - 1)/\gamma$ ($\log x$ for $\gamma = 0$) for every $x > 0$
  [dHF, Definition~B.2.1], [BGT, Ch.~3].
\end{definition}

\begin{theorem}[Properties of extended regular variation]
  \label{thm:extended-rv}
  \uses{def:extended-rv}
  \lean{ExtremeValueTheory.IsExtendedRegularlyVarying.isRegularlyVarying,
    ExtremeValueTheory.IsExtendedRegularlyVarying.isRegularlyVarying_of_pos,
    ExtremeValueTheory.IsExtendedRegularlyVarying.exists_tendsto_of_neg,
    ExtremeValueTheory.IsExtendedRegularlyVarying.isRegularlyVarying_sub_of_neg, ExtremeValueTheory.IsExtendedRegularlyVarying.log}
  \leanok
  The auxiliary function is regularly varying with index $\gamma$. For monotone $f$: if $\gamma > 0$
  then $f \in RV_\gamma$ and $a/f \to \gamma$; if $\gamma < 0$ then $f(\infty)$ exists and
  $f(\infty) - f \in RV_\gamma$. If $f > 0$, $\log f$ is extended regularly varying with index
  $\min(\gamma, 0)$ and auxiliary function $a/f$ [dHF, Lemma~3.5.1].
\end{theorem}
\begin{proof}
  \leanok
  Standard manipulations of the limit relation, with Potter-type bounds for the uniformity.
\end{proof}

\begin{definition}[Second-order regular variation]
  \label{def:second-order-rv}
  \uses{def:regularly-varying}
  \lean{ExtremeValueTheory.boxCox, ExtremeValueTheory.IsSecondOrderRegularlyVarying}
  \leanok
  $U$ is second-order regularly varying with indices $\gamma > 0$, $\rho \le 0$ and auxiliary function
  $A \to 0$ of constant sign with $|A| \in RV_\rho$ if
  $(U(tx)/U(t) - x^\gamma)/A(t) \to x^\gamma (x^\rho - 1)/\rho$ for every $x > 0$ [dHF, \S2.3].
\end{definition}

\begin{theorem}[Uniform second-order bound]
  \label{thm:second-order-bound}
  \uses{def:second-order-rv}
  \lean{ExtremeValueTheory.IsSecondOrderRegularlyVarying.tendsto_log, ExtremeValueTheory.IsSecondOrderRegularlyVarying.of_tendsto_div,
    ExtremeValueTheory.IsSecondOrderRegularlyVarying.exists_measurable, ExtremeValueTheory.IsSecondOrderRegularlyVarying.exists_bound}
  \leanok
  For measurable $U$ and $A$ and every $\delta > 0$ there are $C$ and $t_0$ with
  $|\log U(xt) - \log U(t) - \gamma \log x| \le C |A(t)| x^\delta$ for $t \ge t_0$, $x \ge 1$
  [dHF, Theorem~2.3.9] (one-sided version).
\end{theorem}
\begin{proof}
  \leanok
  Uniform convergence on compacts plus a telescoping argument over dyadic scales.
\end{proof}

\section{Multivariate regular variation}

\begin{definition}[Regularly varying random vector]
  \label{def:rv-vector}
  \uses{def:regularly-varying}
  \lean{ExtremeValueTheory.VanishesNearZero, ExtremeValueTheory.IsRegVarVector, ExtremeValueTheory.IsRegVarVector.tendsto_measureReal,
    ExtremeValueTheory.IsRegVarVector.unique}
  \leanok
  A random vector $X$ in a normed space is regularly varying with index $\alpha$ and exponent measure
  $\mu$ if $\PR[\|X\| > x] \in RV_{-\alpha}$, $\mu(\{0\}) = 0$, $\mu$ is finite on
  $\{\|y\| > r\}$ for $r > 0$, and $\E f(X/x) / \PR[\|X\| > x] \to \int f \ud\mu$ for every bounded
  continuous $f$ vanishing near $0$ [Hult--Lindskog 2006], [KS20, \S2.1]. Equivalently,
  $\PR[X/x \in A]/\PR[\|X\| > x] \to \mu(A)$ for $\mu$-continuity sets $A$ bounded away from $0$.
\end{definition}

\begin{theorem}[Homogeneity and polar decomposition]
  \label{thm:rv-vector-polar}
  \uses{def:rv-vector, thm:basrak-segers}
  \lean{ExtremeValueTheory.IsRegVarVector.map_smul, ExtremeValueTheory.IsRegVarVector.measure_smul_set,
    ExtremeValueTheory.IsRegVarVector.measure_lt_norm, ExtremeValueTheory.isRegVarVector_iff,
    ExtremeValueTheory.IsRegVarVector.tendstoLawCond_spectralMeasure, ExtremeValueTheory.isRegVarVector_polarMeasure,
    ExtremeValueTheory.IsRegVarVector.eq_polarMeasure, ExtremeValueTheory.polarMeasure, ExtremeValueTheory.radialMeasure}
  \leanok
  $\mu(tA) = t^{-\alpha}\mu(A)$. $X$ is regularly varying if and only if $\|X\| \in RV_{-\alpha}$ and the
  law of $X/\|X\|$ given $\|X\| > x$ converges weakly to a spectral measure $S$ on the unit sphere;
  then $\mu$ is the image of $\alpha r^{-\alpha-1} \ud r \otimes S$ under $(r, \theta) \mapsto r\theta$
  [Resnick 2007, Theorem~6.1], [BDM 2002].
\end{theorem}
\begin{proof}
  \leanok
  The conditional-limit machinery of Theorem~\ref{thm:basrak-segers}, applied to $X$ itself.
\end{proof}

\begin{theorem}[Linear combinations]
  \label{thm:rv-vector-linear}
  \uses{thm:rv-vector-polar}
  \lean{ExtremeValueTheory.IsRegVarVector.tendsto_linear, ExtremeValueTheory.IsRegVarVector.tendsto_inner,
    ExtremeValueTheory.IsRegVarVector.tendsto_linear_spectral}
  \leanok
  For continuous linear $L$, $\PR[L X > x]/\PR[\|X\| > x] \to \mu\{L > 1\} = \int (L\theta)_+^\alpha \ud S$
  [BDM 2002]. The converse Cram\'er--Wold direction is not formalized.
\end{theorem}
\begin{proof}
  \leanok
  $\{L > 1\}$ is a $\mu$-continuity set bounded away from $0$.
\end{proof}

\begin{proposition}[Asymptotic independence and full dependence]
  \label{prop:rv-vector-examples}
  \uses{thm:rv-vector-polar}
  \lean{ExtremeValueTheory.IsRegVarVector.measure_compl_eq_zero, ExtremeValueTheory.isRegVarVector_diag, ExtremeValueTheory.isRegVarVector_smul_const,
    ExtremeValueTheory.IsRegVarVector.measure_ne_zero_ne_zero_eq_zero, ExtremeValueTheory.isRegVarVector_of_indepFun,
    ExtremeValueTheory.axesSpectralMeasure}
  \leanok
  $(Z, Z)$ has spectral measure $\delta_{(1,1)}$. Independent components put no exponent mass off the
  axes, and for independent nonnegative $X_1, X_2$ with $\PR[X_2 > x]/\PR[X_1 > x] \to c$ the spectral
  measure is $(1+c)^{-1}(\delta_{(1,0)} + c\,\delta_{(0,1)})$ [Resnick 2007, \S6.5].
\end{proposition}
\begin{proof}
  \leanok
  Joint exceedances of independent components are of order $\PR[\|X\| > x]^2$.
\end{proof}


\chapter{Karamata's theorem}

Let $F \colon \bR \to \bR$ be antitone and regularly varying with index $-\alpha$, where
$\alpha > 1$. Typically $F = 1 - G$ is the survival function of a c.d.f.\ $G$ in the
Fréchet domain of attraction (Chapter~\ref{chap:domains-of-attraction}).

\begin{lemma}[Positivity of antitone eventually positive functions]
  \label{lem:pos-of-antitone}
  \lean{ExtremeValueTheory.IsRegularlyVarying.pos_of_antitone}
  \leanok
  An antitone function $F$ with $F(x) > 0$ for all sufficiently large $x$ is positive everywhere.
\end{lemma}
\begin{proof}
  \leanok
  For any $t$ there is $x \ge t$ with $F(x) > 0$, and $F(t) \ge F(x)$.
\end{proof}

\begin{lemma}[Integrability of regularly varying tails]
  \label{lem:rv-integrableOn-Ioi}
  \uses{def:regularly-varying, thm:rv-potter-monotone}
  \lean{ExtremeValueTheory.IsRegularlyVarying.integrableOn_Ioi}
  \leanok
  If $F$ is antitone and regularly varying with index $-\alpha$, $\alpha > 1$, then $F$ is
  integrable on every half-line $(x, \infty)$.
\end{lemma}
\begin{proof}
  \leanok
  By the Potter bounds (Theorem~\ref{thm:rv-potter-monotone}) with
  $\delta = (\alpha - 1)/2$, $F(t) \le C t^{-\alpha+\delta}$ for $t$ large, and
  $-\alpha + \delta < -1$, so $F$ is integrable near $+\infty$. On bounded intervals
  $F$ is bounded and measurable, being monotone.
\end{proof}

\begin{theorem}[Karamata's theorem for tail integrals]
  \label{thm:karamata}
  \uses{def:regularly-varying, thm:rv-potter-monotone, lem:rv-integrableOn-Ioi, lem:pos-of-antitone}
  \lean{ExtremeValueTheory.IsRegularlyVarying.tendsto_integral_Ioi_div}
  \leanok
  Let $F$ be antitone and regularly varying with index $-\alpha$, where $\alpha > 1$. Then
  \begin{align*}
    \lim_{x \to \infty} \frac{\int_x^\infty F(t) \, \ud t}{x \, F(x)} \; = \; \frac{1}{\alpha - 1}
  \end{align*}
  [BGT, Proposition~1.5.10; see also Theorem~1.5.11]. In particular, if $F$ is the survival
  function of a random variable, the mean excess over a high threshold $x$ is asymptotically
  $x/(\alpha - 1)$.
\end{theorem}
\begin{proof}
  \leanok
  After the substitution $t = x u$,
  \begin{align*}
    \frac{\int_x^\infty F(t) \, \ud t}{x \, F(x)} \; = \; \int_1^\infty \frac{F(x u)}{F(x)} \, \ud u .
  \end{align*}
  The integrand tends to $u^{-\alpha}$ pointwise. By the Potter bounds
  (Theorem~\ref{thm:rv-potter-monotone}) with $\delta = (\alpha-1)/2$ and $A = 2$, for $x$
  large it is dominated by $2 u^{-\alpha + \delta}$, which is integrable on $(1, \infty)$.
  Dominated convergence yields the limit $\int_1^\infty u^{-\alpha} \, \ud u = \frac{1}{\alpha - 1}$.
\end{proof}

\begin{theorem}[Karamata's theorem, relaxed hypotheses]
  \label{thm:karamata-relaxed}
  \uses{thm:karamata}
  \lean{ExtremeValueTheory.IsRegularlyVarying.tendsto_integral_Ioi_div_of_measurable,
    ExtremeValueTheory.IsRegularlyVarying.tendsto_integral_Ioi_div_of_antitoneOn}
  \leanok
  The conclusion of Theorem~\ref{thm:karamata} holds for every measurable $F \in RV_{-\alpha}$,
  $\alpha > 1$, and for every $F \in RV_{-\alpha}$ that is antitone on some $[a, \infty)$
  [BGT, Proposition~1.5.10].
\end{theorem}
\begin{proof}
  \leanok
  Potter bounds for measurable regularly varying functions give the domination needed in the
  proof of Theorem~\ref{thm:karamata}.
\end{proof}

\begin{theorem}[Karamata's theorem, converse half]
  \label{thm:karamata-converse}
  \uses{thm:karamata-representation}
  \lean{ExtremeValueTheory.IsRegularlyVarying.of_tendsto_mul_div_integral_Ioi, ExtremeValueTheory.IsRegularlyVarying.of_log_sub_eq_integral,
    ExtremeValueTheory.IsRegularlyVarying.of_tendsto_mul_deriv_div, ExtremeValueTheory.IsRegularlyVarying.of_abs_log_sub_le, ExtremeValueTheory.IsRegularlyVarying.of_tendsto}
  \leanok
  Let $F$ be antitone and integrable on $(a, \infty)$, and $\alpha > 1$. If
  $x F(x) / \int_x^\infty F \to \alpha - 1$, then $F \in RV_{-\alpha}$ [BGT, Theorem~1.6.1].
  Similarly, if $g > 0$ is differentiable with $x g'(x)/g(x) \to \rho$, then $g \in RV_\rho$.
\end{theorem}
\begin{proof}
  \leanok
  With $I(x) = \int_x^\infty F$, $(\log I)'(x) = -F(x)/I(x)$, so $x (\log I)'(x) \to 1 - \alpha$.
  The representation theorem gives $I \in RV_{1-\alpha}$, and $F = (I/x) \cdot (x F/I)$.
\end{proof}

\begin{theorem}[Karamata's theorem for integrals over $[0, x]$]
  \label{thm:karamata-zero}
  \uses{thm:karamata-relaxed}
  \lean{ExtremeValueTheory.IsRegularlyVarying.tendsto_integral_zero_div}
  \leanok
  If $f \in RV_\rho$ is measurable and locally integrable on $[0, \infty)$ and $\rho > -1$, then
  $\int_0^x f / (x f(x)) \to 1/(\rho + 1)$ [BGT, Proposition~1.5.8].
\end{theorem}
\begin{proof}
  \leanok
  Uniform convergence on compacts and Potter bounds near $0$ after rescaling $t = x s$.
\end{proof}


\chapter{Breiman's lemma}

Breiman's lemma (Breiman 1965; see also Jessen and Mikosch, \emph{Regularly varying
functions}, 2006, Lemma~4.2(1)) describes the tail of a product $X Y$ of independent random
variables, where $X$ has a regularly varying tail and $Y \ge 0$ has a slightly higher moment.

\begin{lemma}[Analytic core of Breiman's lemma]
  \label{lem:breiman-analytic}
  \uses{def:regularly-varying, thm:rv-potter-monotone}
  \lean{ExtremeValueTheory.IsRegularlyVarying.tendsto_integral_div}
  \leanok
  Let $F$ be antitone and regularly varying with index $-\alpha$, $\alpha > 0$. Let $\mu$ be a
  finite measure on $\bR$ concentrated on $[0, \infty)$ such that $y \mapsto y^{\alpha + \eps}$
  is $\mu$-integrable for some $\eps > 0$. Then
  \begin{align*}
    \lim_{x \to \infty} \frac{1}{F(x)} \int \indof{\set{y > 0}}(y) \, F(x / y) \, \ud \mu(y)
    \; = \; \int y^\alpha \, \ud \mu(y) .
  \end{align*}
\end{lemma}
\begin{proof}
  \leanok
  For $y > 0$ the integrand divided by $F(x)$ tends to $y^\alpha$ by regular variation
  (with $c = 1/y$), and for $y = 0$ both sides vanish. For domination, choose
  $0 < \delta \le \eps$ with $\delta < \alpha$, and let $x_0$ be given by the Potter bounds
  (Theorem~\ref{thm:rv-potter-monotone}) with $A = 2$. Then:
  \begin{itemize}
    \item if $x, x/y \ge x_0$, then
      $F(x/y)/F(x) \le 2 \max\{ y^{\alpha - \delta}, y^{\alpha + \delta} \}$;
    \item for $x \ge x_0$, $F(x) \ge \frac{F(x_0)}{2} (x / x_0)^{-(\alpha+\delta)}$, so that when
      $x / y < x_0$, $F(x/y)/F(x) \le 2 \frac{F(0)}{F(x_0)} y^{\alpha + \delta}$.
  \end{itemize}
  Both bounds are at most a constant times $1 + y^{\alpha + \eps}$, which is $\mu$-integrable
  since $\mu$ is finite. Dominated convergence concludes.
\end{proof}

\begin{theorem}[Breiman's lemma]
  \label{thm:breiman}
  \uses{lem:breiman-analytic}
  \lean{ExtremeValueTheory.tendsto_measure_mul_div_measure}
  \leanok
  Let $X$ and $Y$ be independent real random variables on a probability space, such that
  $x \mapsto \PR[X > x]$ is regularly varying with index $-\alpha$, $\alpha > 0$, and
  $Y \ge 0$ almost surely with $\EX[Y^{\alpha + \eps}] < \infty$ for some $\eps > 0$. Then
  \begin{align*}
    \lim_{x \to \infty} \frac{\PR[X Y > x]}{\PR[X > x]} \; = \; \EX[Y^\alpha] .
  \end{align*}
  No sign condition on $X$ is needed.
\end{theorem}
\begin{proof}
  \leanok
  By independence the joint law of $(X, Y)$ is the product of the marginals, so
  conditioning on $Y$ gives, for $x > 0$,
  $\PR[X Y > x] = \int \indof{\set{y > 0}}(y) \, \PR[X > x / y] \, \ud \PR_Y(y)$
  (on $\set{Y = 0}$ the product is $0 < x$). The survival function $x \mapsto \PR[X > x]$
  is antitone, so Lemma~\ref{lem:breiman-analytic} applied to $\mu = \PR_Y$ gives the claim.
\end{proof}


\chapter{Domains of attraction}
\label{chap:domains-of-attraction}

A c.d.f.\ $F$ is in the \emph{maximum domain of attraction} of an extreme value
distribution $G$ if there are orientation preserving affine isomorphisms $A_n \in \OriAffR$
with $\big( (A_n . F)(x) \big)^n \to G(x)$ for every $x$, i.e., if the suitably normalized
maximum of $n$ i.i.d.\ samples from $F$ converges in distribution to $G$. This chapter proves
Gnedenko's characterizations of the Fréchet and Weibull domains of attraction in terms of
regular variation. References: de Haan and Ferreira, \emph{Extreme Value Theory: An
Introduction}, Springer (2006) [dHF]; Resnick, \emph{Extreme Values, Regular Variation, and
Point Processes}, Springer (1987) [Resnick]; Embrechts, Klüppelberg and Mikosch,
\emph{Modelling Extremal Events}, Springer (1997) [EKM].

\section{The Fréchet domain of attraction}

\begin{definition}[Fréchet maximum domain of attraction]
  \label{def:mda-frechet}
  \uses{def:cdf, def:std-Frechet-cdf, lem:oriented-affine-action-on-cdf}
  \lean{CumulativeDistributionFunction.IsInMDAFrechet}
  \leanok
  Let $\alpha > 0$. A c.d.f.\ $F$ is in the maximum domain of attraction of the standard
  Fréchet distribution $\FrechetCDF{\alpha}$ (with $\FrechetCDF{\alpha}(x) = \exp(-x^{-\alpha})$
  for $x > 0$ and $0$ for $x \le 0$) if there are $A_n \in \OriAffR$ with
  $\big( (A_n . F)(x) \big)^n \to \FrechetCDF{\alpha}(x)$ for every $x \in \bR$.
  (Since $\FrechetCDF{\alpha}$ is continuous, this is the same as convergence at all
  continuity points.)
\end{definition}

\begin{lemma}[Scaling constants from regular variation]
  \label{lem:frechet-exists-scaling}
  \uses{def:regularly-varying}
  \lean{ExtremeValueTheory.Frechet.exists_tendsto_mul_of_isRegularlyVarying}
  \leanok
  Let $g$ be antitone, right-continuous, with $g(x) \to 0$ as $x \to \infty$, and regularly
  varying with index $-\alpha$. Then there are constants $a_n > 0$ such that
  $n \, g(a_n x) \to x^{-\alpha}$ for every $x > 0$.
\end{lemma}
\begin{proof}
  \leanok
  Take the quantile-type constants $a_n = \inf \set{ x \; | \; g(x) \le 1/n }$ (made positive).
  Right continuity gives $n \, g(a_n) \le 1$, and regular variation at $c \, a_n$ with $c < 1$
  gives $\liminf n \, g(a_n) \ge c^{-\alpha}$ for every $c<1$; hence $n \, g(a_n) \to 1$ (using
  $a_n \to \infty$, which holds as $g > 0$ eventually). Then
  $n \, g(a_n x) = n \, g(a_n) \cdot \frac{g(x a_n)}{g(a_n)} \to x^{-\alpha}$.
\end{proof}

\begin{lemma}[Regular variation from affine normalization]
  \label{lem:frechet-rv-of-affine}
  \uses{def:regularly-varying}
  \lean{ExtremeValueTheory.Frechet.isRegularlyVarying_of_tendsto_mul_affine}
  \leanok
  Let $g$ be antitone, $\alpha > 0$, and $a_n > 0$, $b_n \in \bR$ with
  $n \, g(a_n y + b_n) \to y^{-\alpha}$ for every $y > 0$. Then $g$ is regularly varying with
  index $-\alpha$.
\end{lemma}
\begin{proof}
  \leanok
  Extend $a, b$ to real arguments by $a(t) = a_{\lfloor t \rfloor}$, $b(t) = b_{\lfloor t \rfloor}$,
  so that $t \, g(a(t) y + b(t)) \to y^{-\alpha}$. Comparing the levels $t$ and $2t$ by
  monotonicity of $g$ shows that $a(2t) \ge \rho \, a(t)$ eventually for some $\rho > 1$ (indeed
  $a(2t) \sim 2^{1/\alpha} a(t)$) and that $b(2t) - b(t) = o(a(t))$. A doubling argument
  (summing along dyadic chains $t, 2t, 4t, \ldots$ and using local bounds on a base interval
  $[T, 2T]$) then yields $a(t) \to \infty$ and $b(t)/a(t) \to 0$. Hence the centring may be
  dropped, $n \, g(a_n x) \to x^{-\alpha}$ for $x > 0$. Finally, for $t$ large choose
  $N(t) = \sup \set{ n \; | \; a_n \le t }$, so that $a_{N(t)} \le t < a_{N(t)+1}$; a sandwich
  using monotonicity of $g$ gives $g(c t)/g(t) \to c^{-\alpha}$.
\end{proof}

\begin{theorem}[Gnedenko's theorem, Fréchet case]
  \label{thm:mda-frechet-iff}
  \uses{def:mda-frechet, def:regularly-varying, lem:frechet-exists-scaling,
    lem:frechet-rv-of-affine, thm:tfae-ev-limit}
  \lean{CumulativeDistributionFunction.isInMDAFrechet_iff,
    CumulativeDistributionFunction.isInMDAFrechet_of_isRegularlyVarying,
    CumulativeDistributionFunction.isRegularlyVarying_of_isInMDAFrechet}
  \leanok
  Let $\alpha > 0$. A c.d.f.\ $F$ is in the maximum domain of attraction of $\FrechetCDF{\alpha}$
  if and only if its tail $x \mapsto 1 - F(x)$ is regularly varying with index $-\alpha$
  [dHF, Theorem~1.2.1], [Resnick, Proposition~1.11], [EKM, Theorem~3.3.7].
\end{theorem}
\begin{proof}
  \leanok
  Sufficiency: the tail $g = 1 - F$ is antitone, right-continuous and tends to $0$, so
  Lemma~\ref{lem:frechet-exists-scaling} gives $a_n > 0$ with $n (1 - F(a_n x)) \to x^{-\alpha}
  = -\log \FrechetCDF{\alpha}(x)$ for $x > 0$. With $A_n^{-1}(x) = a_n x$, the equivalence of
  the limit relations (Theorem~\ref{thm:tfae-ev-limit}) gives
  $(A_n . F)(x)^n \to \FrechetCDF{\alpha}(x)$ for $x > 0$. For $x \le 0$ we use
  $F(a_n x) \le F(0) < 1$ (as $1 - F > 0$ everywhere), so $F(a_n x)^n \to 0$.

  Necessity: write $A_n^{-1}(x) = a_n x + b_n$ with $a_n > 0$. By
  Theorem~\ref{thm:tfae-ev-limit}, $n (1 - F(a_n y + b_n)) \to y^{-\alpha}$ for $y > 0$, and
  Lemma~\ref{lem:frechet-rv-of-affine} gives regular variation of $1 - F$.
\end{proof}


\section{The Weibull domain of attraction}

\begin{definition}[Weibull maximum domain of attraction]
  \label{def:mda-weibull}
  \uses{def:cdf, def:std-Weibull-cdf, lem:oriented-affine-action-on-cdf}
  \lean{CumulativeDistributionFunction.IsInMDAWeibull}
  \leanok
  Let $\alpha > 0$. A c.d.f.\ $F$ is in the maximum domain of attraction of the standard
  Weibull distribution $\WeibullCDF{\alpha}$ (with $\WeibullCDF{\alpha}(x) = \exp(-(-x)^\alpha)$
  for $x < 0$ and $1$ for $x \ge 0$) if there are $A_n \in \OriAffR$ with
  $\big( (A_n . F)(x) \big)^n \to \WeibullCDF{\alpha}(x)$ for every $x \in \bR$.
\end{definition}

\begin{lemma}[Regularly varying antitone functions vanish at infinity]
  \label{lem:weibull-tendsto-zero}
  \uses{def:regularly-varying}
  \lean{ExtremeValueTheory.Weibull.tendsto_zero_of_isRegularlyVarying}
  \leanok
  If $g$ is antitone, positive, and regularly varying with index $-\alpha$, $\alpha > 0$,
  then $g(x) \to 0$ as $x \to \infty$.
\end{lemma}
\begin{proof}
  \leanok
  $g$ has a limit $\ell \ge 0$ at $+\infty$ by monotonicity. If $\ell > 0$, then
  $g(2x)/g(x) \to \ell/\ell = 1 \ne 2^{-\alpha}$, a contradiction.
\end{proof}

\begin{lemma}[The endpoint lemma]
  \label{lem:weibull-endpoint}
  \lean{ExtremeValueTheory.Weibull.eventually_eq_zero_of_tendsto_mul_affine}
  \leanok
  Let $u$ be antitone and nonnegative, $\alpha > 0$, $a_n > 0$, $b_n \in \bR$, with
  $n \, u(a_n z + b_n) \to (-z)^\alpha$ for every $z < 0$ and $n \, u(b_n) \to 0$. Then for
  every $c > 0$, eventually $u(b_n + c \, a_n) = 0$.
\end{lemma}
\begin{proof}
  \leanok
  Let $\kappa = 2^{-1/\alpha} < 1$. Comparing the levels $n$ and $2n$ by monotonicity of $u$
  (using $2n \, u(a_{2n} z + b_{2n}) \to 2(-z)^\alpha$) shows that, for every small $\delta > 0$,
  eventually $a_{2n} \le \rho \, a_n$ with $\rho = \kappa + 2\delta < 1$ and
  $b_{2n} - b_n \le 3 \delta \, a_n$. Choosing $\delta$ small relative to $c (1 - \rho)$, the
  potential $\Phi_n = b_n + c \, a_n$ is eventually non-increasing along doublings,
  $\Phi_{2n} \le \Phi_n$. Suppose $\eta = u(\Phi_n) > 0$ for some large $n$. For $m = 2^k n$ we get
  $b_m < \Phi_m \le \Phi_n$, so $u(b_m) \ge \eta$ and $m \, u(b_m) \ge \eta$ for all $k$,
  contradicting $m \, u(b_m) \to 0$. Hence $u(\Phi_n) = 0$ eventually.
\end{proof}

\begin{theorem}[Gnedenko's theorem, Weibull case]
  \label{thm:mda-weibull-iff}
  \uses{def:mda-weibull, def:regularly-varying, lem:frechet-exists-scaling,
    lem:frechet-rv-of-affine, lem:weibull-tendsto-zero, lem:weibull-endpoint, thm:tfae-ev-limit}
  \lean{CumulativeDistributionFunction.isInMDAWeibull_iff,
    CumulativeDistributionFunction.isInMDAWeibull_of_isRegularlyVarying,
    CumulativeDistributionFunction.exists_endpoint_of_isInMDAWeibull}
  \leanok
  Let $\alpha > 0$. A c.d.f.\ $F$ is in the maximum domain of attraction of $\WeibullCDF{\alpha}$
  if and only if $F$ has a finite right endpoint $x_F$ (that is, $F(x_F) = 1$ and $F(x) < 1$ for
  $x < x_F$) and $t \mapsto 1 - F(x_F - t^{-1})$ is regularly varying with index $-\alpha$
  [dHF, Theorem~1.2.1(2)], [EKM, Theorem~3.3.12].
\end{theorem}
\begin{proof}
  \leanok
  The transformation $x \mapsto x_F - x^{-1}$ reduces both directions to the analytic lemmas
  of the Fréchet case, applied to the antitone function $g(t) = 1 - F(x_F - t^{-1})$
  (suitably extended to small $t$).

  Sufficiency: $g$ tends to $0$ by Lemma~\ref{lem:weibull-tendsto-zero}, and
  Lemma~\ref{lem:frechet-exists-scaling} gives $a_n$ with $n \, g(a_n y) \to y^{-\alpha}$ for
  $y > 0$. Taking $A_n^{-1}(x) = x_F + a_n^{-1} x$ and substituting $y = -x^{-1}$ for $x < 0$
  gives $n (1 - F(A_n^{-1}(x))) \to (-x)^\alpha$, hence the limit relation by
  Theorem~\ref{thm:tfae-ev-limit}; for $x \ge 0$, $F(A_n^{-1}(x)) \ge F(x_F) = 1$.

  Necessity: write $A_n^{-1}(x) = a_n x + b_n$. Theorem~\ref{thm:tfae-ev-limit} gives
  $n (1 - F(a_n z + b_n)) \to (-z)^\alpha$ for $z < 0$, while $F(b_n)^n \to 1$ gives
  $n (1 - F(b_n)) \to 0$. The endpoint lemma (Lemma~\ref{lem:weibull-endpoint}) shows that
  $F(b_n + c \, a_n) = 1$ eventually for every $c > 0$; in particular the right endpoint
  $x_F$ is finite, and one deduces $(b_n - x_F)/a_n \to 0$. The centring $b_n$ can then be
  replaced by $x_F$, and Lemma~\ref{lem:frechet-rv-of-affine}, applied after the
  transformation $x \mapsto x_F - x^{-1}$, gives regular variation of $g$.
\end{proof}


\section{The Gumbel domain of attraction}

% TODO: The Gumbel domain of attraction (ExtremeValueTheory/DomainOfAttraction/Gumbel.lean)
% is being formalized separately. Add its blueprint entries here (auxiliary function
% characterization, von Mises condition, examples) once the Lean file is finished.

The Gumbel domain of attraction (sufficient conditions via an auxiliary function and the
von Mises condition) is work in progress.


\chapter{Peaks over threshold}

The Pickands--Balkema--de Haan theorem (Balkema and de Haan, 1974; Pickands, 1975) says
that, for distributions in a maximum domain of attraction, the conditional distribution of
the excess over a high threshold $u$ is, after rescaling, close to a generalized Pareto
distribution. We prove the heavy-tailed (Fréchet) case; see Embrechts, Klüppelberg and
Mikosch, \emph{Modelling Extremal Events} (1997), Theorem~3.4.13 [EKM], and McNeil, Frey
and Embrechts, \emph{Quantitative Risk Management} (2005), \S5.2 [MFE].

Throughout, $F \colon \bR \to \bR$ plays the role of a survival function
$F(x) = \PR[X > x]$.

\begin{definition}[Generalized Pareto survival function]
  \label{def:gpd-survival}
  \lean{ExtremeValueTheory.gpdSurvival}
  \leanok
  For shape $\xi \ge 0$, scale $\beta > 0$ and $x \ge 0$, the survival function of the
  generalized Pareto distribution (GPD) is
  \begin{align*}
    \overline{G}_{\xi,\beta}(x) \; = \;
    \begin{cases}
      (1 + \xi x / \beta)^{-1/\xi} & \text{if } \xi \ne 0, \\
      \exp(-x/\beta) & \text{if } \xi = 0 .
    \end{cases}
  \end{align*}
\end{definition}

\begin{lemma}[Basic properties of the GPD]
  \label{lem:gpd-properties}
  \uses{def:gpd-survival}
  \lean{ExtremeValueTheory.gpdSurvival_of_eq_zero, ExtremeValueTheory.gpdSurvival_of_ne_zero,
    ExtremeValueTheory.gpdSurvival_pos, ExtremeValueTheory.gpdSurvival_le_one,
    ExtremeValueTheory.gpdSurvival_antitoneOn, ExtremeValueTheory.gpdSurvival_tendsto_atTop,
    ExtremeValueTheory.gpdSurvival_continuousOn, ExtremeValueTheory.gpdSurvival_inv_one}
  \leanok
  Let $\xi \ge 0$ and $\beta > 0$. On $[0, \infty)$ the function $\overline{G}_{\xi,\beta}$ takes
  values in $(0, 1]$, is antitone and continuous, and tends to $0$ at $+\infty$. Moreover, for
  $\alpha \ne 0$, $\overline{G}_{1/\alpha, 1}(x) = (1 + x/\alpha)^{-\alpha}$.
\end{lemma}
\begin{proof}
  \leanok
  For $x \ge 0$ the base $1 + \xi x / \beta$ is at least $1$, and the exponent $-1/\xi$ is negative,
  so all properties follow from elementary properties of real powers and of $\exp$.
\end{proof}

\begin{definition}[Excess survival function and mean excess function]
  \label{def:excess-survival}
  \lean{ExtremeValueTheory.excessSurvival, ExtremeValueTheory.meanExcess}
  \leanok
  The \emph{conditional excess survival function} over the threshold $u$ is
  $\overline{F}_u(x) = F(u + x)/F(u)$; if $F(x) = \PR[X > x]$, this is $\PR[X - u > x \mid X > u]$.
  The \emph{mean excess function} is $e(u) = \big( \int_u^\infty F(t) \, \ud t \big) / F(u)$, which
  equals $\EX[X - u \mid X > u]$ when $\EX[X^+] < \infty$.
\end{definition}

\begin{lemma}[Pointwise convergence of excesses]
  \label{lem:pot-pointwise}
  \uses{def:regularly-varying, def:gpd-survival, def:excess-survival, lem:gpd-properties}
  \lean{ExtremeValueTheory.IsRegularlyVarying.tendsto_excessSurvival}
  \leanok
  If $F$ is regularly varying with index $-\alpha$, $\alpha > 0$, then with
  $\beta(u) = u/\alpha$, for each $x \ge 0$,
  \begin{align*}
    \lim_{u \to \infty} \frac{F(u + \beta(u) x)}{F(u)} \; = \; \overline{G}_{1/\alpha, 1}(x) .
  \end{align*}
\end{lemma}
\begin{proof}
  \leanok
  We have $u + \beta(u) x = (1 + x/\alpha) u$, so by regular variation with $c = 1 + x/\alpha > 0$
  the ratio tends to $(1 + x/\alpha)^{-\alpha} = \overline{G}_{1/\alpha,1}(x)$
  (Lemma~\ref{lem:gpd-properties}).
\end{proof}

\begin{theorem}[Pickands--Balkema--de Haan, Fréchet case]
  \label{thm:pickands-balkema-de-haan}
  \uses{lem:pot-pointwise, thm:rv-uct-monotone, lem:gpd-properties}
  \lean{ExtremeValueTheory.IsRegularlyVarying.tendstoUniformlyOn_excessSurvival,
    ExtremeValueTheory.IsRegularlyVarying.tendsto_sup_excessSurvival}
  \leanok
  Let $F$ be antitone and regularly varying with index $-\alpha$, $\alpha > 0$, and let
  $\beta(u) = u/\alpha$. Then $x \mapsto F(u + \beta(u) x)/F(u)$ converges to
  $\overline{G}_{1/\alpha, 1}$ uniformly on $[0, \infty)$ as $u \to \infty$; that is,
  \begin{align*}
    \lim_{u \to \infty} \; \sup_{x \ge 0} \; \Big| \frac{F(u + \beta(u) x)}{F(u)}
      - \overline{G}_{1/\alpha, 1}(x) \Big| \; = \; 0
  \end{align*}
  [EKM, Theorem~3.4.13].
\end{theorem}
\begin{proof}
  \leanok
  Let $\eps > 0$. Since $\overline{G}_{1/\alpha,1}(x) \to 0$, choose $M$ with
  $\overline{G}_{1/\alpha,1}(M) < \eps/2$. On the compact interval $c = 1 + x/\alpha \in [1, 1 + M/\alpha]$
  the uniform convergence theorem for monotone regularly varying functions
  (Theorem~\ref{thm:rv-uct-monotone}) gives uniform convergence for $x \in [0, M]$. For $x > M$,
  both functions are nonnegative and bounded by their values at $x = M$ (by monotonicity), which
  are eventually less than $\eps$; hence the difference is less than $\eps$ there as well.
  The supremum form is a reformulation of uniform convergence.
\end{proof}

\begin{theorem}[Mean excess function of a regularly varying tail]
  \label{thm:mean-excess}
  \uses{def:excess-survival, thm:karamata}
  \lean{ExtremeValueTheory.IsRegularlyVarying.tendsto_meanExcess_div}
  \leanok
  If $F$ is antitone and regularly varying with index $-\alpha$, $\alpha > 1$, then
  $e(u)/u \to 1/(\alpha - 1)$ as $u \to \infty$. This matches the mean
  $\beta(u)/(1 - \xi) = u/(\alpha - 1)$ of the limiting GPD with $\xi = 1/\alpha$ and
  $\beta(u) = u/\alpha$ [EKM, \S3.4], [MFE, \S5.2].
\end{theorem}
\begin{proof}
  \leanok
  $e(u)/u = \big( \int_u^\infty F(t) \, \ud t \big) / \big( u \, F(u) \big)$, so this is exactly
  Karamata's theorem (Theorem~\ref{thm:karamata}).
\end{proof}


\chapter{Order statistics and the Hill estimator}
\label{chap:hill}

This chapter develops the order statistics of a finite sample and proves the weak consistency of
the Hill estimator of the extreme value index $\gamma = 1/\alpha$ of a regularly varying tail. We
follow de Haan and Ferreira, \emph{Extreme Value Theory: An Introduction}, Springer (2006), \S3.2
[dHF], and Mason, \emph{Laws of large numbers for sums of extreme values}, Ann.\ Probab.\ 10
(1982) [Mason].

\section{Order statistics}

\begin{definition}[Descending order statistics]
  \label{def:order-stat}
  \lean{ExtremeValueTheory.orderStat, ExtremeValueTheory.orderStat_of_lt,
    ExtremeValueTheory.orderStat_of_le}
  \leanok
  For a sample $x = (x_0, \ldots, x_{n-1}) \in \bR^n$, the descending order statistics
  $X_{(1)} \ge X_{(2)} \ge \cdots \ge X_{(n)}$ are the entries of $x$ sorted decreasingly. In Lean
  they are indexed from $0$: $\texttt{orderStat}\,x\,j = X_{(j+1)}$ for $j < n$ (and $0$, a junk
  value, for $j \ge n$). They are obtained by reversing the increasing sort of the tuple $x$.
\end{definition}

\begin{definition}[Exceedance count]
  \label{def:exceed-count}
  \lean{ExtremeValueTheory.exceedCount, ExtremeValueTheory.exceedCount_eq_sum}
  \leanok
  The number of exceedances of the level $a$ by the sample $x$ is
  $N_a(x) = \# \set{ i \; | \; x_i > a } = \sum_i \indof{\set{x_i > a}}$.
\end{definition}

\begin{lemma}[Basic properties of order statistics]
  \label{lem:order-stat-basic}
  \uses{def:order-stat}
  \lean{ExtremeValueTheory.orderStat_perm, ExtremeValueTheory.sum_orderStat,
    ExtremeValueTheory.orderStat_anti}
  \leanok
  The order statistics $(X_{(1)}, \ldots, X_{(n)})$ are a permutation of the sample; in particular
  $\sum_{j=1}^n f(X_{(j)}) = \sum_i f(x_i)$ for every $f$. They are decreasing:
  $X_{(j+1)} \le X_{(i+1)}$ for $i \le j < n$.
\end{lemma}
\begin{proof}
  \leanok
  The sort is a permutation, and so is its composition with the reversal $\texttt{Fin.rev}$;
  monotonicity is the monotonicity of the sorted tuple, reversed.
\end{proof}

\begin{lemma}[Duality between order statistics and exceedance counts]
  \label{lem:order-stat-duality}
  \uses{def:order-stat, def:exceed-count, lem:order-stat-basic}
  \lean{ExtremeValueTheory.lt_orderStat_iff, ExtremeValueTheory.orderStat_le_iff}
  \leanok
  For $j < n$ and $a \in \bR$: $a < X_{(j+1)}$ if and only if $N_a(x) \ge j + 1$; equivalently
  $X_{(j+1)} \le a$ if and only if $N_a(x) \le j$.
\end{lemma}
\begin{proof}
  \leanok
  If $a < X_{(j+1)}$, then $X_{(1)}, \ldots, X_{(j+1)}$ all exceed $a$, and these are $j + 1$
  distinct entries of the sample. Conversely, if $X_{(j+1)} \le a$, then only
  $X_{(1)}, \ldots, X_{(j)}$ can exceed $a$, so $N_a(x) \le j$.
\end{proof}

\begin{lemma}[Measurability]
  \label{lem:order-stat-measurable}
  \uses{def:order-stat, def:exceed-count}
  \lean{ExtremeValueTheory.measurable_exceedCount, ExtremeValueTheory.measurable_orderStat}
  \leanok
  If $X_0, \ldots, X_{n-1}$ are measurable random variables, then the exceedance counts and the
  order statistics of $(X_0, \ldots, X_{n-1})$ are measurable.
\end{lemma}
\begin{proof}
  \leanok
  $N_a$ is a finite sum of indicators of measurable sets. For the order statistics, by
  Lemma~\ref{lem:order-stat-duality} the sets $\set{ a < X_{(j+1)} }$ are measurable for all $a$.
\end{proof}

\begin{theorem}[Distribution of exceedance counts and order statistics]
  \label{thm:exceed-count-binomial}
  \uses{lem:order-stat-duality, lem:order-stat-measurable}
  \lean{ExtremeValueTheory.measure_lt_orderStat_eq, ExtremeValueTheory.measure_exceedCount_eq,
    ExtremeValueTheory.measure_lt_orderStat_eq_sum}
  \leanok
  For random variables $X_0, \ldots, X_{n-1}$ and $j < n$,
  $\PR[X_{(j+1)} > a] = \PR[N_a \ge j + 1]$. If moreover the $X_i$ are independent with common
  exceedance probability $p = \PR[X_i > a]$, then $N_a$ is binomial,
  \begin{align*}
    \PR[N_a = m] \; = \; \binom{n}{m} p^m (1 - p)^{n - m} ,
    \qquad
    \PR[X_{(j+1)} > a] \; = \; \sum_{m = j+1}^{n} \binom{n}{m} p^m (1 - p)^{n - m} .
  \end{align*}
\end{theorem}
\begin{proof}
  \leanok
  The first identity is Lemma~\ref{lem:order-stat-duality} pointwise. The event $\set{N_a = m}$ is
  the disjoint union over $m$-element subsets $S$ of indices of the events
  $\bigcap_{i \in S} \set{X_i > a} \cap \bigcap_{i \notin S} \set{X_i \le a}$, each of probability
  $p^m (1 - p)^{n-m}$ by independence. Summing over $m \ge j + 1$ gives the last formula.
\end{proof}

\section{The Hill estimator}

\begin{definition}[Hill estimator]
  \label{def:hill}
  \uses{def:order-stat}
  \lean{ExtremeValueTheory.hill}
  \leanok
  The Hill estimator based on the top $k + 1$ order statistics of a sample $x \in \bR^n$ is
  \begin{align*}
    H_{k,n} \; = \; \frac{1}{k} \sum_{i=1}^{k} \big( \log X_{(i)} - \log X_{(k+1)} \big) .
  \end{align*}
\end{definition}

\begin{definition}[Log-excess]
  \label{def:log-excess}
  \lean{ExtremeValueTheory.logExcess, ExtremeValueTheory.logExcess_nonneg,
    ExtremeValueTheory.logExcess_of_le, ExtremeValueTheory.logExcess_of_ge,
    ExtremeValueTheory.logExcess_anti, ExtremeValueTheory.measurable_logExcess}
  \leanok
  The log-excess of $y$ over the threshold $t > 0$ is
  $\log^+(y/t) = \log\big( \max(y, t)/t \big)$. It is nonnegative and measurable in $y$, equals $0$
  for $y \le t$ and $\log y - \log t$ for $y \ge t$, and is antitone in $t$.
\end{definition}

\begin{lemma}[Log-excess representation of the Hill estimator]
  \label{lem:hill-log-excess}
  \uses{def:hill, def:log-excess, lem:order-stat-basic}
  \lean{ExtremeValueTheory.hill_eq_sum_logExcess}
  \leanok
  If $k < n$ and $X_{(k+1)} > 0$, then
  $H_{k,n} = \frac{1}{k} \sum_{i=0}^{n-1} \log^+ \big( x_i / X_{(k+1)} \big)$.
\end{lemma}
\begin{proof}
  \leanok
  By Lemma~\ref{lem:order-stat-basic}, the sum over the sample equals the sum over the order
  statistics. The terms with $j \ge k + 1$ vanish since $X_{(j)} \le X_{(k+1)}$, the term $j = k + 1$
  vanishes, and for $j \le k$ the term is $\log X_{(j)} - \log X_{(k+1)}$.
\end{proof}

\begin{lemma}[A weak law of large numbers for triangular arrays]
  \label{lem:chebyshev-wlln}
  \lean{ExtremeValueTheory.tendsto_measure_abs_sum_div_sub}
  \leanok
  Let $X_0, X_1, \ldots$ be i.i.d., $g_n$ measurable with $g_n(X_0) \in L^2$, and $b_n > 0$ with
  $n \, \EX[g_n(X_0)]/b_n \to m$ and $n \, \EX[g_n(X_0)^2]/b_n^2 \to 0$. Then
  \begin{align*}
    \frac{1}{b_n} \sum_{i<n} g_n(X_i) \; \to \; m \qquad \text{in probability.}
  \end{align*}
\end{lemma}
\begin{proof}
  \leanok
  By pairwise independence, the variance of the normalized sum is at most
  $n \, \EX[g_n(X_0)^2]/b_n^2 \to 0$, while its mean tends to $m$. Chebyshev's inequality
  concludes.
\end{proof}

\begin{definition}[Survival function of a random variable]
  \label{def:survival}
  \lean{ExtremeValueTheory.survival, ExtremeValueTheory.survival_anti,
    ExtremeValueTheory.survival_nonneg, ExtremeValueTheory.survival_pos,
    ExtremeValueTheory.survival_eq_one_sub_cdf, ExtremeValueTheory.continuousWithinAt_survival,
    ExtremeValueTheory.tendsto_survival_atTop}
  \leanok
  The survival function of a real random variable $Y$ is $\overline{F}(x) = \PR[Y > x]$. It is
  nonnegative, antitone and right-continuous, equals $1 - F_Y(x)$ for the c.d.f.\ $F_Y$ of $Y$,
  and tends to $0$ at $+\infty$. If it is regularly varying, it is everywhere positive.
\end{definition}

\begin{lemma}[Uniform exponential decay along geometric levels]
  \label{lem:survival-exp-le}
  \uses{def:survival, thm:rv-potter-monotone}
  \lean{ExtremeValueTheory.exists_survival_exp_le}
  \leanok
  If $\overline{F}$ is regularly varying with index $-\alpha$, $\alpha > 0$, there is $x_0 > 0$
  such that $\overline{F}(s e^u) \le 2 \, \overline{F}(s) \, e^{-(\alpha/2) u}$ for all $s \ge x_0$
  and $u \ge 0$.
\end{lemma}
\begin{proof}
  \leanok
  This is the Potter bound (Theorem~\ref{thm:rv-potter-monotone}) with $\delta = \alpha/2$ and
  constant $2$, applied with $c = e^u \ge 1$.
\end{proof}

\begin{lemma}[Moments of the log-excess]
  \label{lem:log-excess-moments}
  \uses{def:log-excess, def:survival, lem:survival-exp-le}
  \lean{ExtremeValueTheory.setOf_lt_logExcess, ExtremeValueTheory.integral_logExcess,
    ExtremeValueTheory.integrableOn_mul_exp_neg, ExtremeValueTheory.lintegral_logExcess_sq_le,
    ExtremeValueTheory.exists_memLp_logExcess}
  \leanok
  For $s, u > 0$, $\set{ \log^+(Y/s) > u } = \set{ Y > s e^u }$, and hence (layer cake formula)
  \begin{align*}
    \EX \big[ \log^+(Y/s) \big] \; = \; \int_0^\infty \overline{F}(s e^u) \, \ud u .
  \end{align*}
  If $\overline{F}$ is regularly varying with index $-\alpha$, $\alpha > 0$, there are $x_0, C$ such
  that for $s \ge x_0$ the log-excess $\log^+(Y/s)$ is square integrable with
  $\EX \big[ \log^+(Y/s)^2 \big] \le C \, \overline{F}(s)$.
\end{lemma}
\begin{proof}
  \leanok
  The layer cake formula gives $\EX[\log^+(Y/s)^2] = \int_0^\infty 2u \, \overline{F}(s e^u) \, \ud u$,
  and Lemma~\ref{lem:survival-exp-le} bounds it by
  $4 \overline{F}(s) \int_0^\infty u \, e^{-(\alpha/2)u} \, \ud u < \infty$.
\end{proof}

\begin{lemma}[Integrated regular variation]
  \label{lem:tail-integral-log-excess}
  \uses{lem:log-excess-moments, lem:survival-exp-le, def:regularly-varying}
  \lean{ExtremeValueTheory.integral_exp_neg_mul_Ioi, ExtremeValueTheory.tendsto_integral_logExcess_div}
  \leanok
  If $\overline{F}$ is regularly varying with index $-\alpha$, $\alpha > 0$, then
  \begin{align*}
    \lim_{s \to \infty} \frac{\EX \big[ \log^+(Y/s) \big]}{\overline{F}(s)} \; = \;
    \lim_{s \to \infty} \int_0^\infty \frac{\overline{F}(s e^u)}{\overline{F}(s)} \, \ud u
    \; = \; \int_0^\infty e^{-\alpha u} \, \ud u \; = \; \frac{1}{\alpha} .
  \end{align*}
\end{lemma}
\begin{proof}
  \leanok
  Pointwise $\overline{F}(s e^u)/\overline{F}(s) \to e^{-\alpha u}$ by regular variation, and the
  integrands are dominated by $2 e^{-(\alpha/2) u}$ (Lemma~\ref{lem:survival-exp-le}), so dominated
  convergence applies.
\end{proof}

\begin{lemma}[Intermediate thresholds]
  \label{lem:hill-threshold}
  \lean{ExtremeValueTheory.upperQuantile_spec_of_tendsto_atTop}
  \leanok
  Let $\overline{F}$ be antitone and right-continuous with $\overline{F}(x) \to 0$ as
  $x \to \infty$, and let $0 < p < \overline{F}(M)$. Then $u = \inf \set{ x \; | \; \overline{F}(x) \le p }$
  satisfies $\overline{F}(u) \le p$, $p < \overline{F}(x)$ for $x < u$, and $M \le u$.
\end{lemma}
\begin{proof}
  \leanok
  As in Lemma~\ref{lem:gumbel-quantile}.
\end{proof}

\begin{theorem}[Weak consistency of the Hill estimator]
  \label{thm:hill-consistency}
  \uses{lem:hill-log-excess, lem:chebyshev-wlln, lem:log-excess-moments,
    lem:tail-integral-log-excess, lem:hill-threshold, thm:exceed-count-binomial,
    lem:order-stat-duality, def:regularly-varying}
  \lean{ExtremeValueTheory.tendstoInMeasure_hill}
  \leanok
  Let $X_0, X_1, \ldots$ be i.i.d.\ real random variables whose survival function
  $\overline{F}(x) = \PR[X_0 > x]$ is regularly varying with index $-\alpha$, $\alpha > 0$, and let
  $k = k(n)$ be an intermediate sequence: $k(n) \to \infty$ and $k(n)/n \to 0$. Then the Hill
  estimator $H_{k(n),n}$ of $X_0, \ldots, X_{n-1}$ converges in probability to $1/\alpha$
  [dHF, Theorem~3.2.2], [Mason].
\end{theorem}
\begin{proof}
  \leanok
  Instead of the Rényi representation of exponential spacings used in [dHF], we use the tail
  empirical process argument in the spirit of [Mason], which only needs Chebyshev's inequality.

  \emph{Step 1 (deterministic).} By Lemma~\ref{lem:hill-log-excess},
  $H_{k,n} = \frac{1}{k} \sum_i \log^+(X_i / X_{(k+1)})$, and $s \mapsto \sum_i \log^+(X_i/s)$ is
  antitone (Definition~\ref{def:log-excess}).

  \emph{Step 2 (thresholds).} Choose deterministic thresholds
  $u_n = \inf \set{ x \; | \; \overline{F}(x) \le k/n }$. By Lemma~\ref{lem:hill-threshold},
  right-continuity and regular variation, $u_n \to \infty$, $n \overline{F}(u_n)/k \to 1$, and
  hence $n \overline{F}(c u_n)/k \to c^{-\alpha}$ for every $c > 0$.

  \emph{Step 3 (laws of large numbers).} Fix $c > 0$. By Lemma~\ref{lem:chebyshev-wlln} with
  $b_n = k(n)$:
  the exceedance count satisfies $N_{c u_n}/k \to c^{-\alpha}$ in probability (its summands are
  indicators with mean $\overline{F}(c u_n)$, and the second moment condition holds since
  $k \to \infty$); and $\frac{1}{k} \sum_i \log^+(X_i/(c u_n)) \to c^{-\alpha}/\alpha$ in
  probability, the mean being computed by Lemma~\ref{lem:tail-integral-log-excess} and the second
  moment being $O(\overline{F}(c u_n)) = O(k/n)$ by Lemma~\ref{lem:log-excess-moments}.

  \emph{Step 4 (sandwich).} By the duality
  $X_{(k+1)} > a \iff N_a \ge k + 1$ (Lemma~\ref{lem:order-stat-duality}), Step 3 gives
  $c_1 u_n < X_{(k+1)} \le c_2 u_n$ with probability tending to $1$ for any $c_1 < 1 < c_2$. On
  this event, monotonicity (Step 1) sandwiches $H_{k,n}$ between the averages of log-excesses
  over $c_2 u_n$ and over $c_1 u_n$, which are close to $c_2^{-\alpha}/\alpha$ and
  $c_1^{-\alpha}/\alpha$. Letting $c_1, c_2 \to 1$ concludes.
\end{proof}

\begin{corollary}[Hill estimator for Pareto samples]
  \label{cor:hill-pareto}
  \uses{thm:hill-consistency, lem:rv-examples}
  \lean{ExtremeValueTheory.tendstoInMeasure_hill_pareto}
  \leanok
  If $X_0, X_1, \ldots$ are i.i.d.\ Pareto$(\alpha)$, i.e.\ $\PR[X_0 > x] = x^{-\alpha}$ for
  $x \ge 1$, and $k(n) \to \infty$, $k(n)/n \to 0$, then $H_{k(n),n} \to 1/\alpha$ in probability.
\end{corollary}
\begin{proof}
  \leanok
  The survival function agrees with $x^{-\alpha}$ for $x \ge 1$, hence is regularly varying with
  index $-\alpha$; apply Theorem~\ref{thm:hill-consistency}.
\end{proof}

\section{The Pickands and moment estimators}

\begin{definition}[Pickands estimator]
  \label{def:pickands}
  \uses{def:order-stat}
  \lean{ExtremeValueTheory.pickands}
  \leanok
  $\hat\xi_{k,n} = (\log 2)^{-1} \log \dfrac{X_{(k+1)} - X_{(2k+1)}}{X_{(2k+1)} - X_{(4k+1)}}$
  [dHF, \S3.3], [EKM, \S6.4.2].
\end{definition}

\begin{theorem}[Consistency of the Pickands estimator]
  \label{thm:pickands-consistency}
  \uses{def:pickands, thm:pbdh-general, thm:pbdh-iff}
  \lean{ExtremeValueTheory.tendstoInMeasure_pickands, ExtremeValueTheory.tendstoInMeasure_pickands_exponential,
    ExtremeValueTheory.gpdSurvivalInv_pickands_ratio, ExtremeValueTheory.tendsto_measure_abs_exceedCount_div_sub,
    ExtremeValueTheory.tendsto_measure_not_mem_orderStat}
  \leanok
  For iid samples with c.d.f.\ $F$ in the domain of attraction of $H_\xi$, $\xi \in \bR$
  arbitrary, and $k \to \infty$, $k/n \to 0$, the Pickands estimator converges to $\xi$ in
  probability [EKM, Theorem~6.4.1], [dHF, \S3.3].
\end{theorem}
\begin{proof}
  \leanok
  Let $u_n$ be the $(1 - 8k/n)$-quantile and $\beta$ the scale from
  Pickands--Balkema--de Haan. The number of exceedances of $u_n + \beta(u_n) x$, divided by $k$,
  converges in probability to $8 \bar G_{\xi}(x)$, which localizes
  $(X_{(jk+1)} - u_n)/\beta(u_n)$ near $\bar G_\xi^{-1}(j/8)$ for $j = 1, 2, 4$. The ratio of
  differences of these three points is exactly $2^\xi$.
\end{proof}

\begin{definition}[Log-moments and the moment estimator]
  \label{def:moment-estimator}
  \uses{def:hill}
  \lean{ExtremeValueTheory.logMoment, ExtremeValueTheory.logMoment_one, ExtremeValueTheory.momentEstimator}
  \leanok
  $M^{(j)}_{k,n} = k^{-1} \sum_{i<k} (\log X_{(i+1)} - \log X_{(k+1)})^j$, so $M^{(1)}$ is the
  Hill estimator, and the moment estimator is
  $\hat\gamma = M^{(1)} + 1 - \tfrac12 \big(1 - (M^{(1)})^2/M^{(2)}\big)^{-1}$ [dHF, \S3.5].
\end{definition}

\begin{theorem}[Consistency of the log-moments and the moment estimator, $\xi > 0$]
  \label{thm:moment-consistency}
  \uses{def:moment-estimator, thm:hill-consistency}
  \lean{ExtremeValueTheory.tendstoInMeasure_logMoment, ExtremeValueTheory.tendstoInMeasure_momentEstimator,
    ExtremeValueTheory.tendstoInMeasure_momentEstimator_of_isInMDAGEV, ExtremeValueTheory.tendstoInMeasure_hill_of_isInMDAGEV,
    ExtremeValueTheory.tendstoInMeasure_momentEstimator_pareto, ExtremeValueTheory.tendsto_integral_logExcess_rpow_div}
  \leanok
  Under the hypotheses of Theorem~\ref{thm:hill-consistency},
  $M^{(j)}_{k,n} \to j!/\alpha^j$ in probability for $j \ge 1$, and the moment estimator converges
  to $1/\alpha$. Equivalently, for $F$ in the domain of attraction of $H_\xi$ with $\xi > 0$, both
  the Hill and moment estimators are consistent for $\xi$. The case $\xi \le 0$ is not
  formalized.
\end{theorem}
\begin{proof}
  \leanok
  As for the Hill estimator, with $\E[\log^+(Y/s)^p] / \bar F(s) \to \Gamma(p+1)/\alpha^p$ in place
  of the first log-moment.
\end{proof}

\section{Extreme quantiles and tail probabilities}

\begin{definition}[Weissman and tail-probability estimators]
  \label{def:weissman}
  \uses{def:hill, def:order-stat}
  \lean{ExtremeValueTheory.upperQuantile, ExtremeValueTheory.weissman, ExtremeValueTheory.tailProbEstimator, ExtremeValueTheory.IsBoundedInProbability}
  \leanok
  For $p_n \to 0$, the quantile $x_{p} = \inf\{x : \bar F(x) \le p\}$ is estimated by
  $\hat x_{p_n} = X_{(k+1)} \, (k/(n p_n))^{\hat\gamma}$, and $\bar F(y_n)$ by
  $\hat p(y_n) = (k/n)(y_n / X_{(k+1)})^{-1/\hat\gamma}$, with $\hat\gamma$ the Hill estimator
  [Weissman 1978], [dHF, \S4.3], [EKM, \S6.4.1].
\end{definition}

\begin{theorem}[Ratio consistency of the Weissman estimator]
  \label{thm:weissman}
  \uses{def:weissman, thm:hill-consistency}
  \lean{ExtremeValueTheory.tendstoInMeasure_weissman_div_of_bias, ExtremeValueTheory.tendstoInMeasure_weissman_div,
    ExtremeValueTheory.tendstoInMeasure_weissman_div_of_isBoundedInProbability,
    ExtremeValueTheory.tendstoInMeasure_weissman_div_pareto, ExtremeValueTheory.tendstoInMeasure_orderStat_div_upperQuantile,
    ExtremeValueTheory.tendsto_upperQuantile}
  \leanok
  Let $X_i$ be iid with $\bar F \in RV_{-\alpha}$, $k \to \infty$, $k/n \to 0$ and
  $d_n = k/(n p_n)$. Then $\hat x_{p_n}/x_{p_n} \to 1$ in probability in each of these cases:
  \begin{itemize}
    \item $d_n$ stays in a compact subset of $(0, \infty)$ (intermediate quantiles);
    \item $\sqrt{k}(\hat\gamma - 1/\alpha)$ is bounded in probability, $\log d_n / \sqrt k \to 0$, and the
      bias condition $x_{p_n} / (x_{k/n} d_n^{1/\alpha}) \to 1$ holds, which is automatic for exact
      Pareto tails. This covers $n p_n \to c \in [0, \infty)$, beyond the sample range.
  \end{itemize}
  [dHF, Theorem~4.3.1] (consistency part).
\end{theorem}
\begin{proof}
  \leanok
  Write $\hat x_{p_n}/x_{p_n} = (X_{(k+1)}/x_{k/n}) \cdot d_n^{\hat\gamma - 1/\alpha} \cdot
  x_{k/n} d_n^{1/\alpha} / x_{p_n}$. The first factor tends to $1$ by localization of the order
  statistic, the second because $(\hat\gamma - 1/\alpha) \log d_n \to 0$, and the third is the bias.
\end{proof}

\begin{theorem}[Ratio consistency of the tail-probability estimator]
  \label{thm:tail-prob-estimator}
  \uses{thm:weissman}
  \lean{ExtremeValueTheory.tendstoInMeasure_tailProbEstimator_div_of_bias, ExtremeValueTheory.tendstoInMeasure_tailProbEstimator_div,
    ExtremeValueTheory.tendstoInMeasure_tailProbEstimator_div_of_isBoundedInProbability}
  \leanok
  Under the same three sets of hypotheses (with $e_n = y_n/x_{k/n}$ in place of $d_n$),
  $\hat p(y_n)/\bar F(y_n) \to 1$ in probability [dHF, \S4.4].
\end{theorem}
\begin{proof}
  \leanok
  As for Theorem~\ref{thm:weissman}.
\end{proof}

\begin{definition}[Return levels]
  \label{def:return-level}
  \lean{ExtremeValueTheory.gevFun, ExtremeValueTheory.returnLevel, ExtremeValueTheory.gevFun_returnLevel}
  \leanok
  The $m$-block return level of $H_\xi$ is $z_m = ((-\log(1 - 1/m))^{-\xi} - 1)/\xi$
  ($-\log(-\log(1 - 1/m))$ for $\xi = 0$), the level with $H_\xi(z_m) = 1 - 1/m$ [Coles, \S3.1.3].
\end{definition}

\section{Upper order statistics and records}

\begin{theorem}[Limit law of the $k$-th largest]
  \label{thm:kth-largest}
  \uses{def:order-stat}
  \lean{ExtremeValueTheory.measureReal_orderStat_le_eq_sum, ExtremeValueTheory.tendsto_measureReal_orderStat_le,
    ExtremeValueTheory.tendsto_measureReal_affine_orderStat_le,
    ExtremeValueTheory.exists_tendsto_measureReal_affine_orderStat_le_of_isInMDAGumbel}
  \leanok
  For iid samples with $n(1 - F(u_n)) \to \tau \in [0, \infty)$,
  $\PR[X_{(k)} \le u_n] \to e^{-\tau} \sum_{j < k} \tau^j/j!$. Consequently, if
  $F^n(a_n x + b_n) \to H(x) \in (0, 1)$, then
  $\PR[(X_{(k)} - b_n)/a_n \le x] \to H(x) \sum_{j<k} (-\log H(x))^j/j!$
  [EKM, Theorem~4.2.3], [LLR, Theorem~2.2.2].
\end{theorem}
\begin{proof}
  \leanok
  $\{X_{(k)} \le u_n\}$ is the event that fewer than $k$ observations exceed $u_n$; the
  exceedance count is binomial with mean $\to \tau$, hence asymptotically Poisson.
\end{proof}

\begin{theorem}[Joint law of the two largest]
  \label{thm:two-largest}
  \uses{thm:kth-largest}
  \lean{ExtremeValueTheory.measureReal_orderStat_zero_le_and_orderStat_one_le,
    ExtremeValueTheory.tendsto_measureReal_orderStat_zero_le_and_orderStat_one_le,
    ExtremeValueTheory.tendsto_measureReal_affine_orderStat_zero_one}
  \leanok
  If $v_n \le u_n$, $n(1 - F(u_n)) \to \tau_1$ and $n(1 - F(v_n)) \to \tau_2$, then
  $\PR[X_{(1)} \le u_n, X_{(2)} \le v_n] \to e^{-\tau_2}(1 + \tau_2 - \tau_1)$ [EKM, \S4.2].
\end{theorem}
\begin{proof}
  \leanok
  Exactly: $\PR[X_{(1)} \le u, X_{(2)} \le v] = F(v)^n + n (F(u) - F(v)) F(v)^{n-1}$.
\end{proof}

\begin{theorem}[Records]
  \label{thm:records}
  \lean{ExtremeValueTheory.IsRecord, ExtremeValueTheory.recordCount, ExtremeValueTheory.measureReal_isRecord, ExtremeValueTheory.measureReal_isRecord_and_isRecord,
    ExtremeValueTheory.integral_recordCount, ExtremeValueTheory.integral_sq_recordCount_sub_le,
    ExtremeValueTheory.tendstoInMeasure_recordCount_div_log}
  \leanok
  For an iid sequence with atomless marginal, the $k$-th observation (counting from $1$) is a record
  with probability $1/k$, record indicators are pairwise independent, the number of records
  $N_n$ among the first $n$ observations satisfies $\E N_n = H_n = \sum_{k \le n} 1/k$ and
  $\Var N_n \le H_n$, and $N_n / \log n \to 1$ in probability [EKM, \S5.4], [Resnick, \S4.1].
  Full mutual independence of the indicators (R\'enyi) is not formalized.
\end{theorem}
\begin{proof}
  \leanok
  By exchangeability, the maximum of the first $k$ observations is equally likely to be any of
  them, and there are no ties. A second exchangeability argument gives pairwise independence;
  Chebyshev's inequality concludes.
\end{proof}

\section{Asymptotic normality of the Hill estimator}

\begin{lemma}[R\'enyi representation]
  \label{lem:renyi}
  \uses{def:order-stat}
  \lean{ExtremeValueTheory.lintegral_orderStat_excess, ExtremeValueTheory.map_orderStat_excess}
  \leanok
  For an iid standard exponential sample of size $n$ and $k < n$, the excesses
  $E_{(i)} - E_{(k+1)}$, $i \le k$, are (conditionally on $E_{(k+1)}$) iid standard exponential
  [dHF, Lemma~3.2.3].
\end{lemma}
\begin{proof}
  \leanok
  Count exceedances and split off coordinates, using the memoryless property.
\end{proof}

\begin{theorem}[Asymptotic normality of the Hill estimator]
  \label{thm:hill-normality}
  \uses{lem:renyi, thm:second-order-bound, thm:hill-consistency}
  \lean{ExtremeValueTheory.survivalQuantile, ExtremeValueTheory.tendstoInDistribution_hill_pareto,
    ExtremeValueTheory.tendstoInDistribution_hill_of_tendsto_zero, ExtremeValueTheory.tendstoInDistribution_hill,
    ExtremeValueTheory.tendstoInDistribution_meanExcess, ExtremeValueTheory.map_eq_expMeasure_map_survivalQuantile}
  \leanok
  Let $X_i$ be iid with tail quantile function $U$ second-order regularly varying with indices
  $\gamma > 0$, $\rho \le 0$ and auxiliary function $A$. If $k \to \infty$, $k/n \to 0$ and
  $\sqrt k A(n/k) \to \lambda$, then
  \begin{align*}
    \sqrt k \, (\hat\gamma_{k,n} - \gamma) \; \Rightarrow \; N\big(\lambda/(1 - \rho), \gamma^2\big)
  \end{align*}
  [dHF, Theorem~3.2.5]. For exact Pareto samples, $\sqrt k(\hat\gamma - 1/\alpha) \Rightarrow
  N(0, \alpha^{-2})$ without the condition $k/n \to 0$.
\end{theorem}
\begin{proof}
  \leanok
  Write $X = U(e^E)$ with $E$ standard exponential. On the exponential scale the Hill estimator is
  $\gamma$ times the mean of $k$ iid exponential excesses plus a bias term
  $k^{-1}\sum_i (\log U(e^{E_{(i)}}) - \log U(e^{E_{(k+1)}}) - \gamma(E_{(i)} - E_{(k+1)}))$. The CLT
  handles the first part; the uniform second-order bound and $E_{(k+1)} - \log(n/k) \to 0$ show the
  bias is $A(n/k)/(1-\rho) + o_P(k^{-1/2})$.
\end{proof}

\begin{theorem}[Consistency of the moment estimator for all $\gamma$]
  \label{thm:moment-general}
  \uses{thm:mda-quantile, thm:extended-rv, thm:moment-consistency}
  \lean{ExtremeValueTheory.gpdMoment_one, ExtremeValueTheory.gpdMoment_two, ExtremeValueTheory.tendstoInMeasure_logMoment_of_isInMDAGEV,
    ExtremeValueTheory.tendstoInMeasure_momentEstimator_of_isInMDAGEV'}
  \leanok
  For $F$ in the domain of attraction of $H_\gamma$, $\gamma \in \bR$, with right endpoint $x_F > 0$,
  $k \to \infty$ and $k/n \to 0$: with $\gamma_- = \min(\gamma, 0)$ and $q_n = a(n/k)/U(n/k)$,
  $M^{(1)}/q_n \to 1/(1 - \gamma_-)$ and $M^{(2)}/q_n^2 \to 2/((1-\gamma_-)(1-2\gamma_-))$ in
  probability, and the moment estimator converges to $\gamma$ [dHF, Theorem~3.5.2].
\end{theorem}
\begin{proof}
  \leanok
  $\log U$ is extended regularly varying with index $\gamma_-$; a weak law at the deterministic
  threshold $U(n/k)$ via the layer-cake formula and Potter bounds, then sandwiching the random
  threshold $X_{(k+1)}$ with exceedance counts.
\end{proof}

\begin{theorem}[Weissman estimator under second-order regular variation]
  \label{thm:weissman-second-order}
  \uses{thm:weissman, thm:hill-normality, thm:second-order-bound}
  \lean{MeasureTheory.TendstoInDistribution.exists_measure_le_abs_le,
    ExtremeValueTheory.isBoundedInProbability_of_tendstoInDistribution, ExtremeValueTheory.isBoundedInProbability_hill,
    ExtremeValueTheory.IsSecondOrderRegularlyVarying.exists_bound_of_neg,
    ExtremeValueTheory.survivalQuantile_of_one_lt, ExtremeValueTheory.upperQuantile_eq_survivalQuantile,
    ExtremeValueTheory.IsSecondOrderRegularlyVarying.tendsto_div_mul_rpow, ExtremeValueTheory.tendsto_upperQuantile_div_of_secondOrder,
    ExtremeValueTheory.tendstoInMeasure_weissman_div_of_secondOrder,
    ExtremeValueTheory.tendstoInMeasure_tailProbEstimator_div_of_secondOrder, ExtremeValueTheory.tendstoInDistribution_weissman}
  \leanok
  Assume the tail quantile function is second-order regularly varying with $\gamma > 0$ and
  $\rho < 0$, $k \to \infty$, $k/n \to 0$ and $\sqrt k A(n/k) \to \lambda$. If $d_n = k/(n p_n) \ge 1$ and
  $\log d_n / \sqrt k \to 0$, then $\hat x_{p_n}/x_{p_n} \to 1$ in probability, and similarly for the
  tail-probability estimator. If moreover $d_n \to \infty$,
  \begin{align*}
    \frac{\sqrt k}{\log d_n} \Big( \frac{\hat x_{p_n}}{x_{p_n}} - 1 \Big)
      \; \Rightarrow \; N\big(\lambda/(1-\rho), \gamma^2\big)
  \end{align*}
  [dHF, \S4.3]. The case $\rho = 0$ is not covered.
\end{theorem}
\begin{proof}
  \leanok
  Convergence in distribution implies tightness, so the Hill CLT gives the rate. For $\rho < 0$ the
  uniform bound $|\log U(xt) - \log U(t) - \gamma \log x| \le C|A(t)|$ for all $x \ge 1$ makes the bias
  ratio tend to $1$. For normality, $\log(\hat x/x) = (\hat\gamma - \gamma)\log d_n + O_P(k^{-1/2})$ and the
  delta method for $\exp$.
\end{proof}

\section{The Hill estimator for dependent data}

\begin{theorem}[Hill consistency from a variance bound]
  \label{thm:hill-variance}
  \uses{thm:hill-consistency}
  \lean{ExtremeValueTheory.tailCount, ExtremeValueTheory.tendstoInMeasure_hill_of_variance, ExtremeValueTheory.tendsto_measure_abs_tailCount_div_sub,
    ExtremeValueTheory.logExcess_le_sum, ExtremeValueTheory.sum_le_logExcess}
  \leanok
  Let $X_i$ have identical marginals with $\bar F \in RV_{-\alpha}$, $k \to \infty$, $k/n \to 0$. If
  $\Var(\sum_{i<n} \mathbf 1\{X_i > u_n\}) / k_n^2 \to 0$ along every $u_n \to \infty$ with
  $n \bar F(u_n)/k_n \to c > 0$, then the Hill estimator converges to $1/\alpha$ in probability. No
  independence, stationarity or mixing is assumed.
\end{theorem}
\begin{proof}
  \leanok
  Chebyshev for the exceedance counts. The log-excess sum is squeezed deterministically between Riemann
  sums of exceedance counts on a geometric grid of thresholds; the remainder is controlled by Markov's
  inequality and Potter bounds, which use only the marginal law.
\end{proof}

\begin{theorem}[Hill consistency under m-dependence and strong mixing]
  \label{thm:hill-dependent}
  \uses{thm:hill-variance, def:mixing}
  \lean{ExtremeValueTheory.variance_tailCount_le_of_isMDependent, ExtremeValueTheory.tendstoInMeasure_hill_of_isMDependent,
    ExtremeValueTheory.variance_tailCount_le_of_isStationary, ExtremeValueTheory.tendstoInMeasure_hill_of_alphaMixing,
    ExtremeValueTheory.tendstoInMeasure_hill_of_alphaMixing_rate}
  \leanok
  The Hill estimator is consistent for $m$-dependent sequences with identical marginals (stationarity is
  not needed), and for stationary $\alpha$-mixing sequences with block length $r_n$ such that
  $\sum_{1 \le h \le r_n} \PR[X_h > u_n \mid X_0 > u_n] = o(k_n)$ and $(n/k_n)^2 \alpha(r_n + 1) \to 0$
  [Resnick--St\u{a}ric\u{a} 1995], [Rootz\'en--Leadbetter--de Haan 1998].
\end{theorem}
\begin{proof}
  \leanok
  $\Var \le (2m+1) n \bar F(u)$ under $m$-dependence. Under stationarity,
  $\Var \le n(\bar F(u) + 2\sum_{h \le r} \PR[X_0 > u, X_h > u] + 2 \sum_{h > r} \alpha(h))$.
\end{proof}

\begin{theorem}[Extreme quantiles at $\rho = 0$: bias cancellation]
  \label{thm:weissman-rho-zero}
  \uses{thm:weissman-second-order}
  \lean{ExtremeValueTheory.tendstoInMeasure_weissman_div_of_le, ExtremeValueTheory.tendstoInMeasure_weissman_div_of_antitoneOn,
    ExtremeValueTheory.tendstoInMeasure_tailProbEstimator_div_of_le, ExtremeValueTheory.tendstoInDistribution_weissman_of_tendsto_bias,
    ExtremeValueTheory.tendstoInDistribution_weissman_of_rho_eq_zero,
    ExtremeValueTheory.tendstoInDistribution_tailProbEstimator_of_tendsto_bias, ExtremeValueTheory.tendstoInDistribution_tailProbEstimator,
    ExtremeValueTheory.tendstoInDistribution_tailProbEstimator_of_rho_eq_zero}
  \leanok
  The Weissman and tail-probability results extend to $\rho \le 0$ given a majorant of $|A|$ on the
  extrapolation range (automatic when $|A|$ is eventually monotone). At $\rho = 0$, with $A$ asymptotically
  constant on that range, the Hill bias $+\lambda$ and the extrapolation bias $-\lambda$ cancel:
  $\sqrt k/\log d_n\,(\hat x/x_p - 1) \Rightarrow N(0, \gamma^2)$ and
  $\sqrt k/\log e_n\,(\hat p/p - 1) \Rightarrow N(0, \gamma^{-2})$. For $\rho < 0$ the tail-probability
  estimator satisfies $\sqrt k/\log e_n\,(\hat p/p - 1) \Rightarrow N(\lambda/((1-\rho)\gamma^2), \gamma^{-2})$
  [dHF, \S4.3--4.4].
\end{theorem}
\begin{proof}
  \leanok
  A uniform bound $|\log U(xt) - \log U(t) - \gamma \log x| \le C(1 + \log x)\,a$ when $|A| \le a$ on $[t, xt]$, and
  the delta method. A bound $C|A(t)|\log x$ fails in general for slowly varying $|A|$.
\end{proof}


\chapter{Subexponential distributions}

We follow Embrechts, Klüppelberg and Mikosch, \emph{Modelling Extremal Events for Insurance
and Finance} (Springer, 1997), \S1.3 [EKM], and Foss, Korshunov and Zachary, \emph{An
Introduction to Heavy-Tailed and Subexponential Distributions} (2nd ed., Springer, 2013).

Measures live on $\bR$; a measure $\mu$ is \emph{concentrated on $[0, \infty)$} if
$\mu\big((-\infty, 0)\big) = 0$. Convolution $\mu * \nu$ is the push-forward of
$\mu \otimes \nu$ under addition, i.e., the law of $X_1 + X_2$ for independent
$X_1 \sim \mu$, $X_2 \sim \nu$.

\section{Tails and convolutions}

\begin{definition}[Tail of a measure]
  \label{def:tail}
  \lean{ExtremeValueTheory.tail}
  \leanok
  The \emph{tail} (survival function) of a measure $\mu$ on $\bR$ is
  $\overline{F}_\mu(x) = \mu\big( (x, \infty) \big)$, regarded as a real number.
\end{definition}

\begin{lemma}[Elementary properties of tails]
  \label{lem:tail-basic}
  \uses{def:tail}
  \lean{ExtremeValueTheory.Subexp.tail_nonneg, ExtremeValueTheory.Subexp.tail_le_one,
    ExtremeValueTheory.Subexp.tail_antitone, ExtremeValueTheory.Subexp.tail_measurable,
    ExtremeValueTheory.Subexp.tail_of_neg, ExtremeValueTheory.Subexp.tendsto_tail_atTop,
    ExtremeValueTheory.Subexp.measureReal_Ioc_eq, ExtremeValueTheory.Subexp.measureReal_Iic_eq}
  \leanok
  For a finite measure $\mu$, the tail $\overline{F}_\mu$ is nonnegative, antitone, measurable,
  and tends to $0$ at $+\infty$, and $\mu\big( (a, b] \big) = \overline{F}_\mu(a) - \overline{F}_\mu(b)$
  for $a \le b$. For a probability measure, $\overline{F}_\mu \le 1$,
  $\mu\big( (-\infty, a] \big) = 1 - \overline{F}_\mu(a)$, and if $\mu$ is concentrated on
  $[0,\infty)$ then $\overline{F}_\mu(x) = 1$ for $x < 0$.
\end{lemma}
\begin{proof}
  \leanok
  Monotonicity and continuity from above of measures.
\end{proof}

\begin{lemma}[Convolution formula for tails]
  \label{lem:tail-conv}
  \uses{def:tail}
  \lean{ExtremeValueTheory.Subexp.conv_apply, ExtremeValueTheory.Subexp.tail_conv,
    ExtremeValueTheory.Subexp.conv_Iio_zero, ExtremeValueTheory.Subexp.integrable_tail_sub}
  \leanok
  For finite measures $\mu, \nu$,
  \begin{align*}
    \overline{F}_{\mu * \nu}(x) \; = \; \int \overline{F}_\nu(x - t) \, \ud \mu(t) .
  \end{align*}
  If $\mu$ and $\nu$ are concentrated on $[0, \infty)$, then so is $\mu * \nu$.
\end{lemma}
\begin{proof}
  \leanok
  By Tonelli, $(\mu * \nu)(S) = \int \nu(S - t) \, \ud \mu(t)$ for measurable $S$; take
  $S = (x, \infty)$. If $s + t < 0$ then $s < 0$ or $t < 0$, which gives the second claim.
\end{proof}

\begin{lemma}[Splitting the convolution integral]
  \label{lem:tail-conv-split}
  \uses{lem:tail-conv, lem:tail-basic}
  \lean{ExtremeValueTheory.Subexp.tail_conv_split, ExtremeValueTheory.Subexp.setIntegral_tail_sub_le,
    ExtremeValueTheory.Subexp.setIntegral_tail_sub_nonneg, ExtremeValueTheory.Subexp.tail_conv_self_ge}
  \leanok
  Let $\mu, \nu$ be probability measures with $\nu$ concentrated on $[0, \infty)$, and $a \le x$.
  Then
  \begin{align*}
    \overline{F}_{\mu * \nu}(x) \; = \; \int_{(-\infty, a]} \overline{F}_\nu(x - t) \, \ud \mu(t)
      + \int_{(a, x]} \overline{F}_\nu(x - t) \, \ud \mu(t) + \overline{F}_\mu(x) .
  \end{align*}
  Each partial integral over a set $S$ lies in $[0, \mu(S)]$. If moreover $\mu$ is concentrated
  on $[0, \infty)$, then writing $\overline{F} = \overline{F}_\mu$,
  \begin{align*}
    \overline{F}_{\mu*\mu}(x) \; \ge \; \overline{F}(x) \big( 1 - \overline{F}(a) \big)
      + \overline{F}(x - a) \big( \overline{F}(a) - \overline{F}(x) \big) + \overline{F}(x) .
  \end{align*}
\end{lemma}
\begin{proof}
  \leanok
  Split the integral of Lemma~\ref{lem:tail-conv} over $(-\infty, a]$, $(a, x]$ and $(x, \infty)$;
  on the last set $x - t < 0$, so the integrand equals $1$ and the integral is $\overline{F}_\mu(x)$.
  For the lower bound use $\overline{F}(x - t) \ge \overline{F}(x)$ on $(-\infty, a]$ and
  $\overline{F}(x - t) \ge \overline{F}(x - a)$ on $(a, x]$, together with
  Lemma~\ref{lem:tail-basic}.
\end{proof}

\begin{definition}[Convolution powers]
  \label{def:convPow}
  \lean{ExtremeValueTheory.Subexp.convPow}
  \leanok
  The $n$-fold convolution power of $\mu$ is defined by $\mu^{*0} = \delta_0$ and
  $\mu^{*(n+1)} = \mu * \mu^{*n}$. We write $\overline{F}^{*n}$ for its tail.
\end{definition}

\begin{lemma}[Properties of convolution powers]
  \label{lem:convPow-basic}
  \uses{def:convPow, lem:tail-conv}
  \lean{ExtremeValueTheory.Subexp.convPow_zero, ExtremeValueTheory.Subexp.convPow_succ,
    ExtremeValueTheory.Subexp.isProbabilityMeasure_convPow, ExtremeValueTheory.Subexp.convPow_one,
    ExtremeValueTheory.Subexp.convPow_Iio_zero}
  \leanok
  Convolution powers of a probability measure are probability measures, $\mu^{*1} = \mu$, and
  if $\mu$ is concentrated on $[0, \infty)$, then so is every $\mu^{*n}$.
\end{lemma}
\begin{proof}
  \leanok
  Induction on $n$ using Lemma~\ref{lem:tail-conv}, and $\mu * \delta_0 = \mu$.
\end{proof}

\begin{lemma}[Sums of i.i.d.\ variables]
  \label{lem:map-pi-sum-convPow}
  \uses{def:convPow}
  \lean{ExtremeValueTheory.map_pi_sum_eq_convPow}
  \leanok
  Under the product measure $\mu^{\otimes n}$ on $\bR^n$, the law of the sum of the
  coordinates $X_1 + \cdots + X_n$ is $\mu^{*n}$.
\end{lemma}
\begin{proof}
  \leanok
  Induction on $n$, splitting off the first coordinate: $\mu^{\otimes (n+1)} \cong \mu \otimes \mu^{\otimes n}$,
  and the push-forward of a product measure under addition is the convolution of the marginals.
\end{proof}


\section{Subexponentiality}

\begin{definition}[Subexponential distribution]
  \label{def:subexponential}
  \uses{def:tail}
  \lean{ExtremeValueTheory.IsSubexponential}
  \leanok
  A probability measure $\mu$ on $\bR$ is \emph{subexponential} if it is concentrated on
  $[0, \infty)$, its tail $\overline{F} = \overline{F}_\mu$ is positive everywhere, and
  \begin{align*}
    \lim_{x \to \infty} \frac{\overline{F}_{\mu * \mu}(x)}{\overline{F}(x)} \; = \; 2
  \end{align*}
  [EKM, Definition~1.3.3].
\end{definition}

\begin{theorem}[Regularly varying tails are subexponential]
  \label{thm:subexp-of-rv}
  \uses{def:subexponential, def:regularly-varying, lem:tail-conv-split}
  \lean{ExtremeValueTheory.IsSubexponential.of_isRegularlyVarying}
  \leanok
  If $\mu$ is a probability measure concentrated on $[0, \infty)$ whose tail $\overline{F}$ is
  regularly varying with index $-\alpha$, then $\mu$ is subexponential [EKM, Corollary~1.3.2].
\end{theorem}
\begin{proof}
  \leanok
  Positivity: $\overline{F}$ is antitone and eventually positive.
  Lower bound: Lemma~\ref{lem:tail-conv-split} with $a = x$ gives
  $\overline{F}_{\mu*\mu}(x) \ge 2 \overline{F}(x) - \overline{F}(x)^2$, so the ratio is at least
  $2 - \overline{F}(x) \to 2$.
  Upper bound: for $c \in (0,1)$, if $X_1 + X_2 > x$ then $X_1 > cx$, or $X_2 > cx$, or both
  $X_1, X_2 > (1-c)x$; hence
  $\overline{F}_{\mu*\mu}(x) \le 2 \overline{F}(cx) + \overline{F}((1-c)x)^2$. Dividing by
  $\overline{F}(x)$, the right-hand side tends to $2 c^{-\alpha} + (1-c)^{-\alpha} \cdot 0$ by
  regular variation. Letting $c \uparrow 1$ shows that the $\limsup$ is at most $2$.
\end{proof}

\begin{theorem}[Subexponential distributions are long-tailed]
  \label{thm:subexp-long-tailed}
  \uses{def:subexponential, lem:tail-conv-split}
  \lean{ExtremeValueTheory.IsSubexponential.tendsto_tail_sub_div}
  \leanok
  If $\mu$ is subexponential, then $\overline{F}(x - y)/\overline{F}(x) \to 1$ as $x \to \infty$,
  for every $y \in \bR$ [EKM, Lemma~1.3.5(a)].
\end{theorem}
\begin{proof}
  \leanok
  For $y \ge 0$: the ratio is at least $1$ by monotonicity. The lower bound of
  Lemma~\ref{lem:tail-conv-split} with $a = y$ gives
  \begin{align*}
    \frac{\overline{F}_{\mu*\mu}(x)}{\overline{F}(x)} \; \ge \; 1 - \overline{F}(y) + 1
      + \frac{\overline{F}(x - y)}{\overline{F}(x)} \big( \overline{F}(y) - \overline{F}(x) \big) ,
  \end{align*}
  and since the left-hand side tends to $2$ and $\overline{F}(y) - \overline{F}(x) \to \overline{F}(y) > 0$,
  we get $\limsup \overline{F}(x-y)/\overline{F}(x) \le 1$. For $y < 0$, substitute $x' = x - y$
  and invert.
\end{proof}

\begin{lemma}[Induction step]
  \label{lem:subexp-conv-step}
  \uses{def:subexponential, thm:subexp-long-tailed, lem:tail-conv-split}
  \lean{ExtremeValueTheory.IsSubexponential.tendsto_tail_conv_div}
  \leanok
  Let $\mu$ be subexponential and $\nu$ a probability measure concentrated on $[0, \infty)$ with
  $\overline{F}_\nu(x) / \overline{F}_\mu(x) \to c > 0$. Then
  $\overline{F}_{\mu*\nu}(x) / \overline{F}_\mu(x) \to c + 1$.
\end{lemma}
\begin{proof}
  \leanok
  Given $\eps > 0$, choose $v \ge 0$ such that $|\overline{F}_\nu(s)/\overline{F}_\mu(s) - c| < \eps$
  for $s \ge v$. Splitting the convolution integral at $a = x - v$
  (Lemma~\ref{lem:tail-conv-split}) and comparing $\overline{F}_\nu(x - t)$ with
  $(c \pm \eps) \overline{F}_\mu(x - t)$ on $(-\infty, x - v]$ gives
  \begin{align*}
    (c - \eps) \big( \overline{F}_{\mu*\mu}(x) - \overline{F}_\mu(x - v) \big) + \overline{F}_\mu(x)
    \; \le \; \overline{F}_{\mu*\nu}(x) \; \le \;
    (c + \eps) \big( \overline{F}_{\mu*\mu}(x) - \overline{F}_\mu(x) \big) + \overline{F}_\mu(x - v) .
  \end{align*}
  Dividing by $\overline{F}_\mu(x)$ and using subexponentiality and long-tailedness
  (Theorem~\ref{thm:subexp-long-tailed}), the limit points of the ratio lie in
  $[(c - \eps) + 1, (c + \eps) + 1]$. Let $\eps \to 0$.
\end{proof}

\begin{theorem}[Tails of convolution powers]
  \label{thm:subexp-convPow}
  \uses{lem:subexp-conv-step, lem:convPow-basic, lem:map-pi-sum-convPow}
  \lean{ExtremeValueTheory.IsSubexponential.tendsto_tail_convPow_div,
    ExtremeValueTheory.IsSubexponential.tendsto_measure_sum_div_tail}
  \leanok
  If $\mu$ is subexponential, then for every $n \ge 1$,
  $\overline{F}^{*n}(x) / \overline{F}(x) \to n$ as $x \to \infty$ [EKM, Lemma~1.3.4].
  Equivalently, for i.i.d.\ $X_1, \ldots, X_n \sim \mu$,
  $\PR[X_1 + \cdots + X_n > x] / \overline{F}(x) \to n$.
\end{theorem}
\begin{proof}
  \leanok
  Induction on $n \ge 1$: the case $n = 1$ is trivial since $\mu^{*1} = \mu$, and the step
  $\mu^{*(n+1)} = \mu * \mu^{*n}$ is Lemma~\ref{lem:subexp-conv-step} with $\nu = \mu^{*n}$ and
  $c = n$, using Lemma~\ref{lem:convPow-basic}. The random-variable form follows from
  Lemma~\ref{lem:map-pi-sum-convPow}.
\end{proof}

\begin{lemma}[Tail of the maximum]
  \label{lem:tail-max}
  \uses{def:tail}
  \lean{ExtremeValueTheory.tendsto_measure_max_div_tail}
  \leanok
  Let $\mu$ be a probability measure with $\overline{F} > 0$ everywhere and let
  $X_1, \ldots, X_n$ be i.i.d.\ with law $\mu$ (the coordinates under $\mu^{\otimes n}$). Then
  $\PR[\max(X_1, \ldots, X_n) > x] / \overline{F}(x) \to n$.
\end{lemma}
\begin{proof}
  \leanok
  $\PR[\max_i X_i > x] = 1 - (1 - \overline{F}(x))^n = \overline{F}(x) \sum_{k=0}^{n-1} (1 - \overline{F}(x))^k$,
  and the sum tends to $n$ since $\overline{F}(x) \to 0$.
\end{proof}

\begin{theorem}[Single big jump principle]
  \label{thm:single-big-jump}
  \uses{thm:subexp-convPow, lem:tail-max}
  \lean{ExtremeValueTheory.IsSubexponential.tendsto_measure_sum_div_measure_max}
  \leanok
  If $\mu$ is subexponential, $n \ge 1$, and $X_1, \ldots, X_n$ are i.i.d.\ with law $\mu$, then
  \begin{align*}
    \lim_{x \to \infty} \frac{\PR[X_1 + \cdots + X_n > x]}{\PR[\max(X_1, \ldots, X_n) > x]} \; = \; 1
  \end{align*}
  [EKM, Lemma~1.3.1(a) and Lemma~1.3.4].
\end{theorem}
\begin{proof}
  \leanok
  Divide numerator and denominator by $\overline{F}(x)$; both tend to $n \ne 0$ by
  Theorem~\ref{thm:subexp-convPow} and Lemma~\ref{lem:tail-max}.
\end{proof}

\begin{theorem}[Kesten's bound]
  \label{thm:kesten}
  \uses{def:subexponential, def:convPow, lem:tail-conv-split, lem:convPow-basic}
  \lean{ExtremeValueTheory.IsSubexponential.kesten_bound}
  \leanok
  If $\mu$ is subexponential, then for every $\eps > 0$ there is $K$ such that
  $\overline{F}^{*n}(x) \le K (1 + \eps)^n \, \overline{F}(x)$ for all $n \in \bN$ and all
  $x \in \bR$ [EKM, Lemma~1.3.5(c)].
\end{theorem}
\begin{proof}
  \leanok
  If $\overline{F}_\nu \le a \overline{F}$ everywhere ($\nu$ concentrated on $[0,\infty)$), then
  splitting the convolution integral gives
  $\overline{F}_{\mu*\nu}(x) \le a \big( \overline{F}_{\mu*\mu}(x) - \overline{F}(x) \big) + \overline{F}(x)$.
  Choose $x_0$ such that $\overline{F}_{\mu*\mu}(x) - \overline{F}(x) \le (1 + \eps) \overline{F}(x)$
  for $x \ge x_0$, and put $A = 1/\overline{F}(x_0)$, so that $1 \le A \overline{F}(x)$ for $x \le x_0$.
  One proves by induction on $n$ that
  $\overline{F}^{*n}(x) \le \big( K (1 + \eps)^n - A/\eps \big) \overline{F}(x)$ for a suitable
  constant $K$: for $x \ge x_0$ use the displayed estimate with
  $a = K (1+\eps)^n - A/\eps$, and for $x < x_0$ use $\overline{F}^{*(n+1)}(x) \le 1 \le A \overline{F}(x)$.
\end{proof}


\chapter{Stable laws and the generalized central limit theorem}
\label{chap:stable-laws}

A probability measure $\mu$ on $\bR$ is \emph{stable} if sums of i.i.d.\ copies of $X \sim \mu$
have, up to an affine change, the law of $X$. Stable laws are exactly the possible limits of
normalized sums of i.i.d.\ random variables, and for heavy-tailed summands with regularly varying
tails the limit is non-Gaussian. We follow Feller, \emph{An Introduction to Probability Theory
and Its Applications}, Vol.~II (2nd ed., Wiley, 1971) [Feller]; Gnedenko and Kolmogorov,
\emph{Limit Distributions for Sums of Independent Random Variables} (1954) [GK]; Samorodnitsky
and Taqqu, \emph{Stable Non-Gaussian Random Processes} (1994) [ST]; and [EKM, \S2.2].

We write $\varphi_\mu(t) = \int e^{itx} \, \ud \mu(x)$ for the characteristic function and
$\mu^{*n}$ for the convolution power (Definition~\ref{def:convPow}).

\section{Stable and symmetric stable laws}

\begin{lemma}[Characteristic functions of convolution powers and affine images]
  \label{lem:charfun-convpow}
  \uses{def:convPow}
  \lean{ExtremeValueTheory.charFun_convPow, ExtremeValueTheory.charFun_map_mul_add,
    ExtremeValueTheory.ne_dirac_of_norm_charFun_lt_one}
  \leanok
  For a probability measure $\mu$, $\varphi_{\mu^{*n}} = \varphi_\mu^n$, and the law of $aX + b$,
  $X \sim \mu$, has characteristic function $t \mapsto \varphi_\mu(a t) \, e^{i b t}$. A measure
  whose characteristic function has modulus $< 1$ somewhere is not a Dirac mass.
\end{lemma}
\begin{proof}
  \leanok
  The characteristic function of a convolution is the product of the characteristic functions;
  induct on $n$. The affine formula is a change of variables. A Dirac mass $\delta_x$ has
  $|\varphi(t)| = |e^{itx}| = 1$.
\end{proof}

\begin{lemma}[Characteristic function criterion]
  \label{lem:convpow-eq-map-of-charfun}
  \uses{lem:charfun-convpow}
  \lean{ExtremeValueTheory.convPow_eq_map_of_charFun}
  \leanok
  If $\varphi_\mu(t)^n = \varphi_\mu(a t) \, e^{i b t}$ for all $t$, then $\mu^{*n}$ is the law
  of $aX + b$, $X \sim \mu$.
\end{lemma}
\begin{proof}
  \leanok
  Both sides are finite measures with the same characteristic function
  (Lemma~\ref{lem:charfun-convpow}), and characteristic functions determine finite measures.
\end{proof}

\begin{definition}[Stable and strictly stable laws]
  \label{def:stable}
  \lean{ExtremeValueTheory.IsStable, ExtremeValueTheory.IsStrictlyStable,
    ExtremeValueTheory.IsStrictlyStable.isStable}
  \leanok
  A probability measure $\mu$ on $\bR$ is \emph{stable} if it is not a Dirac mass and for every
  $n \ge 1$ there are $a_n > 0$, $b_n \in \bR$ with $\mu^{*n} = \mathrm{law}(a_n X + b_n)$,
  $X \sim \mu$ [Feller, \S VI.1], [GK, \S33], [ST, Definition~1.1.4]. It is \emph{strictly
  stable} if one can take $b_n = 0$. Strictly stable laws are stable.
\end{definition}

\begin{proposition}[Gaussian laws are stable]
  \label{prop:gaussian-stable}
  \uses{def:stable, lem:convpow-eq-map-of-charfun}
  \lean{ExtremeValueTheory.norm_charFun_gaussianReal_one_lt, ExtremeValueTheory.isStable_gaussianReal,
    ExtremeValueTheory.isStrictlyStable_gaussianReal}
  \leanok
  For $v > 0$, the Gaussian law $N(m, v)$ is stable, with
  $N(m, v)^{*n} = \mathrm{law}\big( \sqrt{n} X + (n - \sqrt{n}) m \big)$, and $N(0, v)$ is
  strictly stable.
\end{proposition}
\begin{proof}
  \leanok
  With $\varphi(t) = e^{imt - vt^2/2}$ one checks $\varphi(t)^n = \varphi(\sqrt{n} t) \,
  e^{i (n - \sqrt n) m t}$ and applies Lemma~\ref{lem:convpow-eq-map-of-charfun}; $|\varphi(1)| =
  e^{-v/2} < 1$ excludes Dirac masses.
\end{proof}

\begin{definition}[Symmetric $\alpha$-stable laws]
  \label{def:symm-stable}
  \lean{ExtremeValueTheory.IsSymmStable, ExtremeValueTheory.IsSymmStable.isProbabilityMeasure',
    ExtremeValueTheory.IsSymmStable.charFun_eq}
  \leanok
  A probability measure $\mu$ on $\bR$ is \emph{symmetric $\alpha$-stable with constant $c$},
  written $S\alpha S(c)$, if $\varphi_\mu(t) = \exp(-c |t|^\alpha)$ for all $t$
  [ST, (1.2.3)], [Feller, \S XVII.4], [EKM, \S2.2].
\end{definition}

\begin{lemma}[Closure properties of symmetric stable laws]
  \label{lem:symm-stable-closure}
  \uses{def:symm-stable, lem:charfun-convpow}
  \lean{ExtremeValueTheory.IsSymmStable.unique, ExtremeValueTheory.IsSymmStable.conv,
    ExtremeValueTheory.IsSymmStable.map_mul, ExtremeValueTheory.IsSymmStable.map_neg,
    ExtremeValueTheory.IsSymmStable.convPow}
  \leanok
  An $S\alpha S(c)$ law is determined by $(\alpha, c)$. If $\mu$ is $S\alpha S(c)$ and $\nu$ is
  $S\alpha S(d)$, then $\mu * \nu$ is $S\alpha S(c + d)$, $\mu^{*n}$ is $S\alpha S(nc)$, the law
  of $rX$ is $S\alpha S(c |r|^\alpha)$, and the law of $-X$ is $\mu$.
\end{lemma}
\begin{proof}
  \leanok
  Characteristic functions determine laws; multiply characteristic functions for convolutions
  and use $\varphi_{rX}(t) = \varphi_X(rt)$ (Lemma~\ref{lem:charfun-convpow}).
\end{proof}

\begin{theorem}[Symmetric stable laws are strictly stable]
  \label{thm:symm-stable-strictly-stable}
  \uses{lem:symm-stable-closure, lem:convpow-eq-map-of-charfun, def:stable}
  \lean{ExtremeValueTheory.IsSymmStable.isStrictlyStable, ExtremeValueTheory.IsSymmStable.isStable,
    ExtremeValueTheory.isSymmStable_gaussianReal}
  \leanok
  For $\alpha > 0$ and $c > 0$, an $S\alpha S(c)$ law is strictly stable with
  $a_n = n^{1/\alpha}$ [Feller, \S VI.1], [ST, Theorem~1.1.3]. The centred Gaussian $N(0, v)$ is
  $S2S(v/2)$.
\end{theorem}
\begin{proof}
  \leanok
  $\exp(-c|t|^\alpha)^n = \exp(-c |n^{1/\alpha} t|^\alpha)$, so
  Lemma~\ref{lem:convpow-eq-map-of-charfun} applies with $b = 0$; $|\varphi(1)| = e^{-c} < 1$.
  For the Gaussian, $\varphi(t) = e^{-v t^2/2}$.
\end{proof}

\section{The Cauchy law}

\begin{lemma}[Fourier transform of $e^{-2\pi|v|}$]
  \label{lem:fourier-exp-abs}
  \lean{ExtremeValueTheory.CauchyCharFun.expAbs, ExtremeValueTheory.CauchyCharFun.kernel_Ioi,
    ExtremeValueTheory.CauchyCharFun.kernel_Iic, ExtremeValueTheory.CauchyCharFun.re_neg,
    ExtremeValueTheory.CauchyCharFun.re_pos, ExtremeValueTheory.CauchyCharFun.integrableOn_Ioi,
    ExtremeValueTheory.CauchyCharFun.integrableOn_Iic, ExtremeValueTheory.CauchyCharFun.integrable_expAbs,
    ExtremeValueTheory.CauchyCharFun.continuous_expAbs, ExtremeValueTheory.CauchyCharFun.fourier_expAbs,
    ExtremeValueTheory.CauchyCharFun.integral_cexp_mul_cauchyPDFReal}
  \leanok
  With the normalization $\hat f(w) = \int e^{-2\pi i v w} f(v) \, \ud v$, the Fourier transform of
  $v \mapsto e^{-2\pi |v|}$ is the standard Cauchy density $w \mapsto 1/(\pi(1 + w^2))$.
  Consequently $\int e^{itw} \, \frac{\ud w}{\pi(1 + w^2)} = e^{-|t|}$.
\end{lemma}
\begin{proof}
  \leanok
  Split the integral at $0$; on each half-line the integrand is a complex exponential
  $e^{z v}$ with $\mathrm{Re}\, z$ of the right sign, and
  $\frac{1}{2\pi(1 + iw)} + \frac{1}{2\pi(1 - iw)} = \frac{1}{\pi(1 + w^2)}$. The second claim
  follows by Fourier inversion, $e^{-2\pi|v|}$ being continuous and integrable.
\end{proof}

\begin{theorem}[The Cauchy law is stable]
  \label{thm:cauchy-stable}
  \uses{lem:fourier-exp-abs, lem:convpow-eq-map-of-charfun, def:stable, def:symm-stable}
  \lean{ExtremeValueTheory.charFun_cauchyMeasure, ExtremeValueTheory.norm_charFun_cauchyMeasure_one_lt,
    ExtremeValueTheory.isStrictlyStable_cauchyMeasure, ExtremeValueTheory.isStable_cauchyMeasure,
    ExtremeValueTheory.isSymmStable_cauchyMeasure}
  \leanok
  The Cauchy law with location $x_0$ and scale $\gamma$ (density
  $\gamma/(\pi((x - x_0)^2 + \gamma^2))$) has characteristic function
  $\varphi(t) = \exp(i t x_0 - \gamma |t|)$ [Feller, \S XV.2]. For $\gamma \ne 0$ it is strictly
  stable with $a_n = n$, i.e.\ $(X_1 + \cdots + X_n)/n \sim X_1$, and $\mathrm{Cauchy}(0, \gamma)$
  is $S1S(\gamma)$.
\end{theorem}
\begin{proof}
  \leanok
  The characteristic function follows from Lemma~\ref{lem:fourier-exp-abs} by an affine change
  of variables. Then $\varphi(t)^n = \varphi(nt)$ and Lemma~\ref{lem:convpow-eq-map-of-charfun}
  applies; $|\varphi(1)| = e^{-\gamma} < 1$.
\end{proof}

\section{Lévy's continuity theorem}

\begin{theorem}[Lévy's continuity theorem]
  \label{thm:levy-continuity}
  \lean{ExtremeValueTheory.continuous_charFun_probabilityMeasure,
    ExtremeValueTheory.exists_tendsto_of_tendsto_charFun}
  \leanok
  Let $\mu_n$ be probability measures on a finite-dimensional real inner product space whose
  characteristic functions converge pointwise to a function $f$ that is continuous at $0$. Then
  there is a probability measure $\nu$ with $\varphi_\nu = f$, and $\mu_n \to \nu$ weakly
  [Feller, \S XV.3, Theorem~2].
\end{theorem}
\begin{proof}
  \leanok
  By the tightness criterion (pointwise convergence of characteristic functions to a function
  continuous at $0$ implies tightness) and Prokhorov's theorem, the closure of
  $\set{\mu_n}$ is compact in the weak topology, so the sequence has a cluster point $\nu$. Since
  $\nu \mapsto \varphi_\nu(t)$ is continuous for the weak topology, $\varphi_\nu = f$. Lévy's
  convergence theorem with known limit then gives $\mu_n \to \nu$.
\end{proof}

\section{The stable constant}

\begin{lemma}[Elementary integral bounds]
  \label{lem:stable-integral-bounds}
  \lean{ExtremeValueTheory.StableIntegral.one_sub_cos_nonneg,
    ExtremeValueTheory.StableIntegral.one_sub_cos_le_two,
    ExtremeValueTheory.StableIntegral.one_sub_cos_le_sq,
    ExtremeValueTheory.StableIntegral.abs_sin_mul_rpow_le,
    ExtremeValueTheory.StableIntegral.one_sub_cos_mul_rpow_le,
    ExtremeValueTheory.StableIntegral.one_sub_cos_mul_rpow_nonneg,
    ExtremeValueTheory.StableIntegral.intervalIntegrable_rpow_one_sub,
    ExtremeValueTheory.StableIntegral.intervalIntegrable_sin_mul_rpow,
    ExtremeValueTheory.StableIntegral.intervalIntegrable_one_sub_cos_mul_rpow,
    ExtremeValueTheory.StableIntegral.integrableOn_one_sub_cos_mul_rpow,
    ExtremeValueTheory.StableIntegral.integral_one_sub_cos_mul_rpow}
  \leanok
  $0 \le 1 - \cos x \le \min(2, x^2/2)$ and $|t \sin(tu) u^{-\alpha}| \le t^2 u^{1-\alpha}$ for
  $u > 0$. For $0 < \alpha < 2$, $u \mapsto (1 - \cos(tu)) u^{-\alpha-1}$ is integrable on
  $(0, \infty)$ and
  \begin{align*}
    \int_0^\infty (1 - \cos(tu)) \, u^{-\alpha-1} \, \ud u \; = \;
    |t|^\alpha \int_0^\infty (1 - \cos u) \, u^{-\alpha-1} \, \ud u .
  \end{align*}
\end{lemma}
\begin{proof}
  \leanok
  Near $0$ the integrand is bounded by $\frac{t^2}{2} u^{1-\alpha}$, integrable since $\alpha < 2$;
  near $\infty$ by $2 u^{-\alpha-1}$, integrable since $\alpha > 0$. The scaling identity is the
  substitution $u \mapsto u/|t|$ (and evenness of $\cos$).
\end{proof}

\begin{definition}[The stable constant]
  \label{def:stable-const}
  \uses{lem:stable-integral-bounds}
  \lean{ExtremeValueTheory.StableIntegral.stableConst, ExtremeValueTheory.StableIntegral.stableConst_pos}
  \leanok
  For $0 < \alpha < 2$ let $C_\alpha = \alpha \int_0^\infty (1 - \cos u) \, u^{-\alpha-1} \, \ud u$
  [Feller, \S XVII.4]. Then $C_\alpha > 0$. (Classically $C_\alpha = \Gamma(1-\alpha)
  \cos(\pi\alpha/2)$ for $\alpha \ne 1$ and $C_1 = \pi/2$; only positivity is used.)
\end{definition}

\begin{lemma}[Integration by parts]
  \label{lem:stable-integral-by-parts}
  \uses{def:stable-const, lem:stable-integral-bounds}
  \lean{ExtremeValueTheory.StableIntegral.intervalIntegral_sin_mul_rpow_eq,
    ExtremeValueTheory.StableIntegral.tendsto_intervalIntegral_sin_mul_rpow}
  \leanok
  For $\alpha < 2$ and $M > 0$,
  \begin{align*}
    \int_0^M t \sin(tu) \, u^{-\alpha} \, \ud u \; = \; (1 - \cos(tM)) M^{-\alpha}
      + \alpha \int_0^M (1 - \cos(tu)) \, u^{-\alpha-1} \, \ud u ,
  \end{align*}
  and for $0 < \alpha < 2$ this tends to $C_\alpha |t|^\alpha$ as $M \to \infty$.
\end{lemma}
\begin{proof}
  \leanok
  Integrate by parts on $[\eps, M]$ with $\frac{\ud}{\ud u}(1 - \cos(tu)) = t \sin(tu)$, and let
  $\eps \to 0$ using $(1 - \cos(t\eps))\eps^{-\alpha} \le \frac{t^2}{2} \eps^{2-\alpha} \to 0$.
  As $M \to \infty$, the boundary term is at most $2 M^{-\alpha} \to 0$, and
  Lemma~\ref{lem:stable-integral-bounds} identifies the limit of the integral.
\end{proof}

\section{The generalized central limit theorem}

\begin{definition}[Two-sided tail]
  \label{def:abs-tail}
  \lean{ExtremeValueTheory.absTail, ExtremeValueTheory.GCLT.measurableSet_abs_gt,
    ExtremeValueTheory.GCLT.absTail_antitone, ExtremeValueTheory.GCLT.absTail_le_one}
  \leanok
  The two-sided tail of a measure $\mu$ on $\bR$ is $G(x) = \mu\set{ y \; | \; |y| > x }$, i.e.\
  $G(x) = \PR[|X| > x]$ for $X \sim \mu$. It is antitone, and at most $1$ for a probability
  measure.
\end{definition}

\begin{lemma}[Layer cake formula and truncation]
  \label{lem:gclt-layer-cake}
  \uses{def:abs-tail}
  \lean{ExtremeValueTheory.GCLT.integrable_one_sub_cos, ExtremeValueTheory.GCLT.integral_one_sub_cos_min,
    ExtremeValueTheory.GCLT.abs_integral_sub_le}
  \leanok
  For $a, M > 0$,
  \begin{align*}
    \int \big( 1 - \cos(t \min(|x|/a, M)) \big) \, \ud \mu(x) \; = \;
      \int_0^M t \sin(tu) \, G(au) \, \ud u ,
  \end{align*}
  and replacing $\min(|x|/a, M)$ by $|x|/a$ changes the left side by at most $2 G(aM)$.
\end{lemma}
\begin{proof}
  \leanok
  Write $1 - \cos(t s) = \int_0^s t \sin(tu) \, \ud u$ and apply Fubini. The two integrands differ
  only on $\set{|x| > aM}$, where both lie in $[0, 2]$.
\end{proof}

\begin{lemma}[Normalizing sequences and Potter domination]
  \label{lem:gclt-domination}
  \uses{def:regularly-varying, thm:rv-potter-monotone}
  \lean{ExtremeValueTheory.GCLT.tendsto_atTop_of_tendsto_mul, ExtremeValueTheory.GCLT.tendsto_mul_comp_mul,
    ExtremeValueTheory.GCLT.eventually_mul_comp_mul_le}
  \leanok
  Let $G$ be antitone, positive, at most $1$, and regularly varying with index $-\alpha$,
  $\alpha > 0$, and let $n \, G(a_n) \to 1$. Then $a_n \to \infty$,
  $n \, G(a_n u) \to u^{-\alpha}$ for every $u > 0$, and for every $\delta > 0$ there is $K$ such
  that eventually $n \, G(a_n u) \le K (u^{-\alpha-\delta} + u^{-\alpha+\delta})$ for all
  $u > 0$.
\end{lemma}
\begin{proof}
  \leanok
  If $a_n$ did not tend to $\infty$, then $n \, G(a_n) \ge n \, G(A) \to \infty$ along a
  subsequence. Then $n \, G(a_n u) = n \, G(a_n) \cdot G(a_n u)/G(a_n) \to u^{-\alpha}$. The
  domination follows from the Potter bounds (Theorem~\ref{thm:rv-potter-monotone}) for
  $a_n u$ large, and from $G \le 1$ together with $n \approx 1/G(a_n)$ for $a_n u$ bounded.
\end{proof}

\begin{lemma}[Dominated convergence on $(0, M)$]
  \label{lem:gclt-dominated}
  \uses{lem:gclt-domination, lem:stable-integral-bounds}
  \lean{ExtremeValueTheory.GCLT.tendsto_setIntegral_sin_mul}
  \leanok
  Under the hypotheses of Lemma~\ref{lem:gclt-domination} with $\alpha < 2$, for every $t$ and
  $M > 0$,
  $\int_0^M t \sin(tu) \, n \, G(a_n u) \, \ud u \to \int_0^M t \sin(tu) \, u^{-\alpha} \, \ud u$.
\end{lemma}
\begin{proof}
  \leanok
  Choose $\delta > 0$ with $\alpha + \delta < 2$. Since $|t \sin(tu)| \le t^2 u$, the integrands
  are dominated by $K t^2 (u^{1-\alpha-\delta} + u^{1-\alpha+\delta})$, integrable on $(0, M)$;
  apply dominated convergence.
\end{proof}

\begin{lemma}[An approximation argument]
  \label{lem:gclt-approx}
  \lean{ExtremeValueTheory.GCLT.tendsto_of_approx}
  \leanok
  Suppose $|L_n - I_M(n)| \le R_M(n)$ for all $M > 0$ and $n$, with $I_M(n) \to J(M)$ and
  $R_M(n) \to r(M)$ as $n \to \infty$, $J(M) \to c$ and $r(M) \to 0$ as $M \to \infty$. Then
  $L_n \to c$.
\end{lemma}
\begin{proof}
  \leanok
  Given $\eps > 0$, choose $M$ with $|J(M) - c| < \eps/4$ and $r(M) < \eps/4$; then eventually
  $|L_n - c| < \eps$.
\end{proof}

\begin{theorem}[Tauberian step]
  \label{thm:gclt-tauberian}
  \uses{lem:gclt-layer-cake, lem:gclt-domination, lem:gclt-dominated, lem:gclt-approx,
    lem:stable-integral-by-parts, def:stable-const}
  \lean{ExtremeValueTheory.GCLT.tendsto_mul_integral_one_sub_cos}
  \leanok
  Let $\mu$ be a probability measure whose two-sided tail $G$ is regularly varying with index
  $-\alpha$, $0 < \alpha < 2$, and let $a_n > 0$ with $n \, G(a_n) \to 1$. Then for every $t$,
  \begin{align*}
    n \int \big( 1 - \cos(t x / a_n) \big) \, \ud \mu(x) \; \to \; C_\alpha |t|^\alpha
  \end{align*}
  [Feller, \S XVII.5], [GK, \S35]. No symmetry is needed.
\end{theorem}
\begin{proof}
  \leanok
  Apply Lemma~\ref{lem:gclt-approx} with $L_n$ the left side, $I_M(n) = n \int_0^M t \sin(tu) \,
  G(a_n u) \, \ud u$ (Lemma~\ref{lem:gclt-layer-cake}) and $R_M(n) = 2 n \, G(a_n M)$. Then
  $I_M(n) \to \int_0^M t \sin(tu) \, u^{-\alpha} \, \ud u =: J(M)$ by
  Lemma~\ref{lem:gclt-dominated}, $R_M(n) \to 2 M^{-\alpha} \to 0$ by
  Lemma~\ref{lem:gclt-domination}, and $J(M) \to C_\alpha |t|^\alpha$ by
  Lemma~\ref{lem:stable-integral-by-parts}.
\end{proof}

\begin{lemma}[Characteristic function of normalized sums]
  \label{lem:gclt-charfun}
  \uses{thm:gclt-tauberian, lem:charfun-convpow}
  \lean{ExtremeValueTheory.GCLT.charFun_eq_integral_cos, ExtremeValueTheory.GCLT.tendsto_charFun_map_convPow,
    ExtremeValueTheory.charFun_map_mul_sum}
  \leanok
  For a symmetric law $\mu$, $\varphi_\mu(s) = \int \cos(sx) \, \ud \mu(x)$ is real. Under the
  hypotheses of Theorem~\ref{thm:gclt-tauberian} with $\mu$ symmetric, the law of
  $(X_1 + \cdots + X_n)/a_n$ has characteristic function
  $\varphi_\mu(t/a_n)^n \to \exp(-C_\alpha |t|^\alpha)$. For i.i.d.\ random variables,
  $\varphi_{c(X_0 + \cdots + X_{n-1})}(t) = \varphi_{X_0}(ct)^n$.
\end{lemma}
\begin{proof}
  \leanok
  Symmetry kills the imaginary part. Write $\varphi_\mu(t/a_n) = 1 + g_n$ with
  $n g_n = -n \int (1 - \cos(tx/a_n)) \, \ud \mu \to -C_\alpha|t|^\alpha$ by
  Theorem~\ref{thm:gclt-tauberian}; then $(1 + g_n)^n \to \exp(-C_\alpha|t|^\alpha)$. The
  characteristic function of the normalized sum is computed with Lemma~\ref{lem:charfun-convpow}
  and independence.
\end{proof}

\begin{theorem}[Generalized central limit theorem, symmetric case]
  \label{thm:gclt}
  \uses{lem:gclt-charfun, thm:levy-continuity, def:symm-stable, def:stable-const,
    def:convergence-in-distribution}
  \lean{ExtremeValueTheory.tendsto_map_convPow_of_isRegularlyVarying,
    ExtremeValueTheory.tendstoInDistribution_inv_mul_sum_of_isRegularlyVarying}
  \leanok
  Let $\mu$ be a symmetric probability law on $\bR$ whose two-sided tail $G(x) = \PR[|X| > x]$ is
  regularly varying with index $-\alpha$, $0 < \alpha < 2$, and let $a_n > 0$ with
  $n \, G(a_n) \to 1$. Then there is a probability measure $\nu$, symmetric $\alpha$-stable with
  constant $C_\alpha$, such that the laws of $(X_1 + \cdots + X_n)/a_n$ converge weakly to
  $\nu$. Equivalently, for i.i.d.\ $X_0, X_1, \ldots$ with law $\mu$ and $Y \sim S\alpha S(C_\alpha)$,
  \begin{align*}
    \frac{X_0 + \cdots + X_{n-1}}{a_n} \; \tendsinlaw \; Y
  \end{align*}
  [Feller, \S XVII.5, Theorem~2], [GK, \S35, Theorem~2], [EKM, Theorem~2.2.8].
\end{theorem}
\begin{proof}
  \leanok
  By Lemma~\ref{lem:gclt-charfun}, the characteristic functions converge pointwise to
  $\exp(-C_\alpha|t|^\alpha)$, which is continuous. Lévy's continuity theorem
  (Theorem~\ref{thm:levy-continuity}) produces the limit $\nu$ with this characteristic function
  and the weak convergence; in particular, the limit exists without assuming the existence of
  stable laws. The random-variable form follows from the formula for characteristic functions of
  normalized i.i.d.\ sums in Lemma~\ref{lem:gclt-charfun} and uniqueness of $S\alpha S$ laws
  (Lemma~\ref{lem:symm-stable-closure}).
\end{proof}

\begin{lemma}[Symmetrized Pareto law]
  \label{lem:symm-pareto}
  \uses{def:abs-tail, def:regularly-varying}
  \lean{ExtremeValueTheory.SymmPareto.law, ExtremeValueTheory.SymmPareto.pareto_Ioi,
    ExtremeValueTheory.SymmPareto.pareto_Iio_one, ExtremeValueTheory.SymmPareto.isProbabilityMeasure,
    ExtremeValueTheory.SymmPareto.map_neg, ExtremeValueTheory.SymmPareto.absTail_eq,
    ExtremeValueTheory.SymmPareto.isRegularlyVarying, ExtremeValueTheory.SymmPareto.tendsto_mul_absTail}
  \leanok
  For $\alpha > 0$, the symmetrized Pareto law $\frac{1}{2}\big( \mathrm{Pareto}(1, \alpha) +
  \mathrm{law}(-X) \big)$ is a symmetric probability measure with $\PR[|X| > x] = x^{-\alpha}$ for
  $x \ge 1$; its two-sided tail is regularly varying with index $-\alpha$, and
  $a_n = (n + 1)^{1/\alpha}$ satisfies $n \, G(a_n) \to 1$.
\end{lemma}
\begin{proof}
  \leanok
  The Pareto law has $\PR[X > x] = x^{-\alpha}$ for $x \ge 1$ and no mass below $1$; then
  $n \, G(a_n) = n/(n + 1) \to 1$.
\end{proof}

\begin{theorem}[Existence of symmetric stable laws]
  \label{thm:exists-symm-stable}
  \uses{thm:gclt, lem:symm-pareto, lem:symm-stable-closure, def:stable-const,
    thm:symm-stable-strictly-stable}
  \lean{ExtremeValueTheory.exists_isSymmStable_stableConst, ExtremeValueTheory.exists_isSymmStable}
  \leanok
  For every $0 < \alpha \le 2$ and $c \ge 0$ there is a probability measure on $\bR$ with
  characteristic function $t \mapsto \exp(-c |t|^\alpha)$ [Feller, \S XVII.4], [ST, \S1.1].
\end{theorem}
\begin{proof}
  \leanok
  For $\alpha < 2$, the generalized central limit theorem (Theorem~\ref{thm:gclt}) applied to the
  symmetrized Pareto law (Lemma~\ref{lem:symm-pareto}) yields an $S\alpha S(C_\alpha)$ law;
  scaling by $r = (c/C_\alpha)^{1/\alpha}$ (Lemma~\ref{lem:symm-stable-closure}, using
  $C_\alpha > 0$) gives constant $c$. For $\alpha = 2$, take the Gaussian $N(0, 2c)$
  (Theorem~\ref{thm:symm-stable-strictly-stable}).
\end{proof}

\section{Asymmetric tails}

\begin{definition}[Stable characteristic exponent]
  \label{def:has-stable-charfun}
  \lean{ExtremeValueTheory.stableExponent, ExtremeValueTheory.HasStableCharFun,
    ExtremeValueTheory.StableIntegral.stableSkewConst}
  \leanok
  For $\alpha > 0$, $c \ge 0$ and $e \in \bR$, a probability measure $\mu$ has stable
  characteristic function with parameters $(\alpha, c, e)$ if
  \begin{align*}
    \varphi_\mu(t) \; = \; \exp\big( -c |t|^\alpha + i \, e \, \mathrm{sign}(t) \, |t|^\alpha \big) .
  \end{align*}
  The skewness constant is $E_\alpha = \alpha \int_0^\infty \sin u \, u^{-\alpha-1} \, \ud u$ for
  $\alpha < 1$ and $E_\alpha = -\int_0^\infty (1 - \cos u) \, u^{-\alpha} \, \ud u$ for
  $1 < \alpha < 2$. Classically $E_\alpha = C_\alpha \tan(\pi\alpha/2)$; this closed form is not
  formalized.
\end{definition}

\begin{theorem}[Laws with stable characteristic function are strictly stable]
  \label{thm:has-stable-charfun-stable}
  \uses{def:has-stable-charfun}
  \lean{ExtremeValueTheory.HasStableCharFun.isStrictlyStable,
    ExtremeValueTheory.HasStableCharFun.isStable,
    ExtremeValueTheory.HasStableCharFun.iff_isSymmStable}
  \leanok
  If $c > 0$, a law with stable characteristic function $(\alpha, c, e)$ is strictly stable with
  $a_n = n^{1/\alpha}$. With $e = 0$ these are exactly the $S\alpha S(c)$ laws.
\end{theorem}
\begin{proof}
  \leanok
  The exponent is homogeneous of degree $\alpha$ in $t$, so $\varphi(t)^n = \varphi(n^{1/\alpha} t)$.
\end{proof}

\begin{theorem}[Generalized central limit theorem, asymmetric tails]
  \label{thm:gclt-asymmetric}
  \uses{thm:gclt, def:has-stable-charfun, thm:has-stable-charfun-stable}
  \lean{ExtremeValueTheory.tendsto_map_convPow_of_isRegularlyVarying_of_lt_one,
    ExtremeValueTheory.tendsto_map_convPow_sub_of_isRegularlyVarying_of_one_lt}
  \leanok
  Let $\PR[|X| > x]$ be regularly varying with index $-\alpha$, let
  $\PR[X > x]/\PR[|X| > x] \to p$, and let $a_n > 0$ with $n \PR[|X| > a_n] \to 1$. Put
  $\beta = 2p - 1$. If $0 < \alpha < 1$, then $S_n / a_n$ converges weakly to the law with stable
  characteristic function $(\alpha, C_\alpha, \beta E_\alpha)$. If $1 < \alpha < 2$, the same
  holds for $(S_n - n \E X)/a_n$ [Feller, \S XVII.5, Theorem~2], [EKM, Theorem~2.2.8]. The case
  $\alpha = 1$ is not formalized.
\end{theorem}
\begin{proof}
  \leanok
  Split $n(\varphi(t/a_n) - 1)$ (minus the centering term) into its real part, handled as in the
  symmetric case, and its imaginary part, which is an integral of $\sin$ against the difference
  of the one-sided tails. By tail balance the right and left tails are regularly varying with
  weights $p$ and $q = 1 - p$. A layer-cake formula and Potter domination give the limit, and
  L\'evy's continuity theorem concludes.
\end{proof}

\begin{theorem}[Closed forms of the stable constants]
  \label{thm:stable-constants}
  \uses{def:stable-const, def:has-stable-charfun}
  \lean{ExtremeValueTheory.StableIntegral.stableConst_eq_pi_div,
    ExtremeValueTheory.StableIntegral.stableConst_one,
    ExtremeValueTheory.StableIntegral.stableConst_eq_Gamma_mul_cos,
    ExtremeValueTheory.StableIntegral.stableSkewConst_eq_Gamma_mul_sin,
    ExtremeValueTheory.StableIntegral.stableSkewConst_eq_mul_tan,
    ExtremeValueTheory.StableIntegral.integral_rpow_div_one_add_sq}
  \leanok
  For $0 < \alpha < 2$, $C_\alpha = \pi / (2 \sin(\pi\alpha/2) \Gamma(\alpha))$; in particular
  $C_1 = \pi/2$, and for $\alpha \ne 1$, $C_\alpha = \Gamma(1-\alpha)\cos(\pi\alpha/2)$,
  $E_\alpha = \Gamma(1-\alpha)\sin(\pi\alpha/2) = C_\alpha \tan(\pi\alpha/2)$ [ST, \S1.1],
  [Feller, \S XVII.4].
\end{theorem}
\begin{proof}
  \leanok
  Write $u^{-\alpha-1} = \Gamma(\alpha+1)^{-1} \int_0^\infty x^\alpha e^{-xu} \, \ud x$ and swap the
  integrals. The Laplace transforms of $1 - \cos$ and $\sin$ are rational, which reduces both
  constants to $\int_0^\infty x^{s-1}/(1+x^2) \, \ud x = \pi / (2 \sin(\pi s/2))$, itself obtained
  from $1/(1+x^2) = \int_0^\infty e^{-(1+x^2)y} \, \ud y$ and the reflection formula.
\end{proof}

\begin{corollary}[Generalized CLT, textbook form]
  \label{cor:gclt-textbook}
  \uses{thm:gclt-asymmetric, thm:stable-constants}
  \lean{ExtremeValueTheory.tendsto_map_convPow_of_isRegularlyVarying_of_lt_one',
    ExtremeValueTheory.tendsto_map_convPow_sub_of_isRegularlyVarying_of_one_lt',
    ExtremeValueTheory.stableExponent_mul_tan}
  \leanok
  Under the hypotheses of Theorem~\ref{thm:gclt-asymmetric}, the limit law has characteristic
  function $\exp\big(-C_\alpha |t|^\alpha (1 - i \beta \, \mathrm{sign}(t) \tan(\pi\alpha/2))\big)$
  with $C_\alpha = \Gamma(1-\alpha)\cos(\pi\alpha/2)$ and $\beta = 2p - 1$.
\end{corollary}
\begin{proof}
  \leanok
  Substitute Theorem~\ref{thm:stable-constants} into Theorem~\ref{thm:gclt-asymmetric}.
\end{proof}

\begin{theorem}[Generalized central limit theorem, $\alpha = 1$]
  \label{thm:gclt-one}
  \uses{thm:gclt-asymmetric, thm:stable-constants}
  \lean{ExtremeValueTheory.tendsto_map_convPow_sub_of_isRegularlyVarying_one,
    ExtremeValueTheory.tendsto_map_convPow_sub_of_isRegularlyVarying_one', ExtremeValueTheory.stableExponentOne,
    ExtremeValueTheory.HasStableOneCharFun, ExtremeValueTheory.StableIntegral.cosLogConst, ExtremeValueTheory.StableIntegral.stableLogConst}
  \leanok
  Under the hypotheses of Theorem~\ref{thm:gclt-asymmetric} with $\alpha = 1$, let
  $b_n = n \E[X \mathbf 1\{|X| \le a_n\}]$. Then $(S_n - b_n)/a_n$ converges weakly to the law with
  characteristic function
  \begin{align*}
    \exp\big( -\tfrac{\pi}{2} |t| + i (2p - 1) (c\, t - t \log |t|) \big),
  \end{align*}
  where $c = 1 + \int_0^1 (\cos v - 1)/v \, \ud v + \int_1^\infty \cos v / v \, \ud v$
  (classically $1 - \gamma$; this value is not formalized) [Feller, \S XVII.5, Theorem~2],
  [EKM, Theorem~2.2.8].
\end{theorem}
\begin{proof}
  \leanok
  Decompose $\E \sin(sX) - s \E[X \mathbf 1\{|X| \le a\}]$ into one-sided terms of the continuous
  odd function $Q_t(v) = \sin(tv) - t \, \mathrm{clamp}(v, -1, 1)$ plus a boundary term
  $t(\PR[X > a] - \PR[X < -a]) \to t(p - q)$. A truncated layer-cake formula and Potter
  domination give $n \E Q_t(X^+/a_n) \to p\, t (c' - \log|t|)$, where the improper integral
  $\int_0^\infty (t \cos(tu) - t \mathbf 1\{u < 1\})/u \, \ud u = t(c' - \log |t|)$ comes from the
  substitution $v = tu$. The centering is small: $n (t \E[X \mathbf 1\{|X|\le a_n\}]/a_n)^2 \to 0$.
  The exponent is continuous since $t \log |t| \to 0$, so L\'evy's continuity theorem applies.
\end{proof}

\begin{corollary}[The $\alpha = 1$ constant is $1 - \gamma$]
  \label{cor:gclt-one-euler}
  \uses{thm:gclt-one}
  \lean{ExtremeValueTheory.StableIntegral.cosLogConst_eq, ExtremeValueTheory.StableIntegral.stableLogConst_eq,
    ExtremeValueTheory.tendsto_map_convPow_sub_of_isRegularlyVarying_one_eulerMascheroni}
  \leanok
  $\int_0^1 (\cos v - 1)/v \ud v + \int_1^\infty \cos v / v \ud v = -\gamma$, so the constant $c$ in
  Theorem~\ref{thm:gclt-one} is $1 - \gamma$, with $\gamma$ the Euler--Mascheroni constant.
\end{corollary}
\begin{proof}
  \leanok
  Compare with $\int_0^\infty (e^{-v} - \mathbf 1\{v < 1\})/v \ud v = -\gamma$ and show
  $\int_0^\infty (\cos v - e^{-v})/v \ud v = 0$.
\end{proof}


\chapter{Extremes of stationary sequences}
\label{chap:dependence}

This chapter extends the limit theory of maxima from iid samples to stationary sequences, following
Leadbetter, Lindgren and Rootz\'en, \emph{Extremes and Related Properties of Random Sequences and
Processes}, Springer (1983) [LLR], Chapter~3, and [EKM, \S8.1]. Long-range independence of
maxima is the condition $D(u_n)$, and local clustering is controlled by $D'(u_n)$; the extremal
index $\theta$ measures how much extremes cluster.

\section{Stationary sequences and Leadbetter's conditions}

\begin{definition}[Stationary sequence]
  \label{def:stationary}
  \lean{ProbabilityTheory.IsStationary, ProbabilityTheory.IsStationary.map_shift, ProbabilityTheory.IsStationary.identDistrib,
    ProbabilityTheory.IsStationary.factor, ProbabilityTheory.isStationary_of_iIndepFun}
  \leanok
  A sequence $(X_t)_{t \in \bN}$ is stationary if the law of the shifted path $(X_{t+1})_t$
  equals the law of $(X_t)_t$. Then every shift preserves the law, the marginals are identically
  distributed, sliding-window functions of a stationary sequence are stationary, and iid sequences
  are stationary.
\end{definition}

\begin{definition}[Conditions $D(u_n)$ and $D'(u_n)$]
  \label{def:leadbetter}
  \uses{def:stationary}
  \lean{ExtremeValueTheory.partialMax, ExtremeValueTheory.maxLE, ExtremeValueTheory.ConditionD, ExtremeValueTheory.ConditionD'}
  \leanok
  Let $M_n = \max_{t < n} X_t$. Condition $D(u_n)$ holds if there are $l_n = o(n)$ and
  $\alpha_n \to 0$ such that for all $A, B \subseteq \{0, \dots, n-1\}$ separated by at least $l_n$,
  \begin{align*}
    \big| \PR[\max_{A \cup B} X \le u_n] - \PR[\max_A X \le u_n] \, \PR[\max_B X \le u_n] \big|
      \; \le \; \alpha_n .
  \end{align*}
  Condition $D'(u_n)$ holds if
  $\limsup_n n \sum_{1 \le j < n/k} \PR[X_0 > u_n, X_j > u_n] \to 0$ as $k \to \infty$
  [LLR, \S3.2, (3.4.3)].
\end{definition}

\begin{theorem}[Block approximation]
  \label{thm:block-approx}
  \uses{def:leadbetter}
  \lean{ProbabilityTheory.IsStationary.abs_measureReal_maxLE_range_sub_pow_le,
    ExtremeValueTheory.tendsto_measureReal_partialMax_sub_pow, MeasureTheory.sum_measureReal_sub_sum_inter_le_measureReal_biUnion}
  \leanok
  For a stationary sequence satisfying $D(u_n)$ with $n \PR[X_0 > u_n] \to \tau$, and fixed
  $k \ge 1$, $\PR[M_n \le u_n] - \PR[M_{\lfloor n/k \rfloor} \le u_n]^k \to 0$ [LLR, Lemma~3.3.1].
\end{theorem}
\begin{proof}
  \leanok
  Split $\{0, \dots, n-1\}$ into $k$ blocks separated by gaps of length $l_n$. The gaps cost at most
  $2 k l_n \PR[X_0 > u_n] \to 0$, and $D(u_n)$ factorizes the blocks with error $k \alpha_n$.
\end{proof}

\begin{theorem}[Leadbetter's theorem]
  \label{thm:leadbetter}
  \uses{thm:block-approx}
  \lean{ExtremeValueTheory.tendsto_measureReal_partialMax_le_exp, ProbabilityTheory.IsStationary.measureReal_maxLE_range_bounds}
  \leanok
  If a stationary sequence satisfies $D(u_n)$ and $D'(u_n)$ and $n \PR[X_0 > u_n] \to \tau$, then
  $\PR[M_n \le u_n] \to e^{-\tau}$, as for iid sequences [LLR, Theorem~3.4.1], [EKM, Theorem~4.4.3].
\end{theorem}
\begin{proof}
  \leanok
  By Bonferroni, $1 - \frac{n}{k} p_n \le \PR[M_{n/k} \le u_n] \le 1 - \frac{n}{k} p_n + \frac{n}{k}
  \sum_{1 \le j < n/k} \PR[X_0 > u_n, X_j > u_n]$ with $p_n = \PR[X_0 > u_n]$. Raise to the $k$-th
  power, use Theorem~\ref{thm:block-approx}, and let $k \to \infty$ using $D'(u_n)$.
\end{proof}

\section{The extremal index}

\begin{definition}[Extremal index]
  \label{def:extremal-index}
  \uses{def:leadbetter}
  \lean{ExtremeValueTheory.HasExtremalIndex}
  \leanok
  A stationary sequence has extremal index $\theta \in [0, 1]$ if for every $\tau > 0$ and every
  sequence $u_n$ with $n \PR[X_0 > u_n] \to \tau$, $\PR[M_n \le u_n] \to e^{-\theta\tau}$
  [LLR, \S3.7], [EKM, Definition~8.1.1].
\end{definition}

\begin{theorem}[Maxima of stationary sequences]
  \label{thm:stationary-max}
  \uses{def:extremal-index, thm:leadbetter, thm:three-types-of-extr-val-distr}
  \lean{ExtremeValueTheory.hasExtremalIndex_one, ExtremeValueTheory.tendsto_measureReal_partialMax_le_pow,
    ExtremeValueTheory.isExtremeValueDistr_pow, CumulativeDistributionFunction.IsExtremeValueDistr.continuous}
  \leanok
  Under $D$ and $D'$ at all levels, the extremal index is $1$. If the marginal $F$ satisfies
  $F^n(a_n x + b_n) \to H(x)$ for an extreme value distribution $H$ and the sequence has extremal
  index $\theta > 0$, then $\PR[M_n \le a_n x + b_n] \to H(x)^\theta$, and $H^\theta$ is an extreme
  value distribution of the same type as $H$ [LLR, Theorem~3.7.1], [EKM, Theorem~8.1.2].
\end{theorem}
\begin{proof}
  \leanok
  $n(1 - F(a_n x + b_n)) \to -\log H(x)$, so the definition of $\theta$ applies with
  $\tau = -\log H(x)$. For the type, $H^\theta = H^{\theta}$ is max-stable, and the
  classification gives an affine map with $A \cdot H = H^\theta$.
\end{proof}

\begin{proposition}[Examples]
  \label{prop:extremal-index-examples}
  \uses{thm:stationary-max}
  \lean{ExtremeValueTheory.measureReal_maxLE_of_iIndepFun, ExtremeValueTheory.conditionD_of_iIndepFun,
    ExtremeValueTheory.conditionD'_of_iIndepFun, ExtremeValueTheory.hasExtremalIndex_one_of_iIndepFun, ExtremeValueTheory.isStationary_movingMax,
    ExtremeValueTheory.hasExtremalIndex_movingMax}
  \leanok
  Iid sequences satisfy $D$ and $D'$ and have extremal index $1$. The moving maximum
  $X_t = \max(Z_t, Z_{t+1})$ of an iid sequence is stationary with extremal index $1/2$
  [EKM, Example~8.1.3].
\end{proposition}
\begin{proof}
  \leanok
  For the moving maximum, $\PR[M_n \le u] = G(u)^{n+1}$ with $G$ the law of $Z_0$, while
  $\PR[X_0 > u] \sim 2(1 - G(u))$.
\end{proof}

\section{Regularly varying time series and the tail process}

We follow Basrak and Segers, \emph{Regularly varying multivariate time series}, Stoch.\ Proc.\ Appl.119 (2009) [BS09], Planini\'c and Soulier, \emph{The tail process revisited}, Extremes 21 (2018)
[PS18], and Kulik and Soulier, \emph{Heavy-Tailed Time Series}, Springer (2020) [KS20].

\begin{definition}[Conditional limits, tail process and spectral tail process]
  \label{def:tail-process}
  \uses{def:regularly-varying}
  \lean{ExtremeValueTheory.TendstoLawCond, ExtremeValueTheory.HasSpectralTailProcess, ExtremeValueTheory.HasTailProcess,
    ExtremeValueTheory.HasShiftInvariantWindows}
  \leanok
  A sequence $(X_t)_{t \in \bZ}$ has spectral tail process $(\Theta_t)$ with index $\alpha > 0$ if
  $\PR[|X_0| > x]$ is regularly varying with index $-\alpha$ and, for every finite window
  $S \subset \bZ$, the conditional law of $(X_t/|X_0|)_{t \in S}$ given $|X_0| > x$ converges weakly to
  the law of $(\Theta_t)_{t \in S}$ as $x \to \infty$. It has tail process $(Y_t)$ if $|Y_0|$ is
  Pareto($\alpha$) and the conditional law of $(X_t/x)_{t\in S}$ given $|X_0| > x$ converges to that of
  $(Y_t)_{t \in S}$ [BS09, \S2], [KS20, Ch.~5].
\end{definition}

\begin{theorem}[Polar decomposition (Basrak--Segers)]
  \label{thm:basrak-segers}
  \uses{def:tail-process}
  \lean{ExtremeValueTheory.tendstoLawCond_polar, ExtremeValueTheory.isRegularlyVarying_of_tendstoLawCond,
    ExtremeValueTheory.tendstoLawCond_spectral_of_tendstoLawCond, ExtremeValueTheory.eq_map_prod_of_tendstoLawCond, ExtremeValueTheory.polar_iff,
    ExtremeValueTheory.HasSpectralTailProcess.hasTailProcess, ExtremeValueTheory.HasTailProcess.hasSpectralTailProcess,
    ExtremeValueTheory.HasTailProcess.map_window_eq, ExtremeValueTheory.HasSpectralTailProcess.ae_abs_eq_one}
  \leanok
  A sequence has a spectral tail process if and only if it has a tail process, and then
  $Y_t = |Y_0| \Theta_t$ with $|Y_0|$ Pareto($\alpha$) independent of $(\Theta_t)$, and $|\Theta_0| = 1$
  [BS09, Theorem~2.1].
\end{theorem}
\begin{proof}
  \leanok
  Regular variation makes $|X_0|/x$ asymptotically Pareto given $|X_0| > x$; conditioning further on
  the radial part shows that it decouples from the angular part $X/|X_0|$ in the limit.
\end{proof}

\begin{theorem}[Time-change formula]
  \label{thm:time-change}
  \uses{thm:basrak-segers, def:stationary}
  \lean{ExtremeValueTheory.HasSpectralTailProcess.lintegral_timeChange}
  \leanok
  For a stationary sequence with spectral tail process $\Theta$, all $i \in \bZ$, finite windows
  $S$ and measurable $g \ge 0$,
  \begin{align*}
    \E\big[ g((\Theta_{t-i})_{t\in S}) \, \mathbf 1\{\Theta_{-i} \ne 0\} \big]
      \; = \; \E\big[ g((\Theta_t / |\Theta_i|)_{t \in S}) \, |\Theta_i|^\alpha \big]
  \end{align*}
  [BS09, Theorem~3.1], [KS20, Theorem~5.3.1].
\end{theorem}
\begin{proof}
  \leanok
  Stationarity moves the conditioning from $|X_0| > x$ to $|X_i| > x$. Comparing the two
  conditionings in the tail-process limit at every level $\varepsilon > 0$ gives a measure identity,
  which is then rewritten in terms of $\Theta$.
\end{proof}

\begin{definition}[Candidate extremal index]
  \label{def:candidate-extremal-index}
  \uses{def:tail-process}
  \lean{ExtremeValueTheory.candidateExtremalIndex}
  \leanok
  $\vartheta = \E\big[\sup_{k \ge 0} |\Theta_k|^\alpha - \sup_{k \ge 1} |\Theta_k|^\alpha\big]$, with
  suprema in $[0, \infty]$ and $\infty - \infty = 0$ [KS20, \S7.5], [PS18].
\end{definition}

\begin{theorem}[Properties of the candidate extremal index]
  \label{thm:candidate-extremal-index}
  \uses{def:candidate-extremal-index, thm:time-change}
  \lean{ExtremeValueTheory.candidateExtremalIndex_le_one, ExtremeValueTheory.candidateExtremalIndex_eq_measure,
    ExtremeValueTheory.HasTailProcess.candidateExtremalIndex_eq, ExtremeValueTheory.HasSpectralTailProcess.ae_eq_zero_of_indepFun,
    ExtremeValueTheory.HasSpectralTailProcess.candidateExtremalIndex_eq_one_of_indepFun,
    ExtremeValueTheory.HasSpectralTailProcess.candidateExtremalIndex_pos}
  \leanok
  $\vartheta \le 1$ and $\vartheta = \PR[\sup_{k \ge 1} |Y_k| \le 1]$. If $X_t$ is independent of $X_0$
  then $\Theta_t = 0$; in particular $\vartheta = 1$ for iid sequences. For a stationary sequence
  with $\Theta_t \to 0$ as $t \to \infty$, $\vartheta > 0$ [PS18, \S4], [KS20, \S7.5].
\end{theorem}
\begin{proof}
  \leanok
  $|\Theta_0| = 1$ gives $\vartheta \le 1$. Writing $Y = R\Theta$ with $R$ Pareto independent of
  $\Theta$, $\PR[\sup_{k\ge1}|Y_k| \le 1 \mid \Theta] = \sup_{k\ge0}|\Theta_k|^\alpha -
  \sup_{k\ge1}|\Theta_k|^\alpha$. Positivity: decompose according to the last time $|Y_t| > 1$, which
  exists since $\Theta_t \to 0$, and apply the time-change formula to each piece.
\end{proof}

\section{The extremal index and the tail process}

\begin{lemma}[Stationarity notions agree]
  \label{lem:stationary-windows}
  \uses{def:stationary, def:tail-process}
  \lean{ExtremeValueTheory.hasShiftInvariantWindows_iff_isStationary, ExtremeValueTheory.HasShiftInvariantWindows.isStationary_nat,
    ExtremeValueTheory.HasShiftInvariantWindows.map_comp}
  \leanok
  For measurable $(X_t)_{t \in \bZ}$ under a finite measure, shift-invariance of all finite windows is
  equivalent to stationarity of every shifted $\bN$-indexed sequence $(X_{n+k})_{n \in \bN}$.
\end{lemma}
\begin{proof}
  \leanok
  Finite-dimensional laws determine the law of the path.
\end{proof}

\begin{definition}[Anticlustering and condition $\mathcal A$]
  \label{def:anticlustering}
  \uses{def:tail-process}
  \lean{ExtremeValueTheory.AntiClusteringAt, ExtremeValueTheory.ForwardAntiClusteringAt, ExtremeValueTheory.AntiClustering, ExtremeValueTheory.ConditionA}
  \leanok
  $AC(r_n)$ at levels $u_n$: $\lim_m \limsup_n \PR[\max_{m \le |t| \le r_n} |X_t| > u_n \mid |X_0| > u_n] = 0$
  (the forward version only uses $m \le t \le r_n$). Condition $\mathcal A(u_n)$:
  $\PR[M_n \le u_n] - \PR[M_{r_n} \le u_n]^{\lfloor n/r_n \rfloor} \to 0$ [BS09, \S4], [KS20, \S6.1].
\end{definition}

\begin{theorem}[Anticlustering forces $\Theta_t \to 0$]
  \label{thm:ac-theta-zero}
  \uses{def:anticlustering, thm:candidate-extremal-index}
  \lean{ExtremeValueTheory.HasSpectralTailProcess.tendsto_zero_of_forwardAntiClusteringAt,
    ExtremeValueTheory.HasSpectralTailProcess.candidateExtremalIndex_pos_of_forwardAntiClusteringAt,
    ExtremeValueTheory.HasSpectralTailProcess.candidateExtremalIndex_pos_of_antiClustering}
  \leanok
  Forward anticlustering at a single level sequence $u_n \to \infty$ with $r_n \to \infty$ implies
  $\Theta_t \to 0$ almost surely as $t \to \infty$; neither stationarity nor $r_n/n \to 0$ is used.
  For stationary sequences it follows that $\vartheta > 0$ [BS09, Proposition~4.2].
\end{theorem}
\begin{proof}
  \leanok
  The tail-process limit of the anticlustering probabilities at level $1$, extended to all levels
  through the independent Pareto radius $Y = R\Theta$.
\end{proof}

\begin{theorem}[The extremal index equals the candidate extremal index]
  \label{thm:theta-eq-vartheta}
  \uses{thm:ac-theta-zero, thm:time-change, def:extremal-index, lem:stationary-windows}
  \lean{ExtremeValueTheory.HasShiftInvariantWindows.measureReal_exists_lt_eq_sum, ExtremeValueTheory.HasSpectralTailProcess.tendsto_blockRatio,
    ExtremeValueTheory.HasSpectralTailProcess.tendsto_partialMax_abs, ExtremeValueTheory.HasSpectralTailProcess.hasExtremalIndex_abs, ExtremeValueTheory.HasSpectralTailProcess.hasExtremalIndex_of_nonneg,
    ExtremeValueTheory.HasSpectralTailProcess.tendsto_partialMax_abs_of_antiClustering}
  \leanok
  For a stationary sequence with spectral tail process $\Theta$, $u_n \to \infty$, $r_n \to \infty$
  and forward anticlustering at $u_n$,
  \begin{align*}
    \frac{\PR[\max_{t < r_n} |X_t| > u_n]}{r_n \PR[|X_0| > u_n]} \; \to \; \vartheta ;
  \end{align*}
  no mixing and no condition $r_n/n \to 0$ is needed for this. If in addition
  $n \PR[|X_0| > u_n] \to \tau$, $r_n/n \to 0$ and $\mathcal A(u_n)$ holds, then
  $\PR[\max_{t<n} |X_t| \le u_n] \to e^{-\vartheta\tau}$. Under $AC(r_n)$ and $\mathcal A(a_n x)$ at all
  levels, $|X|$ (and $X$ itself if $X \ge 0$) has extremal index $\vartheta$, and
  $\PR[M_n \le a_n x] \to \exp(-\vartheta x^{-\alpha})$ [BS09, Proposition~4.2], [KS20, \S7.5].
\end{theorem}
\begin{proof}
  \leanok
  Exact first-exceedance decomposition
  $\PR[M_r > x] = \sum_{k<r} \PR[|X_0| > x, \max_{0 < t \le k} |X_t| \le x]$ by stationarity; each term,
  divided by $\PR[|X_0| > x]$, tends to $\PR[\max_{0<t\le k} |Y_t| \le 1]$, and anticlustering controls
  the terms with $k \ge m$. The limit of the Ces`aro average is $\vartheta$. Condition $\mathcal A$
  then gives $(1 - r_n \vartheta \tau / n)^{n/r_n} \to e^{-\vartheta\tau}$.
\end{proof}

\begin{theorem}[$\theta \le \vartheta$ without anticlustering]
  \label{thm:theta-le-vartheta}
  \uses{thm:theta-eq-vartheta}
  \lean{ExtremeValueTheory.HasSpectralTailProcess.le_candidateExtremalIndex_of_tendsto_blockRatio, ExtremeValueTheory.HasSpectralTailProcess.le_candidateExtremalIndex_of_tendsto,
    ExtremeValueTheory.HasSpectralTailProcess.le_candidateExtremalIndex, ExtremeValueTheory.tendsto_measureReal_forall_le_candidateExtremalIndex}
  \leanok
  Every limit point of the cluster ratio is at most $\vartheta$, using only stationarity, the tail
  process, $u_n \to \infty$ and $r_n \to \infty$. Consequently, if $|X|$ has extremal index $\theta$ and
  condition $\mathcal A(u_n)$ holds along one level sequence with $n \PR[|X_0| > u_n] \to \tau > 0$,
  $r_n \to \infty$ and $r_n/n \to 0$, then $\theta \le \vartheta$. Anticlustering is not needed. (The
  hypotheses of the literature version, [KS20, Lemma~7.5.4], were not checked against the source.)
\end{theorem}
\begin{proof}
  \leanok
  Without anticlustering, the first-exceedance terms with $k \ge m$ are bounded above by their limits
  but may lose mass, which gives the inequality.
\end{proof}

\begin{proposition}[Moving maximum: $\theta = \vartheta = 1/2$]
  \label{prop:moving-max-vartheta}
  \uses{thm:candidate-extremal-index, prop:extremal-index-examples}
  \lean{ExtremeValueTheory.movingMax, ExtremeValueTheory.isRegularlyVarying_movingMax, ExtremeValueTheory.hasSpectralTailProcess_movingMax,
    ExtremeValueTheory.candidateExtremalIndex_movingMax, ExtremeValueTheory.hasExtremalIndex_movingMax_candidateExtremalIndex}
  \leanok
  For iid Pareto($\alpha$) $Z$ and $X_t = \max(Z_t, Z_{t+1})$, $\PR[X_0 > x] = 2x^{-\alpha} - x^{-2\alpha}$,
  the spectral tail process is the indicator of $\{-1, 0\}$ or of $\{0, 1\}$ with probability $1/2$
  each, and $\vartheta = 1/2$, matching the directly computed extremal index $\theta = 1/2$.
\end{proposition}
\begin{proof}
  \leanok
  An exceedance of $X_0$ comes from $Z_0$ or $Z_1$, each with conditional probability $1/2$ in the limit.
\end{proof}

